\section{Part A. Point sections in characteristic \texorpdfstring{$p$}{p}}
\label{sec:partA}

We prove A1 for every curve with an abelian prime-to-$p$ Galois
cover of a genus-zero curve, in every characteristic. This includes
every geometrically hyperelliptic curve in odd characteristic.
We also prove it for every abelian prime-to-$p$ Galois cover of a
genus-one curve in odd characteristic, and in characteristic two
when that quotient is supersingular. Other results include a
nonabelian $S_3$ class and explicit ordinary characteristic-two
families. The universal question for arbitrary curves remains
unresolved. Both covering conditions in A2 are
equivalent to equality of full prime-to-$p$ section classes. Their final
implication $P=Q$ therefore has exactly the same remaining scope as A1.
The finite joint norm kernel in the first route of the prompt is
determined exactly and is larger than the proposed subgroup. All
statements concern geometric conjugacy of sections and retain the full
arithmetic Galois group.

\subsection{The covering condition is exactly section equality}

\begin{claim}{A.1}{Proved; verification of R2}
Let $K$ be finitely generated over $\F_p$, and let $G/K$ be a
semiabelian variety. For the continuous prime-to-$p$ Kummer map
\[
 \kappa_G:G(K)\longrightarrow
 H^1_{\mathrm{cont}}\left(G_K,\prod_{\ell\ne p}T_\ell G\right),
 \qquad
 \ker\kappa_G=G(K)[p^\infty].
\]
In particular, equality of point sections of a smooth curve after
geometric prime-to-$p$ abelianization implies that their difference in
its generalized Jacobian is $p$-primary torsion. The modulus is the
reduced boundary of the smooth compactification.
\end{claim}
\begin{proof}
We check the arithmetic input in the previous return. The
Mordell--Weil--Lang--N\'eron theorem, in the form of
\citet[I, Theorem 16.7]{A_Milne2008}, says that $A(K)$ is finitely
generated for an abelian variety $A$ over such a field $K$.
If $L/K$ is a finite splitting field for a torus $T$, a normal proper
model of $L$ over its finite constant field gives
\[
 1\longrightarrow k_0^\times\longrightarrow L^\times
 \xrightarrow{\operatorname{div}}\operatorname{Div}(X).
\]
The image in the free divisor group is free, so $L^\times$ is finite
torsion plus a free abelian group. Thus the subgroup
$T(K)\subset (L^\times)^{\dim T}$ has this same property.
For $1\to T\to G\to A\to1$, the image $B\subset A(K)$ is finitely
generated. Choose lifts of its generators and let $H\subset G(K)$ be
the group they generate. There is a finitely generated direct summand
$T_0$ of $T(K)$ containing $H\cap T(K)$ and all its torsion. Write
$T(K)=T_0\oplus T_1$. Then
\[
 G(K)=(H+T_0)\oplus T_1.
\]
The first summand is finitely generated and the second is free. It
follows, coordinate by coordinate, that
$\bigcap_{(n,p)=1}nG(K)=G(K)[p^\infty]$.

Multiplication by each $n$ prime to $p$ is a finite \'etale isogeny;
the Kummer exact sequence identifies the kernel of
$G(K)\to H^1(G_K,G[n])$ with $nG(K)$. These maps are continuous
cohomological maps for the natural Galois action. A point in their
common kernel has compatible rational division points: any finite
collection of compatibility requirements is realized at a common
multiple of the indices, and compactness of the product of their
finite fibres supplies the inverse limit. Hence the same kernel
statement holds for the full prime-to-$p$ Tate module.

For the final assertion use the Abel map
$U\to J_D$, $R\mapsto[R-P]$, defined over $K$ by a chosen
$P\in U(K)$. Pullback of the $[n]$-torsors gives a continuous
morphism from the arithmetic fundamental extension of $U$ to that of
$J_D$. Equality of point sections therefore kills every Kummer class
of $[Q-P]$. Only this direction is used below; no nonabelian conclusion
is inferred from a single abelianization.
\end{proof}

\begin{claim}{A.2}{Proved; uniqueness of a torsion-compatible lift}
Let $\Omega$ be an algebraically closed field of characteristic $p>0$.
Let $V$ be a smooth connected curve over $\Omega$, with smooth proper
compactification $\overline V$ and reduced boundary $D$. Let a cyclic
group generated by an automorphism $a$ of order $m>1$, with $(m,p)=1$,
act freely on $V$. For $R\in V(\Omega)$ and an integer $n$ coprime to
$m$, the class $[aR-R]\in J_D(\Omega)$ is not killed by $n$.
Consequently, if $f:V\to U$ is a connected \'etale Galois cover over
$\Omega$ with group of order prime to $p$, and $P'\in V(\Omega)$ is
fixed, then above any given $Q\in U(\Omega)$ there is at most one
$Q'$ satisfying $[Q'-P']\in J_D(\Omega)[p^\infty]$.
\end{claim}
\begin{proof}
We use the induced automorphism $\sigma$ on the function field and
divisors, choosing its orientation so that $\sigma R=aR$; replacing
$a$ by $a^{-1}$ has no effect on the assertion. Suppose
$n[aR-R]=0$. The divisor description of the generalized Jacobian
gives $f\in\Omega(\overline V)^\times$ such that
\[
 \operatorname{div}(f)=n(\sigma R-R),\qquad
 f(d)=1\quad(d\in D).
\]
This description is the one in \citet[III, \S8]{A_Milne2008}:
principal divisors are divided out only when the function is congruent
to $1$ along the modulus. The second requirement is absent if $D$ is
empty.

The divisor of $\prod_{i=0}^{m-1}\sigma^i(f)$ is zero, so this product
is a nonzero constant. If $D\ne\varnothing$, evaluation at any
boundary point shows that it is $1$. If $D=\varnothing$, multiply
$f$ by a constant whose $m$-th power is its inverse. This leaves its
divisor unchanged and makes its norm equal to $1$.
Multiplicative Hilbert 90 for
$\Omega(\overline V)/\Omega(\overline V)^{\langle\sigma\rangle}$
then supplies $g$ with $f=\sigma(g)/g$. For clarity, this elementary
form follows by applying the operator
$h\mapsto\sum_{i=0}^{m-1}(f\sigma(f)\cdots\sigma^{i-1}(f))
\sigma^i(h)$ to an $h$ for which it is nonzero; independence of
distinct field automorphisms gives such an $h$, and shifting the sum
gives the desired ratio, possibly after taking an inverse.
It follows that
\[
 E:=\operatorname{div}(g)-nR\quad\hbox{satisfies}\quad
 \sigma E=E,
 \qquad \deg E=-n.
\]

Every interior orbit has length $m$, because the action is free on
$V$. Its contribution to the degree of the invariant divisor $E$ is
therefore divisible by $m$. Consider a boundary orbit, let $d$ be
one of its points, and let its stabilizer have order $e$. Its length
is $m/e$, and a generator of its stabilizer is
$\sigma^{m/e}$. Since $(e,p)=1$, the action on a uniformizer $z$ at
$d$ has leading term $\zeta z$, with $\zeta$ of exact order $e$.
Indeed the finite tame stabilizer of a point on a smooth curve acts
faithfully on its tangent line; the kernel is a $p$-group. One can
also obtain a parameter with this linear action by averaging a
parameter against the tangent character.

Write $g=z^v u$, with $u$ a local unit and $v=\operatorname{ord}_d(g)$.
The value at $d$ of $\sigma^{m/e}(g)/g$ is $\zeta^v$. On the other
hand, telescoping $f=\sigma(g)/g$ expresses this ratio as a product
of $m/e$ translates of $f$, all having value $1$ at $d$. Thus
$e\mid v$. The coefficient of $E$ at $d$ is $v$, since $R$ lies in
$V$. The whole orbit contributes $(m/e)v$, again divisible by $m$.
Every orbit contribution is divisible by $m$, whence $m\mid n$.
This contradicts $(m,n)=1$ and $m>1$.

For the last assertion, two putative lifts differ by some nonidentity
deck transformation $a$. Their difference in $J_D$ is $p$-primary,
and the order of $a$ is prime to $p$. The first part, with $n$ a
sufficiently large power of $p$, rules this out. This argument uses
the entire cyclic subgroup generated by $a$; it does not assume that
the full deck group is abelian.
\end{proof}

\begin{claim}{A.3}{Proved; both covering conditions in A2 are equivalent to section equality}
Let $K/\F_p$ be finitely generated, let $Y/K$ be a hyperbolic curve,
and let $P,Q\in Y(K)$. The following conditions are equivalent.
\begin{enumerate}[label=(\roman*)]
\item The continuous sections $s_P,s_Q:G_K\to E_Y^{(p')}$ are
conjugate by $\Pi_Y^{(p')}$.
\item $[Q-P]$ is $p$-primary in $J_Y(K)$, and the covering condition
of A2 holds for every finite extension $L/K$, every finite \'etale
Galois cover $Y'\to Y_L$ of degree prime to $p$, and every
$L$-rational lift $P'$ of $P$.
\item The same condition holds for all finite \'etale geometrically
connected covers whose geometric Galois closure has degree prime
to $p$, without requiring the cover to be Galois over $L$.
\end{enumerate}
In (ii), ``Galois'' may be interpreted in the stricter arithmetic
sense: all deck transformations are defined over $L$. In the reverse
implication it is enough to quantify over finite separable $L/K$.
This answers the comparison requested at cote 149, manuscript
p.~17, scans 35--36.
\end{claim}
\begin{proof}
Suppose (i). Restriction along $G_L\to G_K$ preserves equality of
section classes. For a cover in (iii), the group attached to it is
an open subgroup $H\subset E_{Y_L}^{(p')}$, and its geometric group
is $H\cap\Pi_Y^{(p')}$. The latter really is the full geometric
prime-to-$p$ group of $Y'$: taking the core of a prime-to-$p$ cover
under the finitely many conjugates in a prime-to-$p$ Galois closure
produces a group which is an extension of a subgroup of a product
of prime-to-$p$ groups by a prime-to-$p$ group. Its order is still
prime to $p$. This proves cofinality of the induced topology.

The chosen $P'$ makes a corresponding representative of $s_P(G_L)$
lie in $H$. Choose the representative of $s_Q$ which equals it.
In the finite geometric fibre of $Y'$ over $Q$ this gives a
$G_L$-fixed point, hence $Q'\in Y'(L)$. In the arithmetic extension
for $Y'$, the two lifted section classes are equal. A.1 now gives
$[Q'-P']\in J_{Y'}(L)[p^\infty]$. Applying A.1 also to $Y$ proves
the initial difference condition. This proves (i)$\Rightarrow$(iii),
and (iii)$\Rightarrow$(ii) is immediate.

For (ii)$\Rightarrow$(i), put $\Pi=\Pi_Y^{(p')}$ and use the actual
point $P$ as geometric base point. The extension is
$E_Y^{(p')}=\Pi\rtimes s_P(G_K)$. Choose a cofinal system of open
characteristic subgroups $N\subset\Pi$. Such a system exists since
the geometric prime-to-$p$ fundamental group of a curve is
topologically finitely generated: SGA~1, XIII, Corollary 2.12 gives
its surface presentation with $2g+\nu$ generators and the usual
single relation \citep{A_SGA1}. The intersections of all subgroups
of index at most $d$ give one choice. Write $A_N=\Pi/N$ and let
\[
 H_N=N\,s_P(G_K)\subset E_Y^{(p')},\qquad
 f_N:Y_N\longrightarrow Y
\]
be the associated connected finite \'etale cover. The subgroup is
open; its projection onto $G_K$ is surjective; and
$[E_Y^{(p')}:H_N]=[\Pi:N]$ is prime to $p$. The cover is
geometrically connected and geometrically Galois with group $A_N$.
The coset $H_N$ supplies a distinguished $P_N\in Y_N(K)$ over $P$.
All these data, including the maps for $N'\subset N$, are defined
over $K$.

The conjugation action of $s_P(G_K)$ on the finite group $A_N$ is
continuous. Choose a finite Galois extension $L/K$ whose Galois
subgroup kills it. Over $L$, the subgroup
$N\,s_P(G_L)$ is normal in $\Pi\rtimes s_P(G_L)$, with quotient
$A_N$. Thus $(Y_N)_L\to Y_L$ is an arithmetic Galois cover of
degree $|A_N|$, and (ii) supplies $Q_N\in Y_N(L)$ with
\[
 f_N(Q_N)=Q,\qquad [Q_N-P_N]\in J_{Y_N}(L)[p^\infty].
\]
By A.2, this is the unique geometric point above $Q$ having the
displayed torsion property. Every element of $G_K$ preserves the
fibre, fixes $P_N$, and preserves $p$-primary torsion. Therefore it
fixes $Q_N$. Since $Q_N$ is defined over a finite separable extension,
this says exactly that $Q_N\in Y_N(K)$.

If $N'\subset N$, the finite map $Y_{N'}\to Y_N$ induces the norm
homomorphism of reduced-modulus generalized Jacobians. It sends
$[Q_{N'}-P_{N'}]$ to the difference of the images; the condition
on the reduced modulus is preserved since the norm of a function
taking the value $1$ at every boundary point again takes value $1$
at the boundary. Hence the image of $Q_{N'}$ has $p$-primary
difference from $P_N$, and uniqueness in A.2 makes it equal to
$Q_N$.

We have obtained a compatible $K$-rational point over $Q$ in every
$P$-pointed characteristic neighbourhood. The inverse limit of
these fibres is the prime-to-$p$ \'etale path torsor from $P$ to $Q$.
The compatible points form a $G_K$-fixed path. Equivalently, after
choosing an arbitrary reference path, its coordinate
$\gamma\in\Pi$ satisfies
$s_Q(g)=\gamma s_P(g)\gamma^{-1}$ for every $g\in G_K$.
This is equality of the full section classes, not just of their
finite images. Notice that the initial condition on $[Q-P]$ was
not needed in this direction: it follows from (i).
\end{proof}

\begin{claim}{A.4}{Proved; the neighbourhood form of A1}
For the same $Y,K,P,Q$, equality of the two section classes is
equivalent to the following condition: every finite \'etale
geometrically connected cover with geometrically prime-to-$p$
Galois closure which has a $K$-point over $P$ has a $K$-point over
$Q$. Consequently nonconjugacy is witnessed by one such cover
having a rational point above $P$ and none above $Q$.
\end{claim}
\begin{proof}
One direction is the fixed-fibre argument in A.3. For the reverse
direction use its neighbourhoods $Y_N$. The finite nonempty sets
$f_N^{-1}(Q)(K)$ need not have surjective transition maps, but they
have a nonempty inverse limit: a finite family of compatibility
conditions is realized by a point at a common refining level.
Compactness of their product therefore gives a compatible family.
The path-torsor argument at the end of A.3 supplies the geometric
conjugator. Negating the equivalence gives the last assertion.
\end{proof}

\subsection{The finite joint norm kernel}

\begin{claim}{A.5}{Proved; exact finite norm kernels, refuting the proposed identification}
Let $C$ be a smooth proper genus-two curve over an algebraically closed
field $\Omega$ of characteristic different from $2$. Choose $P\in C(\Omega)$,
let $j(R)=[R-P]$, and put
\[
 Z=C\times_{j,J_C,[2]}J_C,\qquad
 \pi:Z\longrightarrow C,\qquad G=J_C[2](\Omega)\simeq\F_2^4.
\]
For $r=1,2,3$, let $\mathcal H_r$ be all subgroups $H\subset G$
of codimension $r$, put $r_H:Z\to Z/H$, and define the finite
group scheme
\[
 K^{[r]}=\bigcap_{H\in\mathcal H_r}
 \ker\bigl((r_H)_*:J_Z\longrightarrow J_{Z/H}\bigr).
\]
As abstract groups of geometric points, noncanonically,
\[
 \begin{split}
 K^{[1]}&\simeq (\Z/2)^{16}\oplus(\Z/4)^9,\\
 K^{[2]}&\simeq (\Z/2)^9,\\
 K^{[3]}&=0.
 \end{split}
\]
In particular $K_2$ of R4(e) is $K^{[1]}$, of order $2^{34}$
and exponent $4$. The suggestion $K_2=\pi^*J_C[8]$ is false:
the latter has order $2^8$ and is isomorphic to $(\Z/4)^4$.
These are geometric group identifications; no triviality of their
Galois actions is asserted over a nonclosed field.
\end{claim}
\begin{proof}
The Abel map induces an isomorphism on abelian geometric
fundamental groups \citep[III, Proposition 9.1]{A_Milne2008}.
Thus $Z$ is the connected mod-$2$ homology cover of degree $16$;
Riemann--Hurwitz gives $g(Z)=17$. If $H$ is a hyperplane, write
$N_H=\sum_{h\in H}h_*$ on $J_Z$. Counting the hyperplanes
through a vector of $\F_2^4$ gives
\[
 \sum_H N_H=8\,\mathrm{id}+7N_G.
\]
Since $N_H=r_H^*(r_H)_*$ and $N_G=\pi^*\pi_*$, and since each
$\pi_*$ factors through $(r_H)_*$, the identity shows
$K^{[1]}\subset J_Z[8]$. Every $H\in\mathcal H_2$ or
$\mathcal H_3$ is contained in hyperplanes, so its norm factors
the appropriate hyperplane norms. Hence
$K^{[3]}\subset K^{[2]}\subset K^{[1]}$. All three groups are
finite \'etale $2$-primary groups. It therefore suffices to determine
them on the $2$-adic lattices.

Here is the exact lattice whose calculation is supplied, including
the identification with the algebraic kernel. The homomorphism
\[
 \Phi_r:J_Z\longrightarrow\prod_{H\in\mathcal H_r}J_{Z/H}
\]
has kernel $K^{[r]}$. Its dual, under the canonical principal
polarizations, is the sum of the pullback homomorphisms $r_H^*$.
Factor $\Phi_r$ as an isogeny onto its image followed by inclusion
of an abelian subvariety. Dualizing the latter inclusion has connected
kernel and is surjective on $2$-adic Tate modules. Dualizing the
isogeny gives the Cartier dual of its finite kernel
\citep[I, Theorem 9.1]{A_Milne2008}. Consequently the character
group of $K^{[r]}$ has the same invariant factors as
\[
 \frac{T_2J_Z}{\sum_{H\in\mathcal H_r}r_H^*T_2J_{Z/H}}.       \tag{A.1}
\]
This use of duality retains the finite integral cokernel; replacing
Tate modules by rational vector spaces would lose exactly the group
being computed.

In the complex surface model $T_2J$ is the $2$-adic completion of
$H_1$ of the curve. The map induced by $r_H^*$ is the transfer
$r_H^!:H_1(Z/H,\Z)\to H_1(Z,\Z)$. For example this follows
from its defining intersection identity
$\langle r_H^!x,y\rangle_Z=\langle x,(r_H)_*y\rangle_{Z/H}$;
it is the adjoint of pushforward, just as pullback and norm of
Jacobians are adjoint. In every characteristic different from $2$,
SGA~1, XIII, Corollary 2.12 identifies the prime-to-characteristic
geometric group with the corresponding completed surface
presentation. The mod-$2$ homology subgroup and its intermediate
subgroups are intrinsic in this identification. The group transfer
on their abelianizations is defined by summing lifts of a loop and
commutes with completion. Thus (A.1) is the $2$-adic completion of
the same integral transfer quotient in every such characteristic.
This proves the required comparison, without identifying a
characteristic-$p$ full fundamental group with a characteristic-zero
one.

We specify the integral calculation. Take a genus-two surface with
one vertex, oriented edges $a,b,c,d$, and one face attached along
$aba^{-1}b^{-1}cdc^{-1}d^{-1}$. In its mod-$2$ homology cover,
vertices are $v\in\F_2^4$, an oriented edge $(v,i)$ goes from
$v$ to $v+e_i$, and there is one lifted face at each $v$. There
are $16$ vertices, $64$ edges and $16$ faces. A spanning tree has
$15$ edges; the remaining $49$ edges give a basis of the integral
cycle lattice. In this basis the lifted face boundaries have rank
$15$ and primitive image, so $H_1(Z,\Z)$ has rank $34$.

For a quotient $G\to G/H$, the chain transfer sends the edge
$(\bar v,i)$ to $\sum_{v\mapsto\bar v}(v,i)$.
Its boundary is the transfer of the boundary, and hence it sends
cycles to cycles. A spanning tree in each quotient graph supplies
all its graph cycles; quotient face boundaries already land in the
face-boundary lattice of $Z$. Accordingly it is enough to adjoin
these transferred graph cycles to the $16$ lifted face columns.
Let $M_r$ be the resulting integer matrix with $49$ rows. Its
cokernel is precisely
\[
 H_1(Z,\Z)\big/\sum_{H\in\mathcal H_r}r_H^!H_1(Z/H,\Z).
\]
The matrix dimensions and lattice invariants are
\[
\begin{array}{c|r|r|r|l}
 r&|\mathcal H_r|&\#\text{columns}&|\operatorname{coker}M_r|
     &\operatorname{coker}M_r\\\hline
 1&15&121&2^{34}&(\Z/2)^{16}\oplus(\Z/4)^9\\
 2&35&471&2^9&(\Z/2)^9\\
 3&15&391&1&0
\end{array}                                                        \tag{A.2}
\]

The supplied certificates verify these finite integer assertions as
follows. For each $M_r$ they contain integer matrices $X_r,Y_r,H_r$
with $H_r$ square of order $49$ and
\[
 M_rX_r=H_r,\qquad H_rY_r=M_r.                                     \tag{A.3}
\]
These two identities prove equality of column lattices, independently
of any normal-form label. The certificates give
$|\det H_1|=2^{34}$, $4H_1^{-1}$ integral, and
$\operatorname{rank}_{\F_2}(M_1)=24$.
Thus its cokernel has exponent dividing $4$, order $2^{34}$, and
$49-24=25$ generators modulo $2$. If it is
$(\Z/2)^u\oplus(\Z/4)^v$, these give $u+v=25$ and $u+2v=34$,
whence $(u,v)=(16,9)$. Similarly
$|\det H_2|=2^9$ and $2H_2^{-1}$ is integral, yielding
$(\Z/2)^9$; and $|\det H_3|=1$ yields zero. All arithmetic is
exact. The first matrix is also independently reconstructed using
a different spanning tree by the intermediate-quotient program,
and gives the same Smith factors. The files
\texttt{code/A\_norm\_kernel\_certificate.py} and
\texttt{code/A\_intermediate\_certificate.py} produce and verify
(A.3); their JSON matrices are included. The chain maps and cycles
are generated by the explicitly stated cell structure, not
inferred from the expected output.

Finally $\ker(\pi^*:J_C\to J_Z)=J_C[2]$.
A line bundle trivialized on the connected proper Galois cover
descends from a character of $G$, because its descent multipliers
are global units on $Z$, hence constants. This identifies the
kernel with the $16$ characters of $G$, all of order two, and hence
with $J_C[2]$. Therefore
$\pi^*J_C[8]\simeq J_C[8]/J_C[2]\simeq(\Z/4)^4$.
For each hyperplane $H$, norm followed by pullback gives
$(r_H)_*\pi^*a=8f_H^*a=0$ for $a\in J_C[8]$.
This verifies the previously known inclusion and proves that it
is strict by (A.2).
\end{proof}

The explicit curve $y^2=x^5-x$ over $\F_3$, with $P=\infty$,
already supplies a counterexample to the proposed equality of
kernels: the polynomial and its derivative are coprime modulo $3$,
so this is a smooth genus-two curve, and the orders in A.5 are
$2^{34}$ and $2^8$. This finite check is included in
\texttt{code/A\_finite\_route\_checks.py}.

\begin{claim}{A.6}{Proved; what the finite norm tests establish}
In the situation of A.5 in odd characteristic $p$, assume
$c=[Q-P]\in J_C(\Omega)$ is $p$-primary. There is a unique
$p$-primary $b$ with $2b=c$; put
$\widetilde P=(P,0)$, $\widetilde Q=(Q,b)$, and
$\delta=[\widetilde Q-\widetilde P]\in J_Z(\Omega)$.
For every intermediate quotient $Z/H$, a lift of $Q$ having
$p$-primary difference from the image of $\widetilde P$ is
necessarily the image of $\widetilde Q$.
If all degree-two quotient tests succeed, the prime-to-$p$ part
of $\delta$ lies in $K^{[1]}$; all degree-four quotient tests
put it in $K^{[2]}$; and all degree-eight quotient tests force
$\delta$ itself to be $p$-primary. None of these conclusions
asserts that this last $p$-primary class is zero.
\end{claim}
\begin{proof}
Factor $[2]:J_C\to J_C$ through the quotient abelian variety
$J_C/H$. The isogeny $J_C/H\to J_C$ has degree a power of two,
so it has a unique $p$-primary preimage of $c$. A successful point
lift in $Z/H=C\times_{J_C}(J_C/H)$ maps to that preimage under
the second projection, since this projection sends the chosen
point above $P$ to zero and induces a homomorphism of Jacobians.
The pair of its two coordinates therefore determines that point
as the image of $\widetilde Q$. All quotient tests consequently
apply to the same class $\delta$.

In each of the three cases, choose a power $p^a$ killing the
finitely many norm images of $\delta$. Then
$p^a\delta\in K^{[r]}$. In particular $\delta$ is torsion.
Its prime-to-$p$ component belongs to $K^{[r]}$, since multiplication
by $p^a$ is an automorphism of that $2$-primary kernel. Apply A.5.
The surviving $p$-primary component maps to $p$-primary components
under every norm automatically; no norm identity alone forces
that component to vanish.
\end{proof}

\begin{claim}{A.7}{Proved; general incidence identity and the two-prime comparison}
Let $V\to C$ be a connected finite \'etale Galois cover of smooth
proper curves over an algebraically closed field of characteristic
$p$, with group $G=\F_\ell^d$, where $\ell\ne p$ and $d\geq2$.
For the norm maps to all hyperplane quotients, their common kernel
on $J_V$ is killed by $\ell^{d-1}$. If a cover has deck group
$\F_\ell^d\times\F_m^e$ for distinct primes $\ell,m\ne p$,
the intersection of the two common norm kernels, using the
hyperplanes in each factor over the corresponding complementary
base, is zero. Thus if every norm of a point-divisor difference in
both these families is $p$-primary, that difference is $p$-primary.
This does not imply that the difference vanishes.
\end{claim}
\begin{proof}
For $N_H=\sum_{h\in H}h$ in $\Z[G]$, a nonzero vector is in
$c=(\ell^{d-1}-1)/(\ell-1)$ hyperplanes, and zero is in
$c+\ell^{d-1}$ hyperplanes. Hence
\[
 \sum_HN_H=\ell^{d-1}\,1+cN_G.
\]
As in A.5, applying this identity to a common norm-kernel point
gives $\ell^{d-1}x=0$. On a product-group cover, apply the first
assertion to the cover over $V/\F_\ell^d$ and separately over
$V/\F_m^e$. A point in both kernels is killed by two relatively
prime integers and is zero. If all norm images of $\delta$ are
$p$-primary, choose one power $p^a$ killing all of them.
Then $p^a\delta$ lies in the zero intersection, as required.
These are comparisons on one common cover. Tests on the separate
covers of the original curve alone do not supply the intermediate
norm tests on their common refinement.
\end{proof}

\begin{claim}{A.8}{Proved; exact reduction and some decided cases}
The universal assertion A1 over finite fields is equivalent to
injectivity of the point-section map for every hyperbolic curve
over every field finitely generated over $\F_p$. It is also
equivalent to the universal final implication of A2. Independently
of A1, injectivity holds for curves of genus zero and for affine
curves of genus one, and for a proper curve $C/k$ over a finite
field whenever $J_C(k)[p^\infty]=0$. In particular it holds over
every finite constant extension when $C$ has $p$-rank zero.
\end{claim}
\begin{proof}
In genus zero the generalized Jacobian is a torus; it has no
nontrivial geometric $p$-primary torsion, and its Abel map is
injective for a hyperbolic curve. Explicitly, after splitting the
punctures, a cross-ratio coordinate gives an injective coordinate
in that map. A.1 proves injectivity.

For genus one, use $P$ as the origin of the compactification $E$.
If a collision exists, A.1 gives a nonzero class $Q-P$ of order
$p^a$. Choose $m$ prime to $p$ with $m^2>p^a$ and use
$E\setminus[m]^{-1}D\to E\setminus D$.
After a finite extension its deck group and all boundary points
are rational. Equality of sections lifts to a pair $0,Q'$ as in
A.3. The generalized-Jacobian difference has $p$-power order;
its projection to $E$ has order exactly $p^a$, since multiplication
by $m$ is invertible on $p$-primary torsion. The toric kernel has
no such torsion, so the generalized difference also has order
$p^a$. There is consequently a function with divisor
$p^a(Q'-0)$ which is $1$ at every point of $[m]^{-1}D$.
Its degree is $p^a$, but the function minus $1$ has at least
$m^2>p^a$ distinct zeros and no poles there. This is impossible.
This verifies the step of R3 used in the reduction, for every $p$,
including $2$ and $3$.

For genus at least two, the open immersion into the smooth
compactification induces a continuous quotient map of arithmetic
prime-to-$p$ extensions over the identity of $G_K$. A geometric
conjugator maps to a geometric conjugator. Thus an affine collision
would be a proper collision.

We spell out the remaining specialization direction. Spread a
smooth proper pair and the distinct rational sections $P,Q$ over
a smooth connected finite-type scheme $S/\F_q$ with function
field $K$. Shrink $S$ so that the boundary is finite \'etale and
the sections are disjoint. For a closed point $s$ put
$R=\mathcal O_{S,s}^{\mathrm h}$, $L=\operatorname{Frac}(R)$,
and $k=\kappa(s)$. The field $k$ is finite and
$G_L\twoheadrightarrow\pi_1(R)=G_k$ is surjective: a connected
finite \'etale cover of the normal integral henselian local
scheme remains connected over its fraction field.
SGA~1, XIII, Proposition 4.3 supplies the geometric
prime-to-$p$ exact sequence for the family; Examples 4.4 verify
its hypotheses for our smooth proper compactification, relative
normal-crossing boundary and section. Its 4.5.3 identifies the
relative geometric group after base change. In particular,
with base points chosen along $P$, there is a continuous
specialization homomorphism
\[
 \operatorname{sp}:E_{Y_L}^{(p')}\longrightarrow E_{Y_s}^{(p')}
\]
whose restriction to the geometric kernels is an isomorphism,
whose arithmetic quotient is $G_L\twoheadrightarrow G_k$, and
which satisfies
\[
 \operatorname{sp}s_{P_L}=s_{P_s}\operatorname{pr},\qquad
 \operatorname{sp}s_{Q_L}=s_{Q_s}\operatorname{pr}.
\]
A geometric conjugator over $L$ therefore specializes to a
geometric conjugator over $k$, since $G_L\to G_k$ is surjective.
The points remain distinct by the chosen disjointness. A collision
over $K$ thus produces a collision over a finite field of the same
type. The converse universal implication holds because finite
fields are finitely generated over their prime fields.
Combining this with the genus-zero and genus-one cases gives the
first equivalence, and A.3 gives the second.

Finally A.1 puts a collision difference in $J_C(k)[p^\infty]$.
If that group is zero, the Abel embedding forces $P=Q$. If the
$p$-rank is zero, the geometric $p$-primary torsion point group
is zero, so the same argument works over every finite extension.
This statement concerns geometric points of the torsion scheme,
not its possibly nonreduced local part.
\end{proof}

\subsection{An explicit characteristic-two separation}

\begin{claim}{A.9}{Proved; a degree-$27$ cover in characteristic two}
Let $C/\F_2$ be the smooth proper curve with affine equation
\[
 Y^2+(x^2+x)Y=x^5+x^3+x^2+x.
\]
It has genus two and exactly three $\F_2$-points
$P_0=(0,0)$, $P_1=(1,0)$ and $P_\infty$. Their pairwise
differences in $J_C$ have exact order two. Nevertheless their
section classes in the geometrically prime-to-$2$ arithmetic
fundamental group remain pairwise distinct after every finite
constant extension.

More explicitly, over $k=\F_{16}$ there is a connected finite
\'etale cover of degree $27$, with geometric Galois closure of odd
degree, having a $k$-point above $P_\infty$ and none above $P_0$.
Thus this example satisfies the requested whole-group separation;
it is not evidence of conjugacy inferred from finite quotients.
\end{claim}
\begin{proof}
Away from $x=0,1,\infty$ set $y=Y/(x^2+x)$. The equation becomes
\[
 y^2+y=x+\frac1x+\frac1{x+1}.
\]
This Artin--Schreier double cover of the projective line has three
simple poles. At each pole the different exponent is two, giving
$2g(C)-2=-4+6=2$ and hence genus two. Each pole has one rational
point above it; there are no other $\F_2$-values of $x$.
The divisors of $x$, $x+1$ and $x/(x+1)$ give twice the three
pairwise differences. None is principal, since a function with
one simple pole would make $C$ rational. This also verifies the
complete failure of odd abelianized sections to distinguish these
three points.

Choose $\alpha\in\F_4$ with $\alpha^2+\alpha+1=0$, and
$\beta\in k$ with $\beta^2+\beta=\alpha^2$. Define
\[
 \begin{split}
 h&=x^2+x,\qquad f=x^5+x^3+x^2+x,\\
 u&=x^2+\alpha,\\
 A&=x^3+(\beta+1)x^2+\beta x+1.
 \end{split}
\]
A direct expansion in $k[x]$ gives
\[
 A^2+hA+f=u^3.                                                     \tag{A.4}
\]
Put $F=A+Y\in k(C)^\times$. The hyperelliptic involution sends
$Y$ to $Y+h$, so $F\,\iota(F)=u^3$. The function $F$ has a
pole of order six at $P_\infty$ and no other pole.
The only root of $u$ is $r=\alpha^2$, and $h(r)=1$. Thus the
two points above $r$ are unramified for $x$; exactly one, call it
$R$, has $F(R)=0$. Equation (A.4) gives
\[
 \operatorname{div}(F)=6R-6P_\infty.
\]
In particular the normalization $D$ of $C$ in
$k(C)(z)$, $z^3=F$, is \'etale over $C$: at every point the
valuation of $F$ is divisible by three, and three is invertible.
It is geometrically connected. Otherwise $F$ would be a cube
over $\overline k(C)$, whose cube root had its only pole of order
two at $P_\infty$. Such a function is in the span of $1,x$, by
the hyperelliptic degree-two pencil, whereas its cube cannot have
the nonzero $Y$ coefficient occurring in $F$. Since $k$ contains
the cube roots of unity, $D\to C$ is Galois of degree three.
Riemann--Hurwitz now gives $g(D)=4$.

The involution $(x,z)\mapsto(x,u/z)$ lifts the hyperelliptic
involution. Its invariant function $w=z+u/z$ satisfies
\[
 w^3+u w=h.
\]
Consequently $D$ has a degree-two map $r_D$ to the plane cubic
\[
 E:\quad
 W^3+(X^2+\alpha Z^2)W=X^2Z+XZ^2.                                \tag{A.5}
\]
This cubic is geometrically smooth: the derivatives in $X$ and
$Z$ are $Z^2$ and $X^2$, so a singular projective point would
have $X=Z=0$, and the derivative in $W$ then forces $W=0$.
With $O=[1:0:0]$ it is an elliptic curve. The map has degree two
because $z$ satisfies $z^2+wz+u=0$ over $k(E)$, and the displayed
involution is nontrivial.

At $P_\infty$ the ratio $F/x^3$ has residue $1$. The three
points of $D$ above it are rational over $k$; choose $P'$ with
residue $z/x=1$. Since $u/x^2$ has residue $1$, the ratio
$w/x$ has residue zero at this point, so $r_D(P')=O$.
At $P_0$ one has $F=1$, hence the three lifts have
$z=1,\alpha,\alpha^2$. Their images on $E$ are
\[
 [0:\alpha^2:1],\quad[0:\alpha^2:1],\quad[0:0:1].              \tag{A.6}
\]
They are nonzero points of exact order three in $E(k)$.

For completeness the last finite assertion is checked by an
exhaustive exact program, including the group operation. In the
field model $k=\F_2[t]/(t^4+t+1)$ the program takes
$\alpha=t^2+t$ and $\beta=t^2$. It finds exactly nine points on
(A.5) over $k$. For a smooth plane cubic and an arbitrary rational
origin $O$, the sum of two points is obtained by taking the third
intersection of their chord or tangent and then the third
intersection of the line joining that point to $O$.
This rule is valid without assuming $O$ is a flex: subtracting the
two line divisors gives $P+Q\sim O+(P\oplus Q)$.
The program computes all sums, verifies associativity on all
$9^3$ triples, and obtains one point of order one and eight
points of order three. It also checks (A.4), all three values
in (A.6), smoothness, and the base-curve counts
$|C(\F_2)|=3$, $|C(\F_4)|=3$, $|C(\F_{16})|=31$.
The program and its output are
\texttt{code/A\_char2\_cubic\_test.py} and its JSON output.

Now take the fibre product
\[
 Z'=D\times_{r_D,E,[3]}E\longrightarrow D\longrightarrow C.
\]
It has degree $9\cdot3=27$ and is finite \'etale.
It is geometrically connected: norm and pullback for the
degree-two map $D\to E$ compose to $[2]$, so
$T_3J_D\to T_3E$ is surjective; its quotient modulo three is
the monodromy of this pullback of $[3]$. The action is therefore
transitive on its nine geometric sheets.
Its geometric Galois closure has odd degree: take the compositum
of the three conjugates of its degree-nine Galois cover over
the cyclic degree-three cover $D/C$. The resulting group is an
extension of a subgroup of $E[3]^3$ by $\Z/3$.
The pair $(P',O)$ is a $k$-point above $P_\infty$.
Since $E(k)$ has exponent three, $[3]E(k)=\{O\}$; (A.6)
shows that no point above $P_0$ lifts to $Z'(k)$. This proves
the asserted finite separating cover.

To see persistence after arbitrary finite constant extension,
suppose the two section classes became equal. Extend constants
further to contain $k$. The chosen $P'$ on $D$ would then have
a matching lift $Q'$ over $P_0$ whose generalized-Jacobian
difference is $2$-primary, by A.3. Applying $r_{D*}$ would make
$[r_D(Q')-O]$ $2$-primary in $E$, contradicting the exact order
three of every point in (A.6). Thus no such extension exists.
Finally $x\mapsto x+1$ and $x\mapsto1/x$ each change the
Artin--Schreier right side by $1$ and lift over $\F_4$ by a
translation $y\mapsto y+c$, $c^2+c=1$. They permute the three
branch points transitively on unordered pairs. The same
nonconjugacy conclusion therefore holds for each pair among
$P_0,P_1,P_\infty$.
\end{proof}

\begin{claim}{A.10}{Proved; a point-separating quotient criterion}
Let $C/k$ be a smooth proper geometrically connected curve over a
finite field of characteristic $p$. Suppose that, after a finite
constant extension, there is a geometrically connected finite
\'etale cover $D\to C$ with geometrically prime-to-$p$ Galois
closure, together with morphisms
\[
 r_i:D\longrightarrow B_i\qquad(1\le i\le t)
\]
to abelian varieties satisfying
$B_i(\overline k)[p^\infty]=0$. Suppose that the product of the
$r_i$ is injective on geometric points. Then the point-section
map for $C$ is injective after every finite constant extension.

This criterion applies to the ordinary genus-two curve in A.9.
Thus its point-section map is injective over every finite
extension of $\F_2$, for all its rational points, not only for
the three points rational over $\F_2$.
\end{claim}
\begin{proof}
Suppose that $P,Q\in C(L)$ give equal section classes for some
finite constant extension $L/k$. Extend constants further so that
the cover and the maps are defined and some chosen point $P'$ of
$D$ above $P$ is rational. By A.3 there is a point $Q'$ above $Q$
such that $[Q'-P']\in J_D[p^\infty]$. Applying the homomorphism
of Albanese varieties induced by each $r_i$ gives
$r_i(Q')-r_i(P')\in B_i(\overline k)[p^\infty]=0$.
The joint injectivity of these maps implies $Q'=P'$ and hence
$Q=P$. This proves the criterion without asserting that a
single fixed finite cover distinguishes all pairs over all
constant extensions.

We verify the criterion in A.9, retaining its notation. The
elliptic curve $E$ in (A.5) is defined over $\F_4$. The following
complete affine point count uses $\alpha^2+\alpha+1=0$:
\[
\begin{array}{c|c|c|c}
 x&x^2+\alpha&x^2+x&\{w\in\F_4:w^3+(x^2+\alpha)w=x^2+x\}\\
 \hline
 0&\alpha&0&\{0,\alpha^2\}\\
 1&\alpha^2&0&\{0,\alpha\}\\
 \alpha&1&1&\varnothing\\
 \alpha^2&0&1&\{1,\alpha,\alpha^2\}
\end{array}
\]
There are also the two points $O=[1:0:0]$ and
$T=[1:1:0]$ at infinity. Hence $|E(\F_4)|=9$, and the
characteristic polynomial of its $\F_4$-Frobenius $\pi$ is
$(X+2)^2$, by the degree and point-count formula for Frobenius
\citep[II, Theorem 1.1]{A_Milne2008}. The action on a prime-to-two
Tate module is faithful, so the characteristic-polynomial
identity gives $(\pi+[2])^2=0$ in $\operatorname{End}(E)$
\citep[I, Theorem 10.15]{A_Milne2008}.
A nonzero endomorphism of an elliptic curve is an isogeny and
cannot have zero square. Therefore $\pi=[-2]$. Frobenius is
purely inseparable, so $[2]$ has no nonzero geometric kernel
point. Iterating this observation proves
\[
 E(\overline{\F}_2)[2^\infty]=0.                                  \tag{A.7}
\]

Put $v=u/z$ on $D$, and let $d:z\mapsto\alpha z$ generate the
cubic deck group. Define two morphisms to this same elliptic
curve by the rational formulas
\[
 r_0=(x,w_0),\quad w_0=z+v,
 \qquad
 r_1=(x,w_1),\quad w_1=\alpha z+\alpha^{-1}v.
\]
They extend to morphisms on the smooth proper curve $D$; they
are $r_D$ and $r_D\circ d$. We check that they jointly separate
all geometric points, including points where these formulas
have zeros or poles.

At a finite value of $x$ with $u\ne0$, both $z$ and $v$ are
regular units. Equality of both image points implies equality
of $x,w_0,w_1$. The coefficient matrix
\[
 \begin{pmatrix}1&1\\ \alpha&\alpha^{-1}\end{pmatrix}
\]
has determinant $\alpha+\alpha^{-1}=1$. It therefore recovers
both $z$ and $v$. Since $Y=z^3-A$, these values recover the
point on the normalization $D$ at such a point.

At $x=r=\alpha^2$, one has $u=0$ and $h=1$. On the branch
$R$ where $F=0$, the valuations of $u,F,z,v$ are respectively
$2,6,2,0$. Thus $z=0$, $v\in\mu_3$, and
\[
 (w_0,w_1)=(v,\alpha^{-1}v).
\]
On the other branch, $F=1$, so $z\in\mu_3$, $v=0$, and
\[
 (w_0,w_1)=(z,\alpha z).
\]
Each formula distinguishes its three points, and the two sets
of pairs are disjoint because $\alpha\ne\alpha^{-1}$.

Above $x=\infty$ the three points have residues
$\xi=z/x\in\{1,\alpha,\alpha^2\}$. The residue of $v/x$ is
$\xi^{-1}$. The image pairs $(r_0,r_1)$ are, in this order,
\[
 (O,T),\qquad(T,T),\qquad(T,O),
\]
which are distinct. Images of finite points have $Z\ne0$ and
cannot coincide with these. This proves the required geometric
point injectivity. Combining it with (A.7) and the already
proved cubic \'etaleness proves the second assertion.
\end{proof}

\subsection{A limitation of the purely inseparable route}

\begin{claim}{A.11}{Proved; minimal torsion functions are separable}
Let $C$ be a smooth proper curve over an algebraically closed field
$\Omega$ of characteristic $p$, and let $P\ne Q$ have difference
of exact order $p^a$, with $a\ge1$, in $J_C(\Omega)$.
Every function $f$ with divisor $p^a(Q-P)$ defines a separable
morphism $C\to\mathbf P^1$ of degree $p^a$.
If $C$ has genus two and positive $p$-rank, such differences
occur with arbitrarily large $a$. Thus the existence of
infinitely many $p$-primary point differences, by itself, does
not force these functions to be purely inseparable.
\end{claim}
\begin{proof}
If $f=g^p$, then
$\operatorname{div}(g)=p^{a-1}(Q-P)$, contradicting the
exact order. Thus $f\notin\Omega(C)^p$. Over the perfect
field $\Omega$, the kernel of
$d:\Omega(C)\to\Omega^1_{\Omega(C)/\Omega}$ is
$\Omega(C)^p$, so $df\ne0$. The differential criterion for
a function on a smooth curve says that
$\Omega(C)/\Omega(f)$ is separable. Its degree is the pole
degree $p^a$.

For genus two choose a geometric Weierstrass point $B$, so
$K_C\sim2B$, and let $\iota$ be the hyperelliptic involution.
Riemann--Roch shows that each degree-two divisor class is
effective. Hence every $c\in J_C(\Omega)$ has a presentation
\[
 c=[P+R-2B]=[P-\iota R].
\]
In particular the difference map $C\times C\to J_C$ is
surjective on geometric points. Positive $p$-rank means that
$J_C(\Omega)[p^\infty]\simeq(\mathbf Q_p/\mathbf Z_p)^f$
with $f>0$; it contains points of exact order $p^a$ for every
$a$. Each such point has the displayed difference presentation.
Over $\overline{\F}_p$ each pair is defined over some finite
constant extension. This conclusion concerns individual
Jacobian differences and does not supply compatible pairs in
the prime-to-$p$ tower.
\end{proof}

\begin{claim}{A.12}{Proved; necessary growth of lifted torsion orders}
Let $C/\F_q$ be smooth proper geometrically connected of genus
$g\ge2$, and suppose that distinct $P,Q\in C(\F_q)$ have equal
point-section classes in the geometrically prime-to-$p$ group.
Put $J=J_C$ and $|J(\F_q)[p^\infty]|=p^v$. There is an integer
$N\ge1$ such that for every positive $m\equiv1\pmod N$ there is
a geometrically connected finite \'etale cover
$X_m\to C_{\F_{q^m}}$, with abelian geometric Galois group of
order prime to $p$ and degree
\[
 d_m=\frac{|J(\F_{q^m})|}{p^v},
\]
with the following property. For a chosen rational lift $P_m$,
the unique lift $Q_m$ whose difference from $P_m$ is
$p$-primary has exact order $p^{a_m}$ satisfying
\[
 p^{a_m}\ \ge\ \frac{2d_m}{q^m+1}
 \ \ge\ \frac{2(q^{m/2}-1)^{2g}}{p^v(q^m+1)}.                    \tag{A.8}
\]
In particular these orders are unbounded. A uniform bound on
the orders of compatible lifted differences, if established
on these covers, would contradict a collision.
\end{claim}
\begin{proof}
Let $F$ be the $q$-power Frobenius on $J$. Its action on the
finite group $B=J(\overline{\F}_q)[p^{v+1}]$ has finite order;
take that order for $N$. If $m\equiv1\pmod N$, $F^m$ and $F$
have the same fixed points in $B$. This implies
\[
 A_m:=J(\F_{q^m})[p^\infty]=A:=J(\F_q)[p^\infty].
\]
Indeed $A\subset A_m$ and $|A|=p^v$. If $|A_m|>p^v$, then
either $A_m$ is killed by $p^{v+1}$, or an element of larger
order supplies a cyclic subgroup of order $p^{v+1}$ in
$A_m[p^{v+1}]$. In either case $|A_m[p^{v+1}]|>p^v$, contrary
to equality of fixed groups in $B$.

Over $k_m=\F_{q^m}$, the Lang isogeny $F^m-1:J\to J$ is
finite \'etale with constant kernel $J(k_m)$: its derivative
is $-1$, and its geometric kernel is the finite group of
Frobenius-fixed points. Quotient its source by the constant
subgroup $A$ to obtain
\[
 \lambda_m:J/A\longrightarrow J,
 \qquad\ker\lambda_m=J(k_m)/A.
\]
This is finite \'etale of degree $d_m$, prime to $p$, and
\[
 \lambda_m((J/A)(k_m))=A.                                       \tag{A.9}
\]
For if a class $\bar b$ is rational, then $F^m b-b\in A$.
Conversely, if $c\in A$, choose a geometric $b$ with
$F^m b-b=c$; the class of $b$ modulo $A$ is rational and
maps to $c$. This proves (A.9).

Pull back $\lambda_m$ by $j:C\to J$, $R\mapsto[R-P]$, to
obtain $X_m$. The Abel map surjects on the prime-to-$p$
geometric abelian fundamental group
\citep[III, Proposition 9.1 and \S9]{A_Milne2008}; therefore
this pullback of the connected isogeny is geometrically
connected. Its deck group is the constant group $J(k_m)/A$.
The fibre above $P$ has a rational point given by zero and
hence consists of $d_m$ rational points. By A.1,
$[Q-P]\in A$; (A.9) gives a rational point in the fibre above
$Q$, so it too has $d_m$ rational points. Thus
\[
 |X_m(k_m)|\ge2d_m.
\]
Equality of the original section classes restricts to $G_{k_m}$
and lifts with the chosen $P_m$, by A.3. Hence a matching
$Q_m$ exists and is unique by A.2. Its difference is nonzero,
since the Abel map of the positive-genus curve $X_m$ is
injective and $P_m,Q_m$ lie over distinct points.

The exact order $p^{a_m}$ supplies a function with divisor
$p^{a_m}(Q_m-P_m)$. The function can be defined over $k_m$:
the one-dimensional space of functions having that divisor
descends by multiplicative Hilbert 90. It is a nonconstant
map $X_m\to\mathbf P^1$ of degree $p^{a_m}$. Each rational
fibre has at most $p^{a_m}$ rational points, so
$|X_m(k_m)|\le p^{a_m}(q^m+1)$. This proves the first
inequality in (A.8). The second follows from the Frobenius
point-count formula and the absolute value $q^{1/2}$ of its
$2g$ eigenvalues on $J$ \citep[II, Theorem 1.1]{A_Milne2008}.
The lower bound tends to infinity because $g\ge2$.
\end{proof}

\subsection{Hyperelliptic curves in characteristics congruent to two modulo three}

\begin{claim}{A.13}{Proved; a controlled Hesse correspondence}
Let $\Omega$ be algebraically closed of characteristic different
from $2,3$. Let $g:D\to E$ be a ramified degree-two morphism
of smooth proper connected curves, where $E$ has genus one,
and choose a branch point $O\in E(\Omega)$ as its origin.
There is a connected finite \'etale cover $a:W\to D$ of
degree $9$ or $18$, whose geometric Galois closure has order
divisible only by $2,3$, and a morphism
\[
 r:W\longrightarrow E_* ,\qquad E_*:v^2=t^3+1,
\]
such that, for any $S,T\in W(\Omega)$,
\[
 r(S)=r(T)\quad\Longrightarrow\quad
 g(a(T))=g(a(S))\ \hbox{or}\ g(a(T))=2O-g(a(S)).                 \tag{A.10}
\]
All these data descend to a finite constant extension when the
original data are defined over $\overline{\F}_p$.
\end{claim}
\begin{proof}
We use the geometry underlying
\citet[\S4, especially Lemma 4.3]{A_BogomolovTschinkel2002},
keeping track of the degrees and the implication (A.10).
The Hesse normal-form theorem
\citep[Lemma 2.1, with the characteristic hypothesis in \S2]{A_ArtebaniDolgachev2009}
puts the elliptic curve, embedded by $|3O|$, in the form
\[
 E:\quad X^3+Y^3+Z^3+\lambda XYZ=0,
 \qquad O=(1:0:-1).
\]
The coordinate permutations and diagonal cube-root scalings
permute its nine flexes transitively, so the prescribed origin
can be put at the displayed flex. The flexes are precisely
$E[3]$; explicitly, for $\zeta^3=1$, $\zeta\ne1$, they are
the points with one zero coordinate and the ratio of the other
two equal to $-1,-\zeta,-\zeta^2$.

Projection from $O$ gives a degree-two morphism
\[
 \pi:E\longrightarrow\mathbf P^1,
 \qquad t=(X+Z)/Y.
\]
The apparent base point at $O$ is removed by the linear system
$|3O-O|=|2O|$. More explicitly, put $X=1$ and $Z=-1+s$
near $O$. The equation has linear term $3s-\lambda Y$, so
$\pi(O)=\lambda/3$. Smoothness is equivalent to
$\lambda^3\ne-27$; hence this value is distinct from the
four values in (A.11). This also checks the exceptional fibre
at the base flex directly.
It is the quotient by $(X:Y:Z)\mapsto(Z:Y:X)$. This involution
fixes $O$ and hence is inversion for the group law with origin
$O$. Indeed a nonidentity involution fixing the origin is a
group endomorphism $u$ with $(u-1)(u+1)=0$; nonzero
endomorphisms of an elliptic curve are isogenies, so $u=-1$.
The eight nonzero three-torsion points have exactly the four
images
\[
 B_*:=\pi(E[3]\setminus\{O\})
       =\{\infty,-1,-\zeta,-\zeta^2\}.                         \tag{A.11}
\]
Each fibre in (A.11) consists of an opposite pair of nonzero
three-torsion points. It is unramified for $\pi$, since a
nonzero three-torsion point is not two-torsion.

Form $D_3=D\times_{g,E,[3]}E$ and denote its projections by
$a_3:D_3\to D$ and $b_3:D_3\to E$. It is smooth and
connected: the degree-two extension $\Omega(D)/\Omega(E)$
and the degree-nine Galois extension supplied by $[3]$ have
trivial intersection. Thus $a_3$ is a connected \'etale
$E[3]$-torsor of degree nine. Since $g$ is branched at $O$,
$b_3$ has ramification index two above each point of $E[3]$.
Consequently the morphism $\pi b_3:D_3\to\mathbf P^1$ has
even local ramification index at every point above $B_*$.

The smooth double cover $\pi_*:E_*\to\mathbf P^1$ given by
$v^2=t^3+1$ is branched exactly at $B_*$. Normalize the fibre
product $D_3\times_{\mathbf P^1}E_*$ and choose a connected
component $W$. Its map to $D_3$ is \'etale: the pullback
Kummer function has even valuation at every point, and two
is invertible. Its degree is one or two. The resulting
$a:W\to D$ has degree nine or eighteen. The Galois closure
over $D$ is obtained by adjoining the at most nine deck
conjugates of this quadratic extension over $D_3$; its group
is an extension of a subgroup of $(\Z/2)^9$ by $(\Z/3)^2$.
This proves the claimed prime control, including the case of
a split quadratic pullback.

Let $r:W\to E_*$ be the other projection. If $r(S)=r(T)$,
then $\pi(b_3(S))=\pi(b_3(T))$, where the intermediate
projection from $W$ is understood. The quotient description
of $\pi$ gives $b_3(T)=b_3(S)$ or $-b_3(S)$.
Applying $[3]$ gives (A.10). Every map, component and chosen
point used here is defined by finitely many elements of
$\overline{\F}_p$, so a finite constant extension suffices
for all descent assertions needed below.
\end{proof}

\begin{claim}{A.14}{Proved; collisions have equal images in a ramified elliptic quotient}
Let $k$ be finite of characteristic $p\ge5$ with
$p\equiv2\pmod3$, and let $g:D\to E$ be a ramified degree-two
morphism of smooth proper geometrically connected curves over
$k$, where $E$ has genus one. If two point sections of $D$
are equal in the geometrically prime-to-$p$ group, then their
images under $g$ are equal. This holds after every finite
constant extension.
\end{claim}
\begin{proof}
First, $E_*:v^2=t^3+1$ has $p+1$ points over $\F_p$.
The map $t\mapsto t^3$ permutes $\F_p$, since $(3,p-1)=1$;
therefore the quadratic-character sum of $t^3+1$ is zero.
The Frobenius characteristic polynomial is $X^2+p$, giving
$F^2=[-p]$. Thus $[p]$ is purely inseparable and
\[
 E_*(\overline k)[p^\infty]=0.                                  \tag{A.12}
\]
Here the characteristic-polynomial identity and faithfulness
on a prime-to-$p$ Tate module are the same standard facts
used in A.10. This argument verifies the supersingularity
needed in the proof directly.

Suppose that $P,Q\in D(k)$ have equal section classes.
Finite constant extensions preserve that equality. Choose
an origin on $E$ after such an extension and form the
connected \'etale degree-nine cover
\[
 D_3=D\times_{g,E,[3]}E,
 \qquad g_3:D_3\to E.
\]
Connectedness follows from the coprime degrees two and nine,
as in A.13. Choose a geometric lift $P_3$ of $P$ and extend
constants to make it rational. The section-lifting argument
of A.3 gives a lift $Q_3$ of $Q$ whose full prime-to-$p$
section class equals that of $P_3$.

Choose a branch point $b$ of $g$. Every point of $[3]^{-1}b$
is a branch point of $g_3$. In particular we can choose
$O_0,O_1$ there with $O_1-O_0$ of exact order three, and
make these points rational by a further finite extension.
For each $j=0,1$, apply A.13 to $g_3$ with origin $O_j$.
Its cover has geometrically prime-to-$p$ Galois closure
because $p\ne2,3$. Lift the equal section classes with a
chosen lift of $P_3$ to that cover. By A.1 the difference
of the two images on $E_*$ is $p$-primary, and (A.12) makes
the images equal. Formula (A.10) therefore says, writing
$b_P=g_3(P_3)$ and $b_Q=g_3(Q_3)$,
\[
 b_Q=b_P\quad\hbox{or}\quad b_Q=2O_j-b_P
 \qquad(j=0,1).
\]
If $b_Q\ne b_P$, these two equations give
$2(O_1-O_0)=0$, contrary to its exact order three.
Hence $b_Q=b_P$, and applying $[3]$ gives $g(Q)=g(P)$.
All comparisons took place after finite extensions along
actual Galois subgroups; thus the argument applies over
every finite constant extension of the original field.
\end{proof}

\begin{claim}{A.15}{Proved; A1 for all hyperelliptic curves when $p\equiv2\pmod3$, $p\ge5$}
Let $k$ be finite of characteristic $p\ge5$,
$p\equiv2\pmod3$. For every smooth proper geometrically
connected hyperelliptic curve $C/k$ of genus at least two,
the map from rational points to geometric conjugacy classes
of sections of $E_C^{(p')}\to G_k$ is injective, after every
finite constant extension. In particular this proves A1
for every genus-two curve in these characteristics.

By the specialization comparison in A.8, the same conclusion
holds for geometrically hyperelliptic proper curves and their
hyperbolic open subcurves over fields finitely generated over
$\F_p$ in these characteristics.
\end{claim}
\begin{proof}
We first establish a consequence that applies in every
hyperelliptic genus in the stated characteristics. If two
points of a hyperelliptic curve $H$ have equal sections, then
they have the same image under its hyperelliptic map
$x_H:H\to\mathbf P^1$.
To prove this, extend the finite field so that the branch
points and a coordinate on the quotient are defined. Choose
four of the $2g(H)+2$ branch points, and let $E\to\mathbf
P^1$ be the double cover branched at those four points.
The normalization $D$ of $H\times_{\mathbf P^1}E$ is a
connected \'etale double cover of $H$. On function fields
the two quadratic characters have distinct branch sets,
which proves connectedness, and all pullback valuations
are even, which proves \'etaleness. Its map $D\to E$ has
degree two and is ramified above the remaining branch
points of $H$; this set is nonempty because $g(H)\ge2$.
Lift the two equal sections to $D$ with a chosen rational
lift of the first point, using A.3 and a finite extension
for that lift. A.14 says that the resulting points have
equal images on $E$. Their images under $x_H$ are therefore
equal, proving the consequence.

We now prove injectivity by induction on $g=g(C)\ge2$.
After a finite extension split the branch set of the
hyperelliptic map into disjoint subsets
\[
 S=S_2\sqcup S_{2g},\qquad |S_2|=2,
 \quad |S_{2g}|=2g.
\]
Let $B\to\mathbf P^1$ and $H\to\mathbf P^1$ be the double
covers with branch sets $S_2$ and $S_{2g}$, respectively.
Their genera are zero and $g-1$. Since their branch sets
are disjoint, the fibre product
\[
 V=B\times_{\mathbf P^1}H
\]
is smooth and connected: at a branch point of either
factor, the other factor is \'etale; the quadratic
function-field extensions are independent. The diagonal
involution of this fibre product acts freely, and its
quotient is the double cover branched at $S$, hence is
$C$ after the chosen constant extension. Thus $V\to C$
is a connected \'etale double cover.

The map $V\to B\simeq\mathbf P^1$ is a hyperelliptic
double cover with $4g$ branch points, so $g(V)=2g-1$.
If distinct points of $C$ had equal sections, choose their
lifted equal sections at points $P_V,Q_V$ of $V$, as above.
The consequence proved in the first paragraph, applied to
the hyperelliptic curve $V$, gives equal images of these
points on $B$.
If $g=2$, the curve $H$ is elliptic and $V\to H$ is a
ramified double cover, with ramification above the four
points lying over $S_2$; A.14 gives equal images on $H$.
If $g>2$, functoriality gives equal point-section classes
on the hyperelliptic curve $H$ of genus $g-1$, and the
induction hypothesis gives equal images on $H$.
The two projections $V\to B,H$ jointly distinguish every
geometric point, since $V$ is their smooth fibre product.
Therefore $P_V=Q_V$, contradicting their distinct images
on $C$. This proves the induction.

For the final assertion, spread the hyperelliptic double
cover and its distinct branch divisor along with the two
sections. After shrinking the base, the same smooth
hyperelliptic type persists on the closed finite-field
fibre. The specialization map of A.8 carries a putative
collision to a collision on that fibre, which we have
just excluded. An open-curve collision maps to one on
its proper compactification. These are the required
directions of the comparison.
\end{proof}

\begin{claim}{A.16}{Proved; a degree-$16$ or degree-$32$ correspondence}
In A.13 one may replace the target by
$E_\dagger:v^2=t^3-t$, the characteristic assumption by
$\operatorname{char}\Omega\ne2$, and the cover degree by $16$ or
$32$. The geometric Galois closure then has $2$-power order.
Consequently A.14 and A.15 also hold in every characteristic
$p\equiv3\pmod4$, including $p=3$.
\end{claim}
\begin{proof}
Put the elliptic curve with prescribed branch origin in Legendre
form $E:y^2=x(x-1)(x-\lambda)$, where $\lambda\ne0,1$, and
write $\pi=x$. Choose $a,b\in\Omega$ with
$a^2=\lambda$, $b^2=1-\lambda$. The duplication identities are
\[
\begin{aligned}
 x([2]T)&=\frac{(x(T)^2-\lambda)^2}
                   {4x(T)(x(T)-1)(x(T)-\lambda)},\\
 x([2]T)-1&=\frac{(x(T)^2-2x(T)+\lambda)^2}
                   {4x(T)(x(T)-1)(x(T)-\lambda)}.
\end{aligned}
\]
They follow by substituting the tangent slope in the Weierstrass
addition law. Thus the four values
\[
 B_\dagger=\{a,-a,1+b,1-b\}
\]
are the images of four opposite pairs of exact four-torsion
points: halves of $(0,0)$ and $(1,0)$. These values are distinct.
Indeed $a,b\ne0$, and an equality $\pm a=1\pm b$ would imply
$a=0$ or $a=\pm1$ after squaring, contrary to $\lambda\ne0,1$.
They are unramified values of $\pi$. Their ordered cross-ratio is
\[
 \frac{(a-1-b)(-a-1+b)}{(a-1+b)(-a-1-b)}
       =\frac{2ab}{-2ab}=-1,
\]
using $a^2+b^2=1$. A projective transformation therefore sends
these branch values to $\{\infty,0,1,-1\}$. The double cover
branched at $B_\dagger$ is geometrically isomorphic to
$E_\dagger$.

Pull $g:D\to E$ back by $[4]$. The resulting $D_4$ is connected:
the ramified quadratic extension $\Omega(D)/\Omega(E)$ cannot
be contained in the unramified extension defined by $[4]$;
since it is quadratic, their intersection is trivial. Its map
to $D$ is an \'etale $E[4]$-torsor of degree sixteen, and its
map to $E$ is ramified above every point of $E[4]$.
Pulling back the double cover branched at $B_\dagger$ and
normalizing therefore gives an \'etale cover of $D_4$ of
degree one or two. Its Galois closure over $D$ is an extension
of a subgroup of $(\Z/2)^{16}$ by $E[4]$, hence a $2$-group.
The argument using the quotient by inversion gives exactly
the reflection implication (A.10).

If $p\equiv3\pmod4$, the quadratic-character terms for $t$ and
$-t$ in the count of $E_\dagger(\F_p)$ cancel. Hence the
count is $p+1$, its Frobenius satisfies $F^2=[-p]$, and
$E_\dagger(\overline{\F}_p)[p^\infty]=0$. To repeat A.14,
first pull back the ramified double cover by $[4]$, choose two
branch origins whose difference has exact order four, and
apply the correspondence with both origins. The two reflection
alternatives would imply twice this difference is zero, which
is false. This proves the elliptic-image conclusion. The
induction of A.15 only uses characteristic different from two,
so it then proves the hyperelliptic conclusion in the asserted
additional characteristics. The exact duplication and
cross-ratio identities are independently checked in
\texttt{code/A\_hesse\_correspondence\_checks.py}.
\end{proof}

\subsection{Removing the congruence restriction in odd characteristic}

\begin{claim}{A.17}{Proved; a prime-to-$p$ torsion cross-ratio configuration}
Let $\Omega=\overline{\F}_p$ with $p\ne2$, let $E/\Omega$ be
an elliptic curve with origin $O$, and let
$\pi:E\to\mathbf P^1$ be its quotient by inversion. For every
$c\in\Omega\setminus\{0,1\}$ there exist points
$T_1,T_2,T_3,T_4\in E(\Omega)$ of order prime to $p$ such that
none is two-torsion, their images under $\pi$ are distinct,
and the ordered cross-ratio of those images is $c$.
\end{claim}
\begin{proof}
We state the torsion-density theorem being used. If $A$ is an
abelian variety over $\overline{\F}_p$ and $Z\subset A$ is an
irreducible reduced closed subscheme not contained in a
translate of a proper abelian subvariety, then the points of
$Z$ of order prime to $p$ are Zariski dense in $Z$. This is the
contrapositive of the Anderson--Indik consequence stated in
\citet[Proposition 2.3, p.~5]{A_Yang2016}, with its reference to
\citet[p.~704]{A_Tamagawa2004}. We have opened the former text;
the latter is an indirect reference whose original PDF could
not be retrieved. The underlying surjectivity of the primary
projection for generating curves is also stated and proved in
\citet[Theorem 1.1]{A_BogomolovTschinkel2005}. We now verify all
geometric hypotheses of the density statement in this
application.

On $(\mathbf P^1)^4$ use homogeneous coordinates $(X_i:Y_i)$
and put $\Delta_{ij}=X_iY_j-X_jY_i$. Let
\[
 H_c:\quad\Delta_{13}\Delta_{24}
                   -c\Delta_{14}\Delta_{23}=0.
\]
On the locus of four distinct coordinates this is the orbit
of $(\infty,0,1,c)$ under $\operatorname{PGL}_2$; the unique
transformation taking the first three entries to a given
ordered triple determines the fourth. Consequently this open
locus is integral of dimension three. Its complement in $H_c$
is a union of triple-diagonals of dimension at most two: if
one pair of coordinates agrees, the equation and $c\ne0,1$
force a third to agree. Every component of the displayed
nonzero Cartier divisor has dimension three. It follows that
$H_c$ is integral, with line bundle
$\mathcal O(1,1,1,1)$; generic reducedness on the open locus
and the absence of embedded components of a hypersurface
justify the scheme-theoretic assertion.

Let $Z=(\pi^4)^{-1}(H_c)\subset E^4$. We verify irreducibility,
including the boundary, rather than assuming that a finite
pullback preserves it. On the open copy of
$\operatorname{PGL}_2$, the four double covers are defined by
square classes $f(M(t_i))$, where $E$ is written $y^2=f(x)$,
$M\in\operatorname{PGL}_2$ and
$(t_1,t_2,t_3,t_4)=(\infty,0,1,c)$. Fix any finite branch
value $b$ of $\pi$. For each $i$, the divisor
\[
 D_{i,b}=\{M:M(t_i)=b\}
\]
is irreducible: it is a coset of a Borel subgroup. Evaluation
at $t_i$ is smooth, so $f(M(t_i))$ has odd valuation there.
At its generic point the other three functions have valuation
zero, since their arguments avoid the finite branch set.
These four valuations show that the square classes are
independent. Their compositum is therefore a field of degree
$16$ over $\Omega(H_c)$.

The morphism $\pi^4$ is finite flat of degree sixteen. Its
pullback to the integral $H_c$ is finite flat and hence
torsion-free over every affine domain in $H_c$. Each such
coordinate algebra embeds in its generic fibre, the field
just identified. Thus the entire pullback $Z$ is integral;
there are no extra components supported on the diagonals.
It is an ample Cartier divisor in $E^4$, because its line
bundle is the pullback of the ample
$\mathcal O(1,1,1,1)$ under a finite morphism.

An ample integral divisor in this four-dimensional abelian
variety cannot lie in a translate $u+B$ of a proper abelian
subvariety. Such a $B$ would have dimension three, and
$Z=u+B$. Its associated line bundle would be pulled back
from the elliptic quotient $E^4/B$, hence would be trivial
on the positive-dimensional fibres; it would not be ample.
The torsion-density theorem applies to $Z$.

The locus where the four images are distinct and every
$T_i$ avoids $E[2]$ is a nonempty open subset of $Z$. The
finite-flat description, or a generic projective
transformation, proves its nonemptiness. A prime-to-$p$
torsion point of $E^4$ in this open has exactly the four
coordinates asserted. In particular they have a common
finite exponent $n$ with $(n,p)=1$.
\end{proof}

\begin{claim}{A.18}{Proved; supersingular correspondences in every odd characteristic}
Let $\Omega=\overline{\F}_p$ with $p\ne2$, and let
$g:D\to E$ be a ramified degree-two morphism of smooth proper
connected curves with $g(E)=1$. For each chosen branch origin
$O\in E(\Omega)$ there are a connected finite \'etale cover
$a:W\to D$ whose geometric Galois closure has order prime
to $p$, a supersingular elliptic curve $E_0/\Omega$, and a
morphism $r:W\to E_0$ satisfying the reflection implication
(A.10). All data descend to a finite constant extension of
a field of definition of $g$.

Over every finite field of odd characteristic, equal
geometrically prime-to-$p$ point-section classes on $D$
therefore have equal images under $g$.
\end{claim}
\begin{proof}
A supersingular elliptic curve exists in every odd
characteristic. One explicit way to choose one is to take
$c$ a root of the nonconstant polynomial
\[
 H_p(T)=\sum_{i=0}^{m}\binom mi^2T^i,
 \qquad m=(p-1)/2,
\]
and put $E_0:v^2=t(t-1)(t-c)$. The Legendre
Hasse-invariant criterion is the one stated in
\citet[\S1, equations (1.2)--(1.3)]{A_Finotti2009}.
Its use here is valid for every $p>2$, including three.
The root cannot be zero because $H_p(0)=1$, nor one
because Vandermonde's identity gives
$H_p(1)=\binom{p-1}{m}\equiv(-1)^m\pmod p$.
Thus the displayed curve is smooth, and
$E_0(\Omega)[p^\infty]=0$.

Apply A.17 to $(E,O)$ and this $c$. Let $n$ be a common
prime-to-$p$ exponent of the four resulting torsion points.
The set $B=\{\pi(T_i):1\le i\le4\}$ consists of four
unramified values of $\pi$. A projective transformation
identifies its ordered configuration with
$(\infty,0,1,c)$, so the double cover of $\mathbf P^1$
branched at $B$ is isomorphic over $\Omega$ to $E_0$.

Form $D_n=D\times_{g,E,[n]}E$. It is connected even if
$n$ is even: a ramified quadratic extension cannot be a
subextension of the \'etale Galois cover $[n]$. Thus the
quadratic extension and the $[n]$ extension are linearly
disjoint. The map $D_n\to D$ is an \'etale $E[n]$-torsor
of degree $n^2$, while $D_n\to E$ is ramified with index
two above every point of $E[n]$. Every point of
$\pi^{-1}(B)$ belongs to $E[n]$, since these fibres are
the pairs $\{T_i,-T_i\}$. The normalized pullback of the
double cover branched at $B$ is consequently \'etale over
$D_n$. Choose a connected component $W$. Its degree over
$D_n$ is one or two, and its Galois closure over $D$ has
group an extension of a subgroup of $(\Z/2)^{n^2}$ by
$E[n]$. Its order is prime to $p$.

The projection $r:W\to E_0$ has the same reflection property
as in A.13: equal images imply that the two intermediate
points of $E$ are equal or negatives; applying $[n]$ yields
(A.10). All equations, torsion points and chosen components
are defined over a finite subfield of $\Omega$.

For the final assertion, start with a hypothetical equality
of the section classes of $P,Q\in D(k)$. Choose a prime
$\ell\ne2,p$ and first pull back $g$ by $[\ell]$. By A.3,
after a finite extension, there are lifts $P_\ell,Q_\ell$
with equal full prime-to-$p$ section classes. A fibre over
a branch point of $g$ supplies two branch origins of this
new double cover whose difference has exact order $\ell$.
Apply the first assertion with each origin separately.
Their lifts have $p$-primary divisor difference by A.1;
their images on the corresponding supersingular elliptic
curves coincide. The two reflection alternatives therefore
force equal intermediate images on $E$, as in A.14,
since twice an element of exact order $\ell$ is nonzero.
Applying $[\ell]$ gives $g(P)=g(Q)$. Restrictions to finite
constant extensions and lifts of actual equal section
classes are precisely the comparison maps proved in A.3.
\end{proof}

\begin{claim}{A.19}{Proved; A1 for every hyperelliptic curve in odd characteristic}
Let $p\ne2$ and let $K$ be a field finitely generated over
$\F_p$. Let $C/K$ be smooth proper geometrically connected,
geometrically hyperelliptic, and of genus at least two.
The point-section map
\[
 C(K)\longrightarrow
 \{\text{sections of }E_C^{(p')}\to G_K\}/\Pi_C^{(p')}
\]
is injective. The same holds for every hyperbolic open
subcurve of $C$. In particular A1 is proved for every
genus-two curve in every odd characteristic, with no
restriction on the finite field or on its size.
\end{claim}
\begin{proof}
Over finite fields, repeat the proof of A.15 using the
elliptic-image assertion of A.18. All covers in that
induction are \'etale double covers or the covers of A.18,
so their geometric Galois closures remain prime to $p$.
The branch sets can be split and the chosen lifts made
rational after finite constant extensions, which preserve
a hypothetical section equality. The disjoint-branch
fibre product is smooth, and its two projections separate
all geometric points, exactly as proved there. This
settles every hyperelliptic genus by the same induction.
The specialization argument in the last paragraph of
A.15, using A.8, gives the assertion over $K$; functoriality
to the proper compactification gives it for open subcurves.
For these curves over finite fields the exceptional set of
\citet[Proposition 1.8]{A_SaidiTamagawa2009} is therefore empty:
commensurable decomposition groups become equal after
restriction to an open subgroup of $G_k$, which would give
equal point-section classes over a finite constant extension.
\end{proof}

\begin{claim}{A.20}{Proved; elementary abelian two-group covers of a genus-zero curve}
Let $K$ be finitely generated over $\F_p$, $p\ne2$, and let
$C/K$ be smooth proper geometrically connected of genus at
least two. Suppose that $C_{\overline K}$ admits a faithful
action of $(\Z/2)^r$, $r\ge1$, with genus-zero quotient.
Then the geometrically prime-to-$p$ point-section map is
injective for $C$ and for each of its hyperbolic open
subcurves. This includes nonhyperelliptic curves, for example
the smooth plane quartic $X^4+Y^4+Z^4=0$ in every odd
characteristic: changing coordinate signs gives a Klein-four
group, with genus-zero quotient $U^2+V^2+W^2=0$.
\end{claim}
\begin{proof}
We first work over a finite field, extending it to define the
action and all quotients. We induct on $r$. For $r=1$ the
curve is hyperelliptic and A.19 applies. Suppose $r\ge2$,
and that $P,Q$ have equal point-section classes. For each
nonidentity involution $s$, consider the degree-two quotient
\[
 q_s:C\longrightarrow C/\langle s\rangle.
\]
If any quotient has genus zero, $C$ is hyperelliptic and we
are done. If it has genus one, the map is ramified by
Riemann--Hurwitz, since $g(C)\ge2$, and A.18 gives
$q_s(P)=q_s(Q)$. If its genus is at least two, the quotient
has the faithful action of
$(\Z/2)^r/\langle s\rangle$ with the same genus-zero final
quotient. Functoriality of section classes and the induction
hypothesis again give $q_s(P)=q_s(Q)$.

These equalities force $P=Q$ on every fibre, including all
ramified fibres. Indeed choose independent involutions $s,t$.
If $P\ne Q$, equality in the three quotients for $s,t,st$
would give
\[
 Q=sP=tP=(st)P.
\]
But $sP=tP$ implies $(st)P=P$, contradicting the third
equality. This proves the finite-field assertion over all
constant extensions.

For the asserted finitely generated field, first take a
finite extension defining the action. Spread the smooth
proper curve, the faithful action and its quotient maps over
a finite-type base over $\F_p$. After shrinking, their genera
are constant and the action remains faithful; its order is
invertible and the quotients are smooth curves. A hypothetical
collision specializes by A.8 to one on a closed finite-field
fibre in the same class, which has just been excluded.
Functoriality to the compactification proves the open case.
\end{proof}

\begin{claim}{A.21}{Proved; a general elliptic-image test}
Let $k$ be finite of characteristic $p$, with no restriction
on $p$, and let
$g:D\to E$ be a finite separable morphism of smooth proper
geometrically connected curves, with $g(E)=1$. Suppose that
there is a geometric point $b\in E$ and an integer $m>1$
prime to $p$ such that every local ramification index over
$b$ is divisible by $m$. Equal geometrically prime-to-$p$
point-section classes on $D$ then have equal images under
$g$, over every finite constant extension.
\end{claim}
\begin{proof}
Choose a prime $e\mid m$. We adapt the explicit construction
of \citet[\S3.1]{A_Poonen2005}, while verifying prime-to-$p$
Galois closure and keeping the point-image information.
Take $b$ as elliptic origin and choose a prime
$\ell\ge5$, $\ell\ne p$. A connected component $D_\ell$ of
$D\times_{g,E,[\ell]}E$ is an \'etale cover of $D$ whose
geometric Galois closure has order dividing $\ell^2$: the
whole pullback is an $E[\ell]$-torsor and a connected
component is a torsor under its stabilizer. The second
projection $D_\ell\to E$ is surjective and has all local
ramification indices divisible by $e$ above every point of
$E[\ell]$.

Write $x:E\to\mathbf P^1$ for the inversion quotient.
Every nonzero $\ell$-torsion point has trivial stabilizer
under inversion, since $\ell$ is odd. Its image is
therefore an unramified fibre even in characteristic two.
Choose distinct values $a_i$ arising from these pairs.
When $p\ne2$ and $e\ge3$, use four such values and
let
\[
 H:\quad z^e=\frac{(x-a_1)(x-a_2)}{(x-a_3)(x-a_4)}.
\]
It has genus $e-1\ge2$, and its map to the $z$-line has
degree two; hence it is hyperelliptic. For $e=2$, which
necessarily has $p\ne2$, use six
values and the genus-two curve
\[
 H:\quad z^2=
 \frac{(x-a_1)(x-a_2)(x-a_3)}{(x-a_4)(x-a_5)(x-a_6)}.
\]
When $p=2$, the prime $e$ is odd. Use three chosen values,
move them to $0,1,\infty$, and take instead
\[
 H:\quad z^e=x(x-1),\qquad\text{equivalently }x^2+x=z^e.
\]
We verify the required point injectivity for this target
directly, including the elliptic case $e=3$. Its genus is
$g=(e-1)/2$ by tame Riemann--Hurwitz for the degree-$e$
map to the $x$-line. At the unique point at infinity,
$v_\infty(z)=-2$ and $v_\infty(x)=-e$. Since $e$ is odd,
the leading Laurent term gives $v_\infty(dx)=-e-1$.
Differentiation gives $dx=z^{e-1}dz$, hence
$v_\infty(dz)=e-3=2g-2$. The affine curve is \'etale over
the $z$-line, so $z^i\,dz$, $0\le i\le g-1$, form a basis
of regular differentials. Cartier kills the terms with
even $i$ and sends those with odd $i$ to
$z^{(i-1)/2}\,dz$, with square-root semilinearity on
constants. The operator is therefore nilpotent.

There is consequently no nonzero geometric two-torsion
in $J_H$. Indeed such a class represented by a divisor $D$
would give $\operatorname{div}(f)=2D$ with $f$ not a square;
then $df/f$ would be a nonzero regular Cartier-fixed
differential. Regularity follows locally by writing
$f=t^{2a}u$ with $u$ a unit. To check the fixed identity
directly in the separating variable $z$, write
$f=a^2+zb^2$. Then
\[
 \frac{df}{f}=\frac{b^2}{f}\,dz
 =\left(\left(\frac{ab}{f}\right)^2
       +z\left(\frac{b^2}{f}\right)^2\right)dz,
\]
whose Cartier image is itself. Thus
$J_H(\overline k)[2^\infty]=0$, and A.1 and the Abel map
prove point-section injectivity on $H$, for every finite
constant extension. This also treats $e=3$, when $H$ is
an elliptic curve.

In every case, each local valuation of the defining
Kummer function on
$D_\ell$ is divisible by $e$. A connected component of the
normalized pullback is therefore \'etale over $D_\ell$,
of degree one or $e$. The geometric Galois closure of its
composite cover of $D$ has order divisible only by
$\ell,e$, by adjoining its deck conjugates. Both primes
are different from $p$.

Lift equal section classes to this cover using A.3. Their
images on $H$ are equal by A.19 when $p\ne2$ and by the
Cartier calculation just given when $p=2$. Consequently their
intermediate elliptic images have the same $x$-coordinate,
and so differ at most by inversion. Applying $[\ell]$
gives the same reflection alternative (A.10) with origin
$b$. No connectedness of the entire isogeny pullback was
assumed in this argument.

To remove reflection, first pull back the original $g$ by
$[q]$, for a prime $q\ne2,p$, and fix a connected component.
Lift the equal sections there. Every point of $[q]^{-1}b$
is a fibre all of whose ramification indices remain
divisible by $m$. Choose two such points whose difference
has exact order $q$. Apply the construction with each as
origin to the same pair of lifted points. The two
reflection alternatives would force this difference to
be killed by two. Hence the intermediate elliptic images
are equal; applying $[q]$ proves the assertion for $g$.
All finite data and the chosen lifts are made rational by
finite constant extensions, which preserve section equality.
\end{proof}

\begin{claim}{A.22}{Proved; cyclic prime-degree covers of a genus-zero curve}
Let $K$ be finitely generated over $\F_p$, with no restriction
on the characteristic $p$.
Let $C/K$ be smooth proper geometrically connected of genus
at least two. Suppose that, over $\overline K$, it admits a
cyclic cover of prime degree $\ell\ne p$ of a genus-zero
curve. Then its geometrically prime-to-$p$ point-section map
is injective. The conclusion also holds for every hyperbolic
open subcurve of $C$.
\end{claim}
\begin{proof}
We first work over a finite field and allow finite constant
extensions to define all covers and chosen points. The case
$\ell=2$ necessarily has $p\ne2$ and is A.19. In the rest
of the proof $\ell$ is odd.
Every branch fibre of a cyclic $\ell$-cover is totally
ramified, because $\ell$ is prime. If there are $r$ branch
points, tame Riemann--Hurwitz gives
\[
 g(C)=\frac{(\ell-1)(r-2)}2.                                  \tag{A.13}
\]

For $p=2$, the curve $H:z^\ell=x(x-1)$ has
$J_H(\overline k)[2^\infty]=0$, as proved by the Cartier
calculation in A.21. Hence it has point-section
injectivity for every odd $\ell$, including the elliptic
case $\ell=3$.

We first prove a fibre statement, independently of the
induction below: a collision on such a cover has equal
images on its genus-zero base. If $\ell\ge5$, select three
branch points and choose a coordinate sending them to
$0,1,\infty$. The curve
\[
 H:\quad z^\ell=x(x-1)
\]
is a cyclic $\ell$-cover branched at exactly these three
points. It is hyperelliptic: the function $z$ gives a
degree-two map, and its genus is $(\ell-1)/2\ge2$.
If $C$ and $H$ define the same extension of the rational
function field, injectivity follows from A.19 in odd
characteristic and from the preceding Cartier calculation
in characteristic two.
Otherwise their extensions are linearly disjoint, since
both have prime degree and are Galois. The normalization
$V$ of their fibre product is a connected cyclic \'etale
$\ell$-cover of $C$. Tameness and the matching index
$\ell$ at every branch point of $H$ prove \'etaleness,
including over infinity. Equal sections lift to $V$ by
A.3 and have equal images on $H$ by the just-established
injectivity for that target. Their original
$x$-coordinates are therefore equal.

For $\ell=3$, use the same three-branch curve $H$, now of
genus one. Formula (A.13) and $g(C)\ge2$ imply $r\ge4$.
Thus $C$ and $H$ are distinct covers, and the normalized
fibre product $V$ is again a cyclic \'etale cubic cover of
$C$. In characteristic two the target is the
two-rank-zero elliptic curve just treated, so the lifted
images on $H$ are equal. In odd characteristic its cyclic
cubic map to $H$ is ramified over a branch
point of $C$ outside the three selected ones; all its
ramification indices there are three. A.21 gives equality
of the $H$-images of the lifted collision, hence equality
of the original genus-zero coordinates. This proves the
fibre statement for every $\ell$ under consideration.

We now induct on the number $r$ of branch points. The
induction only asserts injectivity for source curves of
genus at least two; lower-genus quotients will be treated
separately. Write the given Kummer extension as
\[
 \overline k(C)=\overline k(x)(y),\qquad y^\ell=c(x).
\]
Choose two branch points $a_1,a_2$. Let $u\not\equiv0$
modulo $\ell$ be the valuation of $c$ at $a_1$. Choose a
rational function $b(x)$ with divisor
$u(a_1-a_2)$, choosing an integer representative for $u$,
and define the cyclic cover
\[
 B:\quad v^\ell=b(x).
\]
It has exactly two branch points and genus zero. Its
function field differs from that of $C$, because $C$ has
at least three branch points. The normalized fibre
product $V$ of $C$ and $B$ is consequently connected and
is a cyclic \'etale $\ell$-cover of $C$.

Put $h(x)=c(x)/b(x)$ and let $H_2$ be the smooth curve
with Kummer equation $w^\ell=h(x)$. The branch at $a_1$
cancels. The branch at $a_2$ either cancels or remains,
and every other branch remains. Thus the number of its
branch points is $r-1$ or $r-2$, and it is nontrivial,
since at least one branch different from $a_1,a_2$
remains. In particular $g(H_2)<g(C)$. On function fields
\[
 \overline k(V)=\overline k(x)(y,v)
              =\overline k(x)(v,w),\qquad w=y/v.                \tag{A.14}
\]
The two Kummer characters $b,h$ are independent; otherwise
$c=b h$ would also be a power of $b$ modulo $\ell$-th
powers and could have no branch outside $\{a_1,a_2\}$.
Thus $V\to B$ and $V\to H_2$ are cyclic covers of degree
$\ell$.

Suppose that distinct $P,Q\in C(k)$ form a collision.
The fibre statement gives a common base value
$x(P)=x(Q)$. This value is not a branch point of $C$,
since a branch fibre contains only one geometric point.
Lift the collision to $P_V,Q_V\in V$ by A.3 after a finite
constant extension. Applying the already proved fibre
statement to the cyclic cover $V\to B$, whose source has
genus $\ell(g(C)-1)+1\ge2$, gives equal $B$-images.

If $g(H_2)\ge2$, induction gives equal $H_2$-images. If
$g(H_2)=1$ and $p$ is odd, the map $V\to H_2$ is fully ramified above
every point over $a_1$: the cover $H_2$ is unramified
there, while $B$ has index $\ell$. A.21 therefore gives
equal $H_2$-images. If $g(H_2)=1$ and $p=2$, formula
(A.13) forces $\ell=3$ and exactly three branch points.
Their three nonzero Kummer exponents sum to zero modulo
three, so they are all one or all two. A projective
coordinate change and, if needed, inversion of the
Kummer generator identify this curve with
$z^3=x(x-1)$. The preceding Cartier calculation gives
its two-rank-zero point injectivity. If $g(H_2)=0$, apply the fibre
statement to $V\to H_2$. In every case both images are
equal.

It remains to check point separation by (A.14), rather
than infer it from a birational function-field identity
at a possibly singular point. Off the branch set of $C$,
both $B$ and $H_2$ are \'etale over the common genus-zero
base. Their fibre product is therefore smooth and normal
there, and agrees with its normalization $V$. Both
$P_V$ and $Q_V$ lie in this locus. Equal images in $B$ and
$H_2$ imply $P_V=Q_V$, contradicting their distinct images
on $C$. This closes the induction, including its
lowest-branch cases, for which $H_2$ has genus zero or one.

Finally spread the cyclic action, the quotient and its
distinct tame branch points over a finite-type model of
$K$, after a finite extension defining them. On a
sufficiently small smooth locus this remains a cyclic
prime-degree cover with the same ramification data.
Specialization of a collision by A.8 contradicts the
finite-field result. Proper compactification gives the
assertion for open subcurves.
\end{proof}

\begin{claim}{A.23}{Proved; all cyclic prime-to-$p$ covers of a genus-zero curve}
Let $K$ be finitely generated over $\F_p$, and let $C/K$
be smooth proper geometrically connected of genus at least
two. Suppose that $C_{\overline K}$ admits a cyclic Galois
map of degree $n>1$, with $p\nmid n$, to a genus-zero
curve. Then its geometrically prime-to-$p$ point-section
map is injective, as is the point-section map of each
hyperbolic open subcurve. There is no restriction on $p$
or on the number of prime factors of $n$.
\end{claim}
\begin{proof}
We work first over a finite field and make finite constant
extensions whenever needed to define the covers, roots of
unity, branch points and chosen lifts. We induct on $n$;
the prime cases have also been proved separately in A.22.
Only curves of genus at least two enter the induction
assertion. A genus-one quotient is treated by A.21, and a
genus-zero quotient supplies a new rational base.

Suppose first that $n=ab$ with $a,b>1$ coprime. Let
$C_a,C_b$ be the intermediate quotient curves of degrees
$a,b$ over the given genus-zero base. Their function fields
have compositum $\overline k(C)$. Let $P,Q$ be a collision.
If either quotient has genus zero, $C$ itself is a cyclic
cover of that quotient of degree $b<n$ or $a<n$, so the
outer induction already gives $P=Q$. Otherwise a quotient
of genus at least two has equal images by the induction
hypothesis. A quotient of genus one also has equal images:
the Galois map from $C$ is ramified by Riemann--Hurwitz,
and every index in a branch fibre is the same nontrivial
divisor of its prime-to-$p$ degree. Thus A.21 applies.

The maps to $C_a,C_b$ jointly distinguish all geometric
points, including points over their common branch values.
To verify this last assertion, take a point of the fibre
product and let $e_a,e_b$ be the two local ramification
indices. They divide the coprime integers $a,b$. Over the
algebraically closed residue field, the completed tame
local product has equation $u^{e_a}=v^{e_b}$, after
absorbing units into the parameters. Its normalization
has exactly $\gcd(e_a,e_b)=1$ point above the selected
point. Since the function fields have the asserted
compositum, $C$ is this normalization and the joint map is
injective on geometric points. Hence $P=Q$.

It remains to treat $n=\ell^a$. Let $C_0$ be the unique
degree-$\ell$ quotient over the genus-zero base. In a
Kummer presentation
\[
 \overline k(C)=\overline k(x)(y),\qquad y^n=c(x),
\]
the branch points of $C_0$ are exactly the points where
the valuation of $c$ is not divisible by $\ell$. These
are exactly the points at which $C$ has full ramification
index $n$. Write $r_{\mathrm{full}}$ for their number. If
$g(C_0)=0$, the cyclic map $C\to C_0$ has degree
$n/\ell<n$, and the outer induction settles $C$.
Otherwise $r_{\mathrm{full}}\ge3$ when $\ell$ is odd and
$r_{\mathrm{full}}\ge4$ when $\ell=2$.

We next prove a fibre statement for degree $n$, using only
the outer induction on smaller degrees: every collision
on $C$ has equal images on the genus-zero base. The case
$g(C_0)=0$ is already settled. If $\ell\ge5$, choose three
full-index branch points and take the cyclic
degree-$\ell$ cover
\[
 H:\quad z^\ell=x(x-1)
\]
after moving those three points to $0,1,\infty$. Its genus
is $(\ell-1)/2$, and it has point-section injectivity by
A.19 for odd $p$ and by the Cartier calculation in A.21
for $p=2$. If $\ell=3$, take the same three-branch cover,
which now has genus one. If $\ell=2$, take the elliptic
double cover branched at four chosen full-index points;
here $p$ is necessarily odd.

Choose a connected component $V$ of the normalized fibre
product of $C$ and $H$. It is an \'etale cyclic cover of
$C$ of degree one or $\ell$. Indeed $H$ is branched only
at the selected points, and its tame ramification index
$\ell$ divides the full index $n$ on $C$; tame base change
cancels the ramification. If $H$ has genus at least two,
a lifted collision has equal $H$-images by its just-cited
injectivity. If $H$ has genus one, the map $V\to H$ is
Galois cyclic, of degree $n$ or $n/\ell$. Its source has
genus at least two since $V\to C$ is \'etale. Thus the
map is ramified, and a branch fibre has a common
nontrivial prime-to-$p$ ramification index. A.21 again
gives equal $H$-images. The original base coordinates are
equal in either case. This proves the fibre statement
without using the inner induction that follows.

For this fixed $n=\ell^a$, we now induct on
$r_{\mathrm{full}}$, with the cases $g(C_0)=0$ provided
by the outer induction. Choose two full-index branch
points $a_1,a_2$. Put $u=v_{a_1}(c)$, which is a unit
modulo $\ell$, and choose $b(x)$ with
\[
 \operatorname{div}(b)=u(a_1-a_2).
\]
The curve $B:v^n=b(x)$ has genus zero and full
ramification at exactly these two points. Its function
field is linearly disjoint from that of $C$. In fact any
nontrivial intersection of the cyclic $\ell$-power
extensions would contain their degree-$\ell$ subfields;
the latter cannot agree, since $C_0$ has at least three
branch points while the corresponding quotient of $B$
has exactly two.

Consequently the normalization $V$ of $C\times_{\mathbf
P^1}B$ is connected and is a cyclic \'etale cover of $C$
of degree $n$. Matching full tame indices cancel
ramification over $a_1,a_2$, and $B$ is unramified
elsewhere. Define $H_2$ by
\[
 w^n=h(x),\qquad h(x)=c(x)/b(x).
\]
At $a_1$ the valuation is zero. At $a_2$ full ramification
may remain, may become partial, or may disappear; no new
branch point occurs. Every other full-index branch point
remains. There is at least one such point, so $h$ has
exact Kummer order $n$ and $H_2$ is a connected cyclic
degree-$n$ cover. Its number of full-index branch points
is strictly smaller. The invertible change of Kummer
generators $w=y/v$ gives
\[
 \overline k(V)=\overline k(x)(y,v)
              =\overline k(x)(v,w).                       \tag{A.15}
\]
The two latter characters are independent: the change
$(c,b)\mapsto(b,c/b)$ is an invertible change of basis of
$(\Z/n)^2$. Thus both maps $V\to B$ and $V\to H_2$ are
cyclic of degree $n$.

Suppose $P\ne Q$ is a collision on $C$. The fibre
statement gives a common base value. It is different from
$a_1,a_2$, since each of those full-index fibres consists
of one geometric point. After a finite constant extension,
lift the collision to $P_V,Q_V$ on $V$ by A.3. Apply the
degree-$n$ fibre statement to $V\to B$: its source has
genus $n(g(C)-1)+1\ge2$, so the $B$-images agree. If
$g(H_2)\ge2$, the inner induction, or its outer-induction
base case, gives equal $H_2$-images. If $g(H_2)=1$, the
map $V\to H_2$ is fully ramified of index $n$ over every
point above $a_1$, because $H_2$ is unramified there and
$B$ has index $n$. A.21 gives the required equality.
If $g(H_2)=0$, apply the already proved degree-$n$ fibre
statement to $V\to H_2$.

These two equalities imply equality of the lifted points.
Indeed outside $\{a_1,a_2\}$ the map $B\to\mathbf P^1$
is \'etale, so $B\times_{\mathbf P^1}H_2$ is \'etale over
the smooth curve $H_2$ and is already normal. By (A.15)
it agrees there with $V$. Both lifted points lie in this
locus. Equal $B$- and $H_2$-images therefore give
$P_V=Q_V$, contrary to their distinct images on $C$.
Notice that no claim of total ramification at the other
branch points was used. This completes both inductions.

The same spreading and specialization as in A.22 applies
to the finite tame cyclic action of order $n$; after
shrinking the model, the quotient and its ramification
indices persist. A.8 therefore gives the assertion over
$K$. Functoriality to the proper compactification gives
the assertion for open subcurves.
\end{proof}

\begin{claim}{A.24}{Proved; all abelian prime-to-$p$ Galois covers of a genus-zero curve}
Let $K$ be finitely generated over $\F_p$. Let $C/K$ be
smooth proper geometrically connected of genus at least
two. Suppose that $C_{\overline K}$ has a faithful action
of a finite abelian group $G$ of order prime to $p$, with
genus-zero quotient. Then its geometrically prime-to-$p$
point-section map is injective. The same holds for every
hyperbolic open subcurve. This theorem includes A.19,
A.20 and A.23, but also allows noncyclic abelian groups
with arbitrary prime-to-$p$ exponents in every characteristic.
\end{claim}
\begin{proof}
Over a finite field, after defining the action and its
quotients by a finite constant extension, induct on $|G|$.
The cyclic case is A.23. For a nontrivial character
$\chi:G\to\Z/n$, understood to have image of order $n$,
put $C_\chi=C/\ker\chi$. It is a cyclic degree-$n$
cover of the common genus-zero quotient. If any
$C_\chi$ has genus zero, then $C\to C_\chi$ has group
$\ker\chi$ of smaller order. The induction hypothesis
settles $C$ at once.

In the remaining case let $P,Q$ be a collision on $C$.
For each $\chi$ with $g(C_\chi)\ge2$, A.23 gives equal
$C_\chi$-images by functoriality. If $g(C_\chi)=1$, the
map $C\to C_\chi$ is ramified by Riemann--Hurwitz. It is
Galois of prime-to-$p$ order, so any branch fibre has
equal nontrivial ramification indices. A.21 gives equal
$C_\chi$-images also in this case. This argument uses
the all-characteristic form of A.21; no classification
of elliptic automorphisms in characteristic two is needed.

It remains to check that these quotient maps jointly
separate geometric points at every ramified fibre. Their
equalities first give the same point of $C/G$, so
$Q=gP$ for some $g\in G$. Let $I\subseteq G$ be the
stabilizer of $P$. Equality of the two images in
$C/\ker\chi$ is exactly the assertion
\[
 g\in(\ker\chi)I.
\]
Every character of the finite abelian group $G/I$ is a
character of $G$ trivial on $I$, and such characters
separate its elements. The displayed containment for
all $\chi$ therefore forces $g\in I$. Hence $Q=P$.

Over a finitely generated field, extend finitely to define
the action and spread the curve, faithful action and
quotient over a finite-type model. Since $|G|$ is
invertible, after shrinking the quotient is again smooth
of genus zero and the action remains faithful. A.8
specializes any collision to the finite-field case just
excluded. Passing to the compactification proves the
open assertion.
\end{proof}

\begin{claim}{A.25}{Proved; ramified double covers of elliptic curves}
Let $K$ be finitely generated over $\F_p$, with $p\ne2$.
Let $C/K$ be smooth proper geometrically connected of genus
at least two. If $C_{\overline K}$ admits a degree-two map
to a genus-one curve, then its geometrically prime-to-$p$
point-section map is injective. More generally the same
holds if $C_{\overline K}$ has a faithful elementary
abelian two-group action with genus-one quotient. The
conclusions include every hyperbolic open subcurve.
\end{claim}
\begin{proof}
First suppose that $C\to E$ is a double cover over a
finite field, allowing finite constant extensions to define
all points and covers. It is ramified because $g(C)\ge2$.
Let $S$ be its reduced branch set, of even cardinality $r$.
Riemann--Hurwitz gives $g(C)=1+r/2$. We induct on $r$.
When $r=2$, the curve has genus two and A.19 applies.

For $r\ge4$, choose distinct $a_1,a_2\in S$. There is a
double cover $B\to E$ branched exactly at these two
points: over the algebraic closure, multiplication by two
on $\operatorname{Pic}^0(E)$ is surjective, so
$\mathcal O_E(a_1+a_2)$ has a square root of degree one.
The section with zero divisor $a_1+a_2$ defines the
required smooth double cover. It is connected because it
is ramified, and it has genus two. These finite data descend
to a finite constant extension.

The quadratic extensions defining $C$ and $B$ differ,
since their branch sets differ. Their normalized fibre
product $V$ is connected and is a double \'etale cover of
$C$: quadratic ramification cancels at $a_1,a_2$, and
$B$ is unramified elsewhere. Let $H_2\to E$ be the third
quadratic quotient, whose Kummer character is the product
of the two original characters. Its branch set is exactly
$S\setminus\{a_1,a_2\}$, so it is connected and has genus
$g(C)-1\ge2$. The induction hypothesis gives point-section
injectivity for $H_2$, while A.19 gives it for $B$.

A collision on $C$ lifts to a collision on $V$ by A.3.
Both quotient images agree by these two injectivity
statements. Their equality forces equality of the points
on $V$: the branch sets of $B$ and $H_2$ over $E$ are
disjoint, so their fibre product is smooth everywhere and
equals its normalization $V$. Thus the two projections
jointly separate every geometric point. This completes
the branch-count induction without needing to infer point
separation from a function-field identity at a singular
point.

For an elementary abelian group $(\Z/2)^r$, induct on
$r$, using the double-cover result for $r=1$. For each
nontrivial involution $s$, the quotient $C/\langle s\rangle$
still maps to the genus-one final quotient. It cannot have
genus zero. If its genus is at least two, the smaller-rank
induction gives equal quotient images of any collision;
if its genus is one, A.21 gives that equality for the
ramified double map from $C$. For $r\ge2$, the three
quotients attached to independent involutions $s,t$ and
$st$ jointly separate every point, by exactly the orbit
argument in A.20. Therefore the collision is trivial.

Spread the tame map or tame finite group action and its
genus-one quotient over a finite-type model of $K$, after
a finite extension defining them. A.8 specializes a
hypothetical collision to the finite-field case just
proved. Functoriality gives the open assertion.
\end{proof}

\subsection{A family of ordinary genus-two curves in characteristic two}

\begin{claim}{A.C2.1}{Proved}
Let $K$ be finitely generated over $\F_2$, let $a,b,c\in K^\times$,
and let $C_{a,b,c}$ be the smooth proper curve with function field
\[
 y^2+y=a x+\frac b{x}+\frac c{x+1}.
\]
It is an ordinary genus-two curve. If
\[
                    (a+b+c)^3+abc=0,                         \tag{C2.1}
\]
then its geometrically prime-to-two point-section map is injective.
The same conclusion holds for every hyperbolic open subcurve and after
any finite extension of $K$. It also holds for any curve that becomes
isomorphic to such a curve after a finite extension.

Over an algebraic closure, condition (C2.1) is equivalent to
\[
 b=a v^3,\qquad c=a(v+1)^3\qquad(v\ne0,1).                    \tag{C2.2}
\]
For each such $v$ there is a connected cyclic \'etale cover $D\to C$
of degree three and two maps $D\to E$ to one supersingular elliptic
curve whose product is injective on geometric points.
\end{claim}
\begin{proof}
Work initially over an algebraic closure $\Omega$ of $K$. The three
simple Artin--Schreier poles are $0,1,\infty$. Their different
exponents are two, so Riemann--Hurwitz gives genus two. Write
\[
 h=x^2+x,\qquad Y=h y,\qquad
 f=a x^5+(a+b+c)x^3+c x^2+b x.
\]
The equation is $Y^2+hY=f$. There is one point $P_\infty$ above
infinity, with pole orders two and five for $x$ and $Y$, respectively.
The regular differentials $dx/h$ and $x\,dx/h$ form a basis. Cartier
fixes each of them: write $1/h=h/h^2$ and use
$\mathcal C((x^2+x)dx)=dx$ and
$\mathcal C((x^3+x^2)dx)=x\,dx$ in characteristic two.
Thus Cartier is bijective and the curve is ordinary.

For the parameterization put $B=b/a$, $C=c/a$. Equation (C2.1)
becomes $(1+B+C)^3+BC=0$. If $B+C\ne0$, set
$v=(C+1)/(B+C)$. Direct substitution gives
$v^2+v=B+C+1$ and $v^3=B$, whence $(v+1)^3=C$.
If $B=C$, the equation implies $B=C=1$; either root of
$v^2+v+1=0$ works. Conversely (C2.2) gives (C2.1) by expansion.
Nonvanishing of $b,c$ excludes $v=0,1$.

Choose $d\in\Omega$ satisfying $d^2+d=a^2v^2+a$, and set
\[
 u=x^2+v^2,\qquad
 A=a x^3+d x^2+(d+a+b+c)x+b.
\]
A coefficient calculation gives
\[
                  A^2+hA+f=a^2u^3.                          \tag{C2.3}
\]
Let $F=(A+Y)/a$. The hyperelliptic involution $\iota$ sends
$Y$ to $Y+h$, so $F\iota(F)=u^3$. At $x=v$, the map
$C\to\mathbf P^1$ is unramified, since $h(v)=v(v+1)\ne0$.
Of its two points, precisely one, denoted $R$, has $F(R)=0$.
The conjugate value is $h(v)/a\ne0$. Because $u$ has order two
there, (C2.3) gives
\[
                 \operatorname{div}(F)=6R-6P_\infty.
\]
The normalization in $z^3=F$ is therefore \'etale. It is connected:
if $F$ were a cube, its cube root would have its only pole at
$P_\infty$ of order two and would lie in
$L(2P_\infty)=\langle1,x\rangle$. Its cube could not contain the
nonzero $Y$ coefficient of $F$. Denote this cyclic cover by $D$;
it has genus four.

Put $w=z+u/z$. The involution $z\mapsto u/z$ gives a map of
degree two from $D$ to
\[
 E:\quad a\bigl(W^3+(X^2+v^2Z^2)W\bigr)=X^2Z+XZ^2,
 \qquad O=[1:0:0].                                          \tag{C2.4}
\]
The derivatives with respect to $X$ and $Z$ are $Z^2$ and
$X^2$. A singular projective point would therefore have $X=Z=0$,
and the derivative with respect to $W$ would then force $W=0$.
Thus this is a smooth plane cubic.

We verify supersingularity without a characteristic restriction beyond
two. For a smooth plane cubic $G(X,W,Z)=0$ in characteristic two,
the scalar of Frobenius on $H^1(E,\mathcal O_E)$ is the coefficient
of $XWZ$ in $G$, in the usual one-dimensional Cech description.
Indeed the exact sequence
$0\to\mathcal O_{\mathbf P^2}(-3)\xrightarrow{G}
\mathcal O_{\mathbf P^2}\to\mathcal O_E\to0$
identifies this group with $H^2(\mathcal O_{\mathbf P^2}(-3))$,
whose generator is represented by $1/(XWZ)$. Frobenius takes it
to $G/(X^2W^2Z^2)$; every cubic monomial except $XWZ$ gives a
zero class in this Cech group. In (C2.4) that coefficient is zero.
Hence the elliptic curve has $2$-rank zero, and
$E(\Omega)[2^\infty]=0$.

Let $\zeta\ne1$ satisfy $\zeta^3=1$ and put $q=u/z$. The two maps
\[
 r_0=(x,z+q),\qquad r_1=(x,\zeta z+\zeta^{-1}q):D\to E
\]
extend to the smooth proper curves. We check their joint injectivity.
For finite $x\ne v$, the two image coordinates recover $z,q$, since
\[
 \det\begin{pmatrix}1&1\\ \zeta&\zeta^{-1}\end{pmatrix}=1.
\]
Together with $Y=a z^3-A$ they recover the point. At $x=v$, let
$\mathcal R=\{s:s^3=h(v)/a\}$; it consists of three distinct
nonzero elements. The three points above $R$ have $z=0$, $q\in
\mathcal R$, and their two $w$-coordinates are
$(q,\zeta^{-1}q)$. The other three have $q=0$, $z\in\mathcal R$,
and coordinates $(z,\zeta z)$. These six pairs are distinct.
At infinity the three residues of $z/x$ are $1,\zeta,\zeta^2$,
while the residues of $q/x$ are their inverses. Writing
$T=[1:1:0]\in E$, their image pairs are
\[
                         (O,T),\quad(T,T),\quad(T,O).
\]
Every finite point maps to the affine chart $Z=1$, so no other
collision with these pairs occurs. The two maps are jointly injective
on all geometric points.

All choices are defined after a finite extension of $K$. A hypothetical
collision of point sections persists over that extension and over a
further extension making a chosen lift to $D$ rational. The fixed-fibre
argument of A.3 supplies a matching lift whose Jacobian difference is
$2$-primary by A.1. Applying $r_0$ and $r_1$ kills this difference,
because $E(\Omega)[2^\infty]=0$. Joint injectivity forces the two
lifts, and then the original points, to coincide. This is the argument
of A.10, with the same Kummer input valid over finitely generated
fields. A collision on an open subcurve maps to one on its proper
compactification. These observations also prove the assertions about
finite extensions and forms.
\end{proof}

\begin{claim}{A.C2.2}{Proved}
For the ordinary genus-two curve $C_{a,b,c}$ over an algebraically
closed field of characteristic two, there exists a connected cyclic
\'etale cubic cover with a supersingular elliptic quotient by a lift
of the hyperelliptic involution exactly when (C2.1) holds.
If (C2.1) fails, every connected cyclic \'etale cubic cover has
ordinary Jacobian and admits only constant maps to abelian varieties
of $2$-rank zero.

For the $2$-rank-one genus-two normal form
$y^2+y=a x^3+b x+c/x$, with $ac\ne0$, every elliptic quotient
of a cyclic cubic \'etale cover by a lifted hyperelliptic
involution is ordinary as well.

In particular take $a=b=1$ and $c=\alpha\in\F_4$ with
$\alpha^2+\alpha+1=0$. Then $(a+b+c)^3+abc=1+\alpha\ne0$.
Thus the construction using a single cyclic cubic cover and
$2$-rank-zero abelian targets cannot settle this curve. This is an
obstruction to that specified construction; covers of other odd
degrees remain available.
\end{claim}
\begin{proof}
We check that the normal form used above exhausts cyclic cubic
\'etale covers. Such covers over an algebraically closed field are
classified, up to choice of generator, by nonzero
$\eta\in J_C[3]$. Riemann--Roch supplies an effective divisor
$D_0$ of degree two with $\eta=[D_0-2P_\infty]$. It is unique:
a degree-two divisor with two independent sections belongs to the
hyperelliptic pencil and would give $\eta=0$.
It cannot contain $P_\infty$. Otherwise
$\eta=[S-P_\infty]$ and $3\eta=0$ would give a function with
its only pole of order three at the Weierstrass point $P_\infty$,
where three is a gap. In the displayed model this also follows from
$L(3P_\infty)=\langle1,x\rangle$.

Choose a function with divisor $3D_0-6P_\infty$.
The space $L(6P_\infty)$ has basis $1,x,x^2,x^3,Y$.
Its $Y$ coefficient cannot vanish: otherwise the divisor is invariant
under the hyperelliptic involution. An invariant effective divisor of
degree two is either a hyperelliptic fibre or a sum of Weierstrass
points, whose class relative to $2P_\infty$ is killed by two.
Neither can represent the nonzero three-torsion class $\eta$.
After scaling, the function is therefore $A+Y$, with
\[
 A=Lx^3+dx^2+ex+k,\qquad L\ne0.
\]
Let $u=x^2+u_1x+u_0$ be the monic polynomial whose zero divisor on
$C$ is $D_0+\iota D_0$. Taking norms and comparing the leading
coefficient gives
\[
                     A^2+hA+f=L^2u^3.                       \tag{C2.5}
\]
The required cubic cover is $z^3=(A+Y)/L$, after a harmless
constant rescaling over the algebraically closed field. Thus no
three-torsion class of smaller reduced degree has been omitted.

The lift $\tau:z\mapsto u/z$ of the hyperelliptic involution has
elliptic quotient. Here is a ramification verification in
characteristic two. The cubic cover has genus four. Over each of
the three branch points of $C\to\mathbf P^1$, the composite
$S_3$-cover has three points with distinct order-two inertia groups.
Exactly one is fixed by the specified $\tau$, and its different
exponent is two, as in the original Artin--Schreier cover. There
are no other fixed points. Riemann--Hurwitz for the quotient by
$\tau$ therefore reads $6=2(2g(E)-2)+6$, so $g(E)=1$.
Its invariant functions $x,w=z+u/z$ give the integral plane cubic
\[
 L\bigl(W^3+(X^2+u_1XZ+u_0Z^2)W\bigr)=X^2Z+XZ^2.            \tag{C2.6}
\]
It has the genus-one function field just computed, so its arithmetic
genus one forces it to be smooth. The Cech calculation in A.C2.1
makes its Hasse scalar $Lu_1$. It is supersingular exactly when
$u_1=0$.

Assume $u_1=0$. The coefficients of $x^5,x,x^3$ in (C2.5) give,
respectively,
\[
                  L=a,\qquad k=b,\qquad e=d+a+b+c.
\]
The coefficients of $x^4,x^2,1$ then imply
\[
 d^2+d+a=a^2u_0,\qquad
 u_0^2+u_0+1=((b+c)/a)^2,\qquad u_0^3=(b/a)^2.               \tag{C2.7}
\]
Choose $v$ with $v^2=u_0$. Since squaring is injective,
(C2.7) gives $v^3=b/a$ and $v^2+v+1=(b+c)/a$.
Consequently $(v+1)^3=c/a$, which implies (C2.1).
Conversely A.C2.1 constructs a connected cubic cover with
supersingular quotient whenever (C2.1) holds.

For the stronger assertion let $D\to C$ be any connected cyclic
cubic cover, let $\delta$ generate its deck group, and use the lift
$\tau$ above. Its Prym variety $P$ has dimension two, and
\[
                       J_D\sim J_C\times E^2.              \tag{C2.8}
\]
We justify the isogeny also in characteristic two. Choose an auxiliary
prime different from two and three and work on rational Tate modules.
The invariant part for $\delta$ is $VJ_C$; its complement is $VP$.
The hyperelliptic involution acts as $-1$ on $VJ_C$. Thus the image
of $VE$ under pullback lies in $VP$ and equals $(VP)^\tau$.
The two subspaces $(VP)^\tau$ and
$\delta(VP)^\tau$ have zero intersection: an element fixed by the
two corresponding reflections is fixed by $\delta$, whereas $VP$
has no $\delta$-invariants. Each has dimension two, and together
they span $VP$. The two pullback maps $E\to P$ consequently give
an isogeny $E^2\to P$. Pullback and norm for $D\to C$, whose
composition is $[3]$, give (C2.8).

If (C2.1) fails, every elliptic quotient (C2.6) is ordinary. The
base curve is ordinary, so (C2.8) makes $J_D$ ordinary. An abelian
quotient of an ordinary abelian variety is ordinary, by Poincare
reducibility and additivity of $2$-rank under products and isogenies.
A positive-dimensional abelian subvariety of a $2$-rank-zero variety
has $2$-rank zero. Hence a homomorphism from $J_D$ to such a variety
has zero image. The Albanese property then makes every map from
$D$ to such a variety constant. Substituting $(1,1,\alpha)$ gives
the asserted explicit example.

For the $2$-rank-one normal form put $Y=xy$, so $h=x$ and
$f=a x^5+b x^3+c x$. The same reduced-divisor argument and
ramification calculation give (C2.5)--(C2.6), with $h=x$.
The right side of the quotient cubic is then $XZ^2$.
The $x^5$ coefficient of the norm on the left is $a$, whereas
the right coefficient is $L^2u_1$. Thus $u_1\ne0$. The
coefficient of $XWZ$ in the quotient cubic remains $Lu_1$,
so it is ordinary. The base curve has $2$-rank one: Cartier
fixes $dx/x$ and kills $dx$, its regular differential basis.
\end{proof}

\begin{claim}{A.C2.3}{Proved}
Let $K$ be finitely generated over $\F_2$ and let $s\in K^\times$.
For the smooth proper genus-two curve
\[
                 C_s:\quad y^2+y=sx^3+x+\frac{s^5}{x},       \tag{C2.9}
\]
which has $2$-rank one, the geometrically prime-to-two point-section
map is injective. The assertion persists after finite extensions,
for finite-extension forms, and for hyperbolic open subcurves.

Equivalently this covers the $2$-rank-one models
$y^2+y=a x^3+b x+c/x$ with $abc\ne0$ and $a^5=b^{16}c$.
There is a connected cyclic \'etale cover $D\to C_s$ of degree
five with two jointly point-injective maps to a genus-two curve
of $2$-rank zero.
\end{claim}
\begin{proof}
The two Artin--Schreier poles have orders three and one, so the
genus is two. With $Y=xy$ the equation is
\[
 Y^2+xY=f,\qquad f=sx^5+x^3+s^5x.
\]
The regular differentials $dx/x,dx$ have Cartier images $dx/x,0$;
hence the $2$-rank is one. The point $P_\infty$ above infinity
again has pole orders two and five for $x,Y$.

Over a finite extension choose $d$ with $d^2+d=s^8$, and put
\[
 u=x^2+s^2,\qquad
 A=x^5+s x^4+x^3+x^2+d x+s^5.
\]
The direct identity
\[
                        A^2+xA+f=u^5                        \tag{C2.10}
\]
gives $F\iota(F)=u^5$ for $F=A+Y$.
At $x=s$, the hyperelliptic map is unramified and the two values
of $F$ differ by $s\ne0$. Exactly one point $R$ has $F=0$.
Since $u$ has order two there and $F$ has pole order ten at
infinity, its divisor is $10R-10P_\infty$. Consequently
$z^5=F$ defines an \'etale cyclic cover $D$ of degree five.
It is connected: a fifth root of $F$ would have its only pole of
order two at $P_\infty$, hence belong to $\langle1,x\rangle$,
which cannot produce the nonzero $Y$ coefficient. The genus of
$D$ is six.

The involution $\tau:z\mapsto u/z$ lifts the hyperelliptic
involution. Put $w=z+u/z$ and $X=x+s$, so $u=X^2$. The
fifth-power identity in characteristic two gives
\[
 w^5+X^2w^3+X^4w=x.
\]
Set $t=X^2+Xw+w^2$. Its left side is $w t^2$, so $x=w t^2$.
Substituting this into the definition of $t$ gives
\[
 (t^4+t^2+1)w^2+s w+(s^2+t)=0.
\]
Writing $H=t^4+t^2+1$ and $V=Hw/s$ exhibits the quotient
function field as that of
\[
 B_s:\quad V^2+V=\frac{(t^4+t^2+1)(t+s^2)}{s^2}.             \tag{C2.11}
\]
Indeed $\Omega(D)^\tau=\Omega(x,w)$, since $z$ satisfies
$z^2+wz+u=0$, and the displayed rational transformations identify
$\Omega(x,w)$ with $\Omega(t,V)$. The right side of (C2.11) has
one pole of order five. Thus $B_s$ is smooth proper of genus two,
and its regular differentials are $dt,t\,dt$. Cartier sends them
to $0,dt$, respectively. It is nilpotent, so $B_s$ has $2$-rank
zero and $J_{B_s}(\overline K)[2^\infty]=0$.

Let $r_0:D\to B_s$ be the quotient map, and let $\delta$ generate
the cyclic deck group of $D/C_s$. Define $r_1=r_0\circ\delta$.
These two maps jointly separate all geometric points by a group
action argument which also handles the ramification points.
A fibre of a quotient by $\langle\tau\rangle$ is a geometric
orbit. If distinct $P,Q$ had both images equal, then
\[
             Q=\tau P=\delta^{-1}\tau\delta P.
\]
The dihedral identity $\tau\delta\tau=\delta^{-1}$ would force
$\delta^2P=P$. This is impossible because a nonidentity deck
transformation of the \'etale cyclic cover $D/C_s$ fixes no point.

Compose $r_0,r_1$ with an Abel embedding of $B_s$ into its
Jacobian, choosing a base point after a finite extension if needed.
The resulting two maps have targets with no geometric
$2$-primary torsion points and are still jointly injective.
The section-lifting and Kummer argument of A.C2.1, or A.10,
therefore proves the asserted point-section injectivity. Finally,
$x'=b x$ transforms $y^2+y=a x^3+b x+c/x$ into
$y^2+y=(a/b^3)x'^3+x'+bc/x'$. The relation
$a^5=b^{16}c$ makes $bc=(a/b^3)^5$, proving the equivalent
parameter statement. The assertions about open subcurves and
finite extensions follow as in A.C2.1.
\end{proof}

\subsection{A thirteen-fold cover of a second characteristic-two curve}

\begin{claim}{A.C2.4}{Proved}
Let $K$ be finitely generated over $\F_2$. The geometrically
prime-to-two point-section map is injective for the curve
\[
                         C_0:\quad y^2+y=x^3+\frac1x,         \tag{C2.12}
\]
and for its finite-extension forms and hyperbolic open subcurves.
There is a connected cyclic \'etale cover of degree thirteen with
two jointly point-injective quotient maps to a genus-six curve of
$2$-rank zero.
\end{claim}
\begin{proof}
The change $u=1/x$, $V=u^2y$ gives
\[
                    V^2+u^2V=u^5+u.                         \tag{C2.13}
\]
Write $P_\infty$ for the unique point above $u=\infty$ and
$P_0$ for that above $u=0$. In this model the Artin--Schreier
pole orders are one at infinity and three at zero. In particular,
$C_0$ has genus two and $2$-rank one. One can see the latter
from its differential basis $du/u^2,du/u$: Cartier kills the
first and fixes the second.

We first exhibit a thirteen-torsion point on a specified translate
of the curve in its Jacobian. Choose
\[
 \alpha^2+\alpha+1=0,\qquad
 \beta^2+\alpha^2\beta=1,
\]
so that $\alpha,\beta\in\F_{16}$, and set $R=(\alpha,\beta)$
on (C2.13). Define
\[
 \begin{split}
 B(u)&=u^4+\alpha u^3+u+1,\\
 A(u)&=\alpha(1+\beta)u^6+\alpha^2\beta u^5+\alpha u^4
       +\alpha(1+\beta)u^3\\
     &\hspace{12mm}+(1+\alpha\beta)u^2+\alpha u+\alpha^2.
 \end{split}
\]
Direct expansion in this field gives
\[
 A^2+u^2AB+(u^5+u)B^2=(u+\alpha)^{13}.                        \tag{C2.14}
\]
For $G=A+BV$, one has $G(R)=0$ and
$G(\iota R)=\alpha\ne0$. Its pole order at $P_\infty$ is
thirteen: $A$ has pole order twelve, while $BV$ has pole
order thirteen. Taking the norm in (C2.14), with the conjugate
factor a unit at $R$, therefore gives
\[
                    \operatorname{div}(G)=13R-13P_\infty.
\]
The class $L=\mathcal O(R-P_\infty)$ is nontrivial, since a
function with one simple pole would make the curve rational.
It has exact order thirteen. The cover $D$ defined by $z^{13}=G$
is \'etale and connected; otherwise $G$ would be a thirteenth
power, again giving a function with divisor $R-P_\infty$.
Its genus is fourteen.

Here is the complete Frobenius calculation which controls the
$2$-rank of this cover. On a smooth curve in characteristic two,
let $\mathcal B$ be the image of
$d:F_*\mathcal O_C\to F_*\Omega_C^1$, with the module structure
induced by absolute Frobenius. The sequence
\[
 0\longrightarrow\mathcal O_C\longrightarrow F_*\mathcal O_C
   \longrightarrow\mathcal B\longrightarrow0                \tag{C2.15}
\]
is exact. Locally in a parameter $t$, the image is the free
rank-one module of differentials $h(t)^2dt$, so $\mathcal B$
is a line bundle. A rational locally exact differential with
ordinary divisor $2D_1$ is a rational section of $\mathcal B$
with divisor $D_1$; this follows from the same square-coefficient
local description.

For (C2.13), the different formula gives
\[
 \operatorname{div}(du)=4P_0-2P_\infty,
 \qquad \operatorname{div}(u)=2P_0-2P_\infty.
\]
The locally exact differential $d(1/u)=du/u^2$ therefore has
divisor $2P_\infty$. Consequently
\[
                         \mathcal B\simeq\mathcal O(P_\infty).
\]
Tensor (C2.15) by $L$ and use the projection formula. The map
on first cohomology is the semilinear Frobenius map
\[
                   H^1(C,L)\longrightarrow H^1(C,L^2).
\]
Both spaces have dimension one, by Riemann--Roch and nontriviality
of these degree-zero line bundles. Its kernel contains, and equals
the image of, $H^0(C,\mathcal B\otimes L)=H^0(C,\mathcal O(R))$,
which has dimension one. Indeed the preceding $H^0(C,L^2)$
is zero. Thus this Frobenius map is zero.

Over an algebraic closure the eigenspace decomposition for the
cyclic cover is
\[
 H^1(D,\mathcal O_D)=\bigoplus_{i\in\Z/13}H^1(C,L^i).
\]
Each nonzero-index summand is one-dimensional. Frobenius sends
the summand indexed by $i$ to that indexed by $2i$. The powers
\[
                1,2,4,8,3,6,12,11,9,5,10,7
\]
are all the nonzero residues modulo thirteen. Thus every
nonzero-index summand lies on a single Frobenius cycle, containing
the zero map just found. The twelfth Frobenius power is zero on
the sum of these summands. The trivial summand is
$H^1(C,\mathcal O_C)$ and has stable Frobenius rank one.
The Artin--Schreier cohomology sequence identifies $2$-rank with
the dimension of Frobenius-fixed cohomology
\citep[Definition 1.1 and the following discussion]{A_Yang2016},
so $D$ has $2$-rank one as well.

The norm identity lifts the hyperelliptic involution to
\[
                  \tau:z\longmapsto (u+\alpha)/z.
\]
If $\delta$ generates the cyclic deck group, then
$\tau\delta\tau=\delta^{-1}$. The quotient $B=D/\langle\tau\rangle$
has genus six. Indeed, above each of the two branch points of
$C\to\mathbf P^1$, exactly one of the thirteen points is fixed
by $\tau$. Their different exponents are respectively four and
two, unchanged by the \'etale local base change
$D\to C$. Riemann--Hurwitz gives
$26=2(2g(B)-2)+6$, hence $g(B)=6$.

The same dihedral representation argument as in A.C2.2 gives
\[
                           J_D\sim J_C\times J_B^2.
\]
For clarity, the complement of the cyclic invariants in a rational
Tate module has dimension twenty-four. Its fixed subspace under
$\tau$, and its translate by $\delta$, each have dimension twelve.
Their intersection is zero, since an element in it is fixed by
$\delta^2$ and therefore by $\delta$, whose order is odd.
These are precisely the two pulled-back copies of $VJ_B$.
The two maps from $J_B$ thus give the claimed isogeny on the Prym.
Additivity of $2$-rank and the calculation $f(D)=f(C)=1$ imply
$f(B)=0$.

Two deck-conjugate quotient maps $D\to B$ jointly separate points:
a nontrivial collision would again give a fixed point of
$\delta^2$, as in A.C2.3. After composing with an Abel embedding
into $J_B$, their targets have no nonzero geometric $2$-primary
torsion points. The section-lifting argument of A.10 proves the
claimed injectivity. All equations, the deck transformation and
chosen base points are defined after finite extensions; this proves
the stated field and form variants.
\end{proof}

\subsection{The exact theta-orbit criterion and its remaining arithmetic condition}

\begin{claim}{A.C2.5}{Proved}
Let $C$ be a smooth proper genus-two curve over an algebraically closed
field $\Omega$ of characteristic two. Let $\mathcal B$ be its line bundle
of locally exact differentials, with the absolute-Frobenius module
structure, and put
\[
 \Theta=\{L\in\operatorname{Pic}^0(C):
                   H^0(C,\mathcal B\otimes L)\ne0\}.
\]
For an odd integer $n>1$, let $L$ have exact order $n$, and let
$D\to C$ be the associated connected cyclic \'etale cover. Write
$\mathcal O_n$ for the orbits of multiplication by two on
$(\mathbf Z/n)\setminus\{0\}$. Then
\[
 f(D)=f(C)+
 \sum_{\substack{O\in\mathcal O_n\\
                  L^i\notin\Theta\ \text{for all }i\in O}}|O|.
                                                               \tag{C2.16}
\]
The set $\Theta$ is symmetric under $L\mapsto L^{-1}$.
If every orbit meets $\{i:L^i\in\Theta\}$, a lift of the
hyperelliptic involution has a quotient $B$ of genus $(n-1)/2$
and $2$-rank zero. Two deck-conjugate quotient maps $D\to B$
jointly separate geometric points. In particular, for an odd prime
$\ell$, the sufficient condition is satisfied whenever
\[
 L\in\Theta[\ell]\setminus\{0\},\qquad
               \langle2,-1\rangle=(\mathbf Z/\ell)^\times.       \tag{C2.17}
\]
Over a field finitely generated over $\F_2$, a curve satisfying
this geometric condition has an injective geometrically
prime-to-two point-section map, as do its finite-extension forms
and hyperbolic open subcurves.
\end{claim}
\begin{proof}
The exact sequence (C2.15) is valid for every smooth curve in
characteristic two. Its local description also proves
$\mathcal B^{\otimes2}\simeq\Omega_C^1$: if $t$ is a local
parameter, a locally exact differential is $h(t)^2dt$, so the
transition functions for $\mathcal B$ square to those for
$\Omega_C^1$. Thus $\deg\mathcal B=1$. Riemann--Roch gives
\[
 h^0(\mathcal B\otimes M)
   =h^0(\Omega_C^1\otimes\mathcal B^{-1}\otimes M^{-1})
   =h^0(\mathcal B\otimes M^{-1})
\]
for every degree-zero $M$, proving the symmetry.

For $0\ne i\in\mathbf Z/n$, both $L^i$ and $L^{2i}$ are
nontrivial of degree zero. They have no global sections, and
their first cohomology spaces have dimension one. Tensoring
(C2.15) by $L^i$ therefore identifies the kernel of
\[
 F_i:H^1(C,L^i)\longrightarrow H^1(C,L^{2i})
\]
with $H^0(C,\mathcal B\otimes L^i)$. Hence $F_i$ is zero
precisely when $L^i\in\Theta$, and otherwise it is a bijective
semilinear map between one-dimensional spaces. The cyclic-cover
decomposition
\[
 H^1(D,\mathcal O_D)=
             \bigoplus_{i\in\mathbf Z/n}H^1(C,L^i)
\]
is respected by Frobenius, with indices multiplied by two.
On an orbit meeting $\Theta$, the $|O|$th Frobenius power is
zero. On an orbit disjoint from $\Theta$, Frobenius is bijective
and its stable rank is $|O|$. The index-zero summand has stable
rank $f(C)$. The Artin--Schreier description of $2$-rank now
proves (C2.16). This proof treats all nonzero characters,
including characters of proper conductor when $n$ is composite.

The hyperelliptic involution $\iota$ acts by $-1$ on
$\operatorname{Pic}^0(C)$, so it lifts to an automorphism
$\tau$ of the cyclic cover which conjugates a deck generator
$\delta$ to $\delta^{-1}$. Such a lift is an involution:
$\tau^2\in\langle\delta\rangle$ commutes with $\tau$, and
inversion fixes no nonidentity element of an odd cyclic group.
At a branch point of $C\to\mathbf P^1$, the action of $\tau$
on the $n$ points above it is an affine inversion of the cyclic
torsor, so exactly one of those points is fixed. The completed
local map $D\to C$ at that point is an isomorphism, and the
different exponent for $D\to D/\langle\tau\rangle$ equals
the exponent for $C\to\mathbf P^1$. The latter different has
total degree six, since $g(C)=2$. Since $g(D)=n+1$,
Riemann--Hurwitz gives
\[
       2n=2(2g(B)-2)+6,\qquad g(B)=(n-1)/2.
\]

The dihedral argument of A.C2.2 applies to every odd $n$:
\[
                         J_D\sim J_C\times J_B^2.
                                                               \tag{C2.18}
\]
Here is the precise linear-algebra verification. In an auxiliary
rational Tate module at a prime different from two and from the
divisors of $n$, let $W$ be the complement of the cyclic
invariants. Its nontrivial cyclic characters occur in pairs
$\lambda,\lambda^{-1}$, and $\tau$ exchanges the members of
each pair. Consequently $W^\tau$ has half the dimension of $W$.
The subspaces $W^\tau$ and $\delta W^\tau$ intersect trivially:
an element in the intersection is fixed by both corresponding
reflections, hence by $\delta^2$, hence by $\delta$, whereas
$W$ has no cyclic invariants. Both subspaces are pulled-back
Tate modules of the quotient Jacobian. They therefore span $W$
and give (C2.18). The cyclic-invariant part is $VJ_C$, on which
$\tau$ acts by $-1$.

If every orbit meets $\Theta$, formula (C2.16) and additivity
of $2$-rank under isogeny give $f(B)=0$. Two quotient maps
associated with the reflections $\tau$ and
$\delta^{-1}\tau\delta$ jointly separate points: a nontrivial
collision would be a fixed point of $\delta^2$, impossible
because the odd cyclic deck action is free. Composing with
Abel embeddings into $J_B$ and applying A.10 proves the
point-section assertion. The covers, quotient maps and chosen
base points descend after a finite extension, as required there.

Finally suppose $n=\ell$ is prime and (C2.17) holds. If
$L\in\Theta$, symmetry also gives $L^{-1}\in\Theta$.
Every nonzero residue modulo $\ell$ then belongs to the
doubling orbit of $1$ or of $-1$, proving the required
orbit condition. In particular, the condition is weaker than
requiring two itself to be a primitive root: for $\ell=7$
there are two doubling orbits, exchanged by sign, and both
are met simultaneously.
\end{proof}

\begin{claim}{A.C2.6}{Unresolved}
For an arbitrary genus-two curve over $\overline{\F}_2$, the
existence of an odd-order torsion point satisfying the
all-orbits condition of A.C2.5 is not established here.
Prime-to-two torsion density on $\Theta$ does not supply the
missing order and orbit restrictions.
\end{claim}
\begin{proof}[Exact remaining obligation]
The theta curve is a translate of the Abel image of $C$ and
generates its Jacobian. Yang's Proposition 2.3 therefore gives
density of odd-order torsion on it over $\overline{\F}_2$
\citep{A_Yang2016}. That statement does not prescribe a prime
order $\ell$, the subgroup $\langle2,-1\rangle$ modulo $\ell$,
or which other multiples of a chosen point remain on $\Theta$.
Formula (C2.16) identifies the exact missing condition: every
nonzero doubling orbit must contain a theta point. For a
composite order, the orbits of proper-order characters must
also be treated. A single point of exact order $n$ only
guarantees a zero map on its own orbit and, by symmetry,
on the opposite orbit.

This distinction is substantive. In the $2$-rank-one genus-two
family of A.C2.2 every cyclic cubic \'etale cover has an ordinary
elliptic involution quotient. Formula (C2.16) therefore implies
$\Theta\cap J_C[3]=\{0\}$ in this family. Thus even the
smallest prescribed odd prime need not occur on the nonzero
theta locus. Any cyclic cover whose degree is divisible by
three retains an ordinary contribution from these order-three
characters, regardless of theta intersections among its
primitive characters. The ordinary family off the hypersurface
of A.C2.1 has the same nonzero-three-torsion obstruction.

Nor does a single theta intersection give a quantitative
small-$2$-rank tower. Formula (C2.16) gives
\[
 \frac{f(D)}n=\frac{f(C)}n+
   \frac1n\sum_{O\text{ disjoint from }\Theta}|O|.
\]
To make this ratio small, theta intersections must eliminate
almost all characters. Density alone supplies no such
proportion. No assertion that arbitrary curves admit these
covers, or that normalized $2$-rank tends to zero in suitable
prime-to-two towers, is used in the positive conclusions.
\end{proof}

\subsection{A weaker criterion from the deck action on two-primary torsion}

\begin{claim}{A.C2.7}{Proved}
Let $C$ be a smooth proper geometrically connected curve of genus at
least two over a field $K$ finitely generated over $\F_2$. Suppose that,
over $\overline K$, it admits a connected nontrivial finite \'etale
Galois cover $D\to C$ of odd degree, with group $G$, for which
\[
             G\longrightarrow
             \operatorname{Aut}_{\F_2}(J_D[2](\overline K))
                                                               \tag{C2.19}
\]
is not faithful. Then the geometrically prime-to-two point-section
map of $C$ is injective. The conclusion also holds for its
finite-extension forms and hyperbolic open subcurves.

More precisely, a pair of distinct points giving equal section classes
would force (C2.19) to be faithful for every such cover. If
$m_2(G)$ denotes the least dimension of a faithful representation of
$G$ over $\F_2$, the same hypothetical pair would force
\[
                            f(D)-f(C)\ge m_2(G).               \tag{C2.20}
\]
In particular, a single such cover with $f(D)=f(C)$ suffices for
point-section injectivity.
\end{claim}
\begin{proof}
Work first over an algebraic closure. The curve $D$ is not
hyperelliptic. Otherwise its hyperelliptic involution is unique,
so every deck transformation descends to its genus-zero quotient.
Choose an element $\sigma\in G$ of odd prime order. Its image on
$\mathbf P^1$ has the same order, since the kernel of descent is
the order-two hyperelliptic group. A finite odd-order automorphism
of $\mathbf P^1$ in characteristic two has a fixed geometric
point. Its fibre on $D$ has at most two geometric points, so
$\sigma$ fixes a point in that fibre. This contradicts the free
deck action of an \'etale Galois cover.

Let $P',Q'\in D(\overline K)$ lie over distinct points of $C$.
The class $T=[Q'-P']$ has trivial stabilizer in $G$. Indeed,
$\delta T=T$ for $1\ne\delta\in G$ would give
\[
                     \delta Q'+P'\sim Q'+\delta P'.
                                                               \tag{C2.21}
\]
The two effective degree-two divisors in (C2.21) are distinct:
their terms lying over each of the two base points would
otherwise force both $\delta P'=P'$ and $\delta Q'=Q'$.
A linear equivalence between distinct effective degree-two
divisors on a curve of genus at least two gives a degree-two
map to $\mathbf P^1$, after removing any common base point.
A degree-one map is excluded by the genus. Thus (C2.21) would
make $D$ hyperelliptic, a contradiction.

The \'etale two-adic Tate module of $J_D$ is free of rank $f(D)$,
and its geometric two-primary torsion is the quotient of its
rational Tate module by this lattice. An odd-order automorphism
acting trivially on $J_D[2](\overline K)$ acts trivially on the
whole two-adic module. This can be checked without any
torsion-free congruence-group citation: if a matrix $A$ of odd
order $m$ satisfies $A\equiv1\pmod{2^r}$, then
\[
 A^m\equiv1+m(A-1)\pmod{2^{r+1}}
\]
implies $A\equiv1\pmod{2^{r+1}}$. Starting at $r=1$ and
iterating proves $A=1$. Thus any element in the kernel of
(C2.19) fixes every two-primary torsion point of $J_D$.

Now suppose that distinct $P,Q\in C(K)$ give equal section
classes. Descend the cover and its deck transformations after
a finite extension of $K$, and extend again to make a chosen
lift $P'$ rational. By A.3 there is a lift $Q'$ with
$T=[Q'-P']$ two-primary. A nonidentity element in the kernel
of (C2.19) would fix $T$, contrary to its trivial stabilizer.
This proves the criterion and the necessity of faithfulness.
The same argument after restriction proves the stated form and
open-subcurve variants.

For the quantitative assertion, the averaging idempotent
$|G|^{-1}\sum_{g\in G}g$ acts integrally on the two-adic Tate
module, since $|G|$ is odd. Its invariant summand is the
pullback of the two-adic Tate module of $J_C$: pullback
followed by norm is multiplication by $|G|$, and norm followed
by pullback is the sum of deck transformations. This summand
has rank $f(C)$. The complementary summand has rank
$f(D)-f(C)$, and its reduction modulo two must be a faithful
$G$-representation because the invariant summand is trivial.
This proves (C2.20).
\end{proof}

\begin{claim}{A.C2.8}{Proved}
In the genus-two cyclic situation of A.C2.5, let $S$ be the
union of the nonzero doubling orbits that miss $\Theta$.
The kernel of the deck action on geometric two-primary torsion
has order
\[
                    d=\gcd\bigl(n,\{i:i\in S\}\bigr),        \tag{C2.22}
\]
where $d=n$ if $S$ is empty. Consequently $d>1$ already implies
point-section injectivity by A.C2.7.

In particular, if $n=\ell^a$ is an odd prime power, it suffices
that every doubling orbit on $(\mathbf Z/n)^\times$ meets
$\{i:L^i\in\Theta\}$. No condition is needed on characters
of smaller order. The simple sufficient condition is
\[
 L\in\Theta\text{ of exact order }\ell^a,\qquad
                 \langle2,-1\rangle=(\mathbf Z/\ell^a)^\times.
                                                               \tag{C2.23}
\]
For a cyclic group of odd prime-power order $\ell^a$ one has
$m_2(G)=\operatorname{ord}_{\ell^a}(2)$.
\end{claim}
\begin{proof}
The Artin--Schreier description of \'etale $H^1$ identifies its
scalar extension with the stable Frobenius part of
$H^1(D,\mathcal O_D)$. The calculation of A.C2.5 therefore
shows that its nontrivial cyclic characters are exactly those
indexed by $S$, with multiplicity one. Dualizing to geometric
two-torsion may invert all characters; it does not change
their common kernel. A power $\delta^k$ acts trivially exactly
when $n\mid ik$ for every $i\in S$. The number of such powers
is the greatest common divisor in (C2.22).

If $n=\ell^a$, this kernel is nontrivial exactly when no
primitive character survives, since a primitive index is
coprime to $n$, whereas all nonprimitive indices are divisible
by $\ell$. Symmetry of $\Theta$ gives (C2.23) as in A.C2.5.
For example at $n=9$, a theta point of exact order nine kills
the single orbit of all six primitive characters. The
order-three block may remain ordinary, but the deck action
then has a kernel of order three, which suffices here. This
does not assert existence of such a nine-torsion theta point
on every curve.

Finally, an irreducible $\F_2$-representation of a cyclic
group corresponds to a doubling orbit of its geometric
characters. A faithful character of order $\ell^a$ has orbit
length $\operatorname{ord}_{\ell^a}(2)$. Any faithful
representation of this cyclic prime-power group must contain
such a character: otherwise all character orders divide
$\ell^{a-1}$ and a subgroup of order $\ell$ acts trivially.
One faithful orbit already gives a faithful representation,
proving the formula for $m_2(G)$.
\end{proof}

The full-quotient obstruction in A.C2.6 and the weaker deck
criterion are compatible. Surviving order-three characters
prevent the entire Prym from having $2$-rank zero, but at
degree $3^a$, $a>1$, they need not make the deck action
faithful. Torsion density still does not guarantee the
prime-power order or simultaneous primitive-orbit
intersections required by (C2.23). A degree-sized lower
bound in place of (C2.20), or a construction forcing a
sufficiently small normalized $2$-rank in general towers,
has not been established here.


\begin{claim}{A.26}{Proved; a nonabelian symmetric-group quotient class}
Let $K$ be finitely generated over $\F_p$, with
$p\nmid6$. Let $C/K$ be smooth proper geometrically
connected of genus at least two. Suppose that
$C_{\overline K}$ has a faithful $S_3$-action with
genus-zero quotient, and that the quotient by a subgroup
generated by a transposition has genus at most two.
Then A1 holds for $C$ and every hyperbolic open subcurve.
The three transposition quotients are isomorphic, so the
genus condition is independent of the selected subgroup.
\end{claim}
\begin{proof}
We first work over a finite field after a finite constant
extension defining the action. Put $N=A_3$ and $B=C/N$.
It is a double cover of the genus-zero final quotient.
If $g(B)=0$, then $C$ is a cyclic cubic cover of a
genus-zero curve and A.23 applies. If $g(B)\ge2$, it is
hyperelliptic and A.19 shows that any collision on $C$
has equal $B$-images. If $g(B)=1$, the cyclic cubic map
$C\to B$ is ramified because $g(C)\ge2$, so A.21 gives
the same equality.

Let $H_s=C/\langle s\rangle$ for each transposition $s$.
If their common genus is zero, $C$ is hyperelliptic and
A.19 settles the result. If it is one, either A.25
settles $C$ directly or A.21 gives equal $H_s$-images.
If it is two, A.19 gives equal $H_s$-images. Thus in the
remaining cases a collision has equal images on $B$ and
all three curves $H_s$.

These maps jointly distinguish every geometric point.
The two points have the same final base value, so write
$Q=gP$ for $g\in S_3$, and let $I$ be the stabilizer of
$P$. Tameness makes $I$ cyclic; it is trivial, a
transposition subgroup, or $N$. If $I=1$, equal images
in $B$ and one $H_s$ force
$g\in N\cap\langle s\rangle=1$. If
$I=\langle s\rangle$, its own quotient has fibre
$\langle s\rangle P=\{P\}$ above the image of $P$, so
equality in $H_s$ gives $Q=P$. If $I=N$, the same
argument applies to the fibre of $C\to B$. Hence no
nontrivial collision exists. Spreading the tame action
and specializing by A.8 gives the assertion over $K$;
functoriality gives the open assertion.

For a concrete branch pattern in this class, let $a$ be
the number of transposition branch points and $b$ the
number of three-cycle branch points. Tame
Riemann--Hurwitz in the regular and degree-three
permutation actions gives
\[
 g(C)=-5+\frac{3a}{2}+2b,\qquad
 g(B)=\frac{a-2}{2},\qquad
 g(H_s)=-2+\frac a2+b.
\]
The pattern $(a,b)=(6,1)$ therefore has genus six, with
$g(B)=g(H_s)=2$. It exists over a finite extension of
every $\F_p$ with $p\nmid6$: the branch-cycle tuple
\[
 (12),(12),(12),(12),(12),(23),(132)
\]
has product one and generates $S_3$. The tame
prime-to-$p$ presentation in SGA~1, XIII, Corollary 2.12
realizes it on a seven-punctured projective line, after
choosing distinct branch points over $\overline{\F}_p$.
The resulting finite cover and its smooth compactification
descend to a finite constant field. No claim about the
full automorphism group of these examples is needed.
\end{proof}

\begin{claim}{A.27}{Proved; abelian covers of genus one with exponent dividing 120}
Let $K$ be finitely generated over $\F_p$, with $p\ne2$.
Let $C/K$ be smooth proper geometrically connected of
genus at least two. Suppose that $C_{\overline K}$ has a
faithful action of a finite abelian group $G$ of order
prime to $p$, with genus-one quotient, and that the
exponent of $G$ divides $120$. Then A1 holds for $C$ and
every hyperbolic open subcurve. In particular this treats
every cyclic cover of a genus-one curve of degree
$3,4,5,$ or $8$ prime to $p$, with arbitrary allowed
partial ramification in degrees four and eight. For $p=3$ or $5$
the prime-to-$p$ hypothesis excludes the corresponding
primary factors of $G$.
\end{claim}
The independent branch-cancellation proof below is retained as an
earlier route. Its conclusion is subsumed by A.29, which removes
the exponent bound by a different divisor construction.
\begin{proof}
We work first over finite fields, making finite constant
extensions to define all data. There are six steps.

\emph{Two fully ramified branch points.}
Let $B\to E$ be a cyclic cover of degree
$n\in\{3,4,5\}$, with $p\nmid n$, branched at exactly
two points $A,B_0$, with full index $n$ at both. The
local Kummer exponents are opposite units modulo $n$.
Choose an elliptic inversion $\iota$ exchanging $A$ and
$B_0$; its origin is a half of their sum. Write the
Kummer equation and its divisor as
\[
 z^n=c,\qquad
 \operatorname{div}(c)=nD+u(A-B_0),\qquad \deg D=0,
 \quad\gcd(u,n)=1.
\]
Since $\iota$ acts as minus one on
$\operatorname{Pic}^0(E)$, the divisor $D+\iota^*D$ is
principal. After adjusting a constant choose $h$ with
\[
 \operatorname{div}(h)=D+\iota^*D,
 \qquad c\,\iota^*c=h^n.
\]
Thus $\iota$ lifts to $\tau(z)=h/z$. If $n$ is odd, the
constant $\iota^*h/h$ has both order dividing $n$ and
order dividing two, so it is one and $\tau^2=1$.
For $n=4$ this constant is a priori $1$ or $-1$. At any
of the four fixed points $T$ of $\iota$, the branch pair
is absent, so $v_T(c)$ is divisible by four and
\[
 v_T(h)=2v_T(c)/4
\]
is even. A local parameter for the tame involution has
$t\mapsto-t$; an anti-invariant rational function has
odd leading valuation. Therefore $\iota^*h=-h$ is
impossible, and again $\tau^2=1$. If $\delta$ generates
the cyclic deck group, then
$\tau\delta\tau=\delta^{-1}$.

Riemann--Hurwitz gives $g(B)=n$. For odd $n$, each of
the four unramified cyclic fibres above a fixed point of
$\iota$ has exactly one $\tau$-fixed point: on a cyclic
torsor of odd order an inversion has one fixed point.
Hence the quotient $H=B/\langle\tau\rangle$ has genus
$(n-1)/2$. For $n=4$ such a fibre has zero or two fixed
points. The total fixed count $r_\tau$ is between zero
and eight, and tame Riemann--Hurwitz reads
\[
 6=2(2g(H)-2)+r_\tau.
\]
Thus $r_\tau$ is two or six and $g(H)$ is two or one.
In all cases the reflection quotient has genus one or
two. When its genus is two, A.19 gives equal $H$-images
of a collision on $B$. When its genus is one, the
degree-two map $B\to H$ is ramified, so A.21 gives the
same conclusion. The same applies to the quotient map
composed with $\delta$.

A.21 applied to $B\to E$ also gives equal $E$-images.
A distinct pair with equal $E$-images lies away from
$A,B_0$, where the cyclic action is free. Equal images
under the two reflection quotients would give
\[
 Q=\tau P=\delta^{-1}\tau\delta P
              =\delta^{-2}\tau P,
\]
so $\delta^2 Q=Q$. This is impossible on that locus,
since $\delta^2$ is nontrivial for $n=3,4,5$. The two
full branch fibres themselves have one geometric point
each. This proves A1 for every such two-branch cover.

\emph{Arbitrary cyclic degrees three and five.}
Let $C\to E$ have prime degree $\ell=3$ or $5$. Its
nonempty branch set has cardinality $r\ge2$: the sum of
its nonzero Kummer residues is zero modulo $\ell$.
The case $r=2$ is the preceding step. For $r\ge3$,
choose two branch points $a_1,a_2$ and a cyclic
degree-$\ell$ cover $B\to E$ branched exactly there,
with exponents $u,-u$, where $u$ matches that of $C$ at
$a_1$. Such a cover exists by divisibility of
$\operatorname{Pic}^0(E)$ by $\ell$; its character can
be represented by a rational function with divisor
$u(a_1-a_2)+\ell D$ for a degree-zero divisor $D$.
It has the two-branch injectivity just proved.

The two prime-degree fields are distinct since their
branch sets differ. Their normalized fibre product $V$
is a connected cyclic \'etale $\ell$-cover of $C$.
The ratio Kummer character defines a third quotient
$H_2\to E$, whose branch set loses $a_1$ and possibly
$a_2$, with every other branch unchanged. It remains a
connected degree-$\ell$ cover and has fewer branch
points, so induction gives its injectivity. A collision
on $C$ has equal $E$-images by A.21. If the points are
distinct, this common image avoids the entire branch
set, since every branch fibre of $C$ is total. Lift the
collision to $V$. Its images on $B$ and $H_2$ agree.
Their fibre product is \'etale over $E$ at the common
base value and is therefore already normal there. It
agrees with $V$ at both points, so those points coincide.
This proves the branch-count induction. A putative
one-point branch set on $H_2$ is excluded by the same
sum-of-residues relation; the inductive endpoint is
again a two-branch cover.

\emph{Arbitrary cyclic degree four.}
Let $C\to E$ have degree four, and let $C_0\to E$ be
its quadratic quotient. The branch points of $C_0$ are
exactly the full-index-four branch points of $C$.
If there are none, $C_0$ is a connected unramified
quadratic cover of $E$, hence has genus one. The double
cover $C\to C_0$ is covered by A.25.

Otherwise the number of full-index points is positive
and even. Choose two of them, $a_1,a_2$, and a cyclic
degree-four cover $B\to E$ with exactly these branch
points, matching the exponent of $C$ at $a_1$ and using
the opposite exponent at $a_2$. We can and do choose
its quadratic quotient different from $C_0$. To justify
this choice even when $C_0$ has only these two branch
points, twist the Kummer character of $B$ by an
unramified cyclic character of order dividing four.
The reduction of these characters modulo two exhausts
the four unramified quadratic characters. Explicitly,
functions with divisor $4(T-O)$, for $[T-O]\in E[4]$,
give the characters, and their quadratic reductions
give $2[T-O]\in E[2]$. Multiplication by two from
$E[4]$ to $E[2]$ is surjective. At most one of these
four quadratic twists equals $C_0$. The twist changes
none of the branch indices or residues of $B$.

Distinct quadratic subfields make the cyclic
degree-four fields of $C$ and $B$ linearly disjoint.
Their normalization $V$ is consequently a connected
cyclic \'etale degree-four cover of $C$. The ratio
character $h=c/b$ has exact order four: the change
$(c,b)\mapsto(b,c/b)$ is an invertible change of basis
in the character group $(\Z/4)^2$. Let $H_2$ be its
degree-four quotient over $E$. The branch at $a_1$
disappears. At $a_2$ the two odd exponents add to an
even exponent, so it is no longer a full-index point;
it may have index two or be unramified. Every other
full-index point is unchanged. Thus their number drops
by two, allowing induction on this number, with the
unramified-quadratic case already supplying its base.

More precisely, a collision on $C$ has equal $E$-images
by A.21; a distinct pair avoids $a_1,a_2$, whose fibres
are total. Lift to $V$. Its $B$-images agree by the
two-branch result. If $g(H_2)\ge2$, induction gives
equal $H_2$-images. If $g(H_2)=1$, the map $V\to H_2$
is fully ramified of index four over every point above
$a_1$, since $H_2$ is unramified there and $B$ has
index four. A.21 gives the same equality. Genus zero
cannot occur because $H_2$ maps separably to $E$.
Away from $a_1,a_2$, $B\to E$ is \'etale, so
$B\times_EH_2$ is already smooth and normal and equals
$V$. Its two projections therefore separate the lifted
points. This proves the degree-four assertion, including
all partial ramification.

\emph{Cyclic degree eight.}
First let $B_8\to E$ have exactly two full-index branch
points. The Kummer/Picard argument in the first step
again lifts inversion exchanging them to an involution
$\tau$: in the notation there, $v_T(h)=2v_T(c)/8$ is
even at each inversion-fixed point, ruling out the
anti-invariant root. Now $g(B_8)=8$, so the fixed count
of any reflection is two or six. The two reflection
classes have fixed counts summing to eight. Indeed on
each of the four fixed elliptic fibres the reflections
act on $\Z/8$ as $j\mapsto a-j$ and $j\mapsto a+1-j$;
exactly one of these has two fixed elements. Choose
$\tau$ with fixed count six. Its quotient $H_8$ has
genus three.

The central involution $\delta^4$ commutes with $\tau$
and descends to a nontrivial involution on $H_8$. Its
quotient is the reflection quotient $H_4$ of the cyclic
degree-four cover $B_4=B_8/\langle\delta^4\rangle$.
At an inversion-fixed elliptic point, the affine
reflection has fixed elements modulo eight exactly when
$a$ is even, and the same is true modulo four. Each
fixed count is two in that case. Thus the induced
reflection on $B_4$ also has total fixed count six,
and $H_4$ has genus one. Consequently
$H_8\to H_4$ is a double cover of a genus-one curve,
and A.25 proves point-section injectivity for $H_8$.
Use this target for two deck-conjugate reflection maps
from $B_8$. Together with the equal elliptic images
from A.21, the same $\delta^2$ argument as in the first
step proves injectivity for $B_8$.

For an arbitrary cyclic degree-eight $C\to E$, induct
on its number of full-index-eight branch points. If
there are none, its quadratic quotient $C_0\to E$ is
unramified and has genus one; the degree-four map
$C\to C_0$ is covered by the preceding step. Otherwise
choose two full-index branch points and a two-branch
cyclic degree-eight auxiliary cover, matching the first
residue and using the opposite residue at the second.
Twisting by unramified degree-eight characters makes
its quadratic quotient different from $C_0$, since
$[4]:E[8]\to E[2]$ is surjective. Distinct quadratic
subfields imply disjointness of the cyclic degree-eight
fields. The normalized fibre product is a connected
cyclic \'etale degree-eight cover of $C$, and its ratio
quotient $H_2\to E$ still has degree eight. The number
of full-index points drops by two, although the second
selected point may retain index two or four.

The rest of the degree-four proof applies with eight
in place of four. A.21 gives equal elliptic images; a
distinct pair avoids the two selected full-index
fibres. The auxiliary images agree by the two-branch
result. The ratio images agree by induction when its
genus is at least two, or by A.21 when its genus is one,
because the map from the common cover is fully ramified
of index eight above the canceled first point. Away
from the two selected points the auxiliary cover is
\'etale, so the final fibre product is already normal
and its projections separate the lifted pair.

\emph{Cyclic degrees dividing 120.}
Degrees two, three, four, five and eight are now settled by
A.25 and the preceding steps. If $n\mid120$ has at
least two distinct prime factors, split it as $n=ab$
with $a,b>1$ coprime and use its two corresponding
quotient curves over $E$. Each has genus at least one.
A quotient of genus at least two is settled by induction
on the number of distinct prime factors; a genus-one quotient
has equal images by A.21, since the Galois map from $C$
is ramified. The two quotient maps jointly separate
geometric points because their local tame indices are
coprime, by the normalization argument in A.23.

\emph{The abelian group and field extensions.}
For a finite abelian $G$ as in the statement, each
nontrivial character quotient is a cyclic cover of $E$
of degree dividing $120$. If its genus is at least two,
the preceding step gives equal images of a collision;
if it is one, A.21 gives the equality. There is no
genus-zero quotient above $E$. With $Q=gP$ and
$I=\operatorname{Stab}(P)$, these equalities give
$g\in(\ker\chi)I$ for every character $\chi$.
Characters of the abelian group $G/I$ separate points,
so $g\in I$ and $P=Q$, including at every ramified
fibre. Finally spread the tame action and its genus-one
quotient and apply A.8 to obtain the result over $K$.
Functoriality to the proper curve gives the open case.
\end{proof}

\subsection{A finite obstruction to the configuration-space exactness route}

\begin{claim}{A.28}{Refuted}
The proposed universal left exactness of the geometrically prime-to-$p$
configuration-space homotopy sequence is false. More precisely, let
$C/\mathbb F_5$ be the smooth proper genus-two curve with affine equation
\[
 y^2=x^5-x+1,
\]
let $P=\infty$, and put $U=(C\times C)\setminus\operatorname{diag}(C)$.
Then the natural map
\[
 \pi_1((C\setminus\{P\})_{\overline{\mathbb F}_5})^{(5')}
 \longrightarrow \pi_1(U_{\overline{\mathbb F}_5})^{(5')}
\]
is not injective. This refutes the sufficient topological condition
considered in \citet[Remark 3.13]{A_SaidiTamagawa2009}; it does not
refute point-section injectivity, which holds for this $C$ by A.19.
\end{claim}
\begin{proof}
We first work with the topological two-point configuration space of an
oriented closed genus-two surface. Write $B$ for its pure braid group
and $F\triangleleft B$ for the kernel of the projection that forgets
the second point. The group $F$ is free on four generators
$a,b,d,e$. Throughout this proof set
\[
 [u,v]=u^{-1}v^{-1}uv,
 \qquad r=([a^{-1},b][d^{-1},e])^{-1}.
\]
Thus $r$ is the peripheral element. We use the corrected presentation
of \citet[Proposition 4, preprint p.~5]{A_BardakovBellingeri2009}.
Its three point-pushing conjugations $A_1,B_2,A_2$ restrict to the
following automorphisms $\alpha,\beta,\gamma$ of $F$:
\[
\begin{array}{c|cccc}
 &a&b&d&e\\ \hline
\alpha&r^{-1}ar&r^{-1}b[a,r]&[r,a]d[r,a]^{-1}&[r,a]e[r,a]^{-1}\\
\beta&a&b&dr&r^{-1}er\\
\gamma&a&b&r^{-1}dr&r^{-1}e[d,r].
\end{array}
\]
The source defines its commutator and its left-conjugation notation
explicitly. It also corrects the earlier exceptional-relation
misprints in Bellingeri's presentation. The certificate below checks
these formulas by free reduction, including explicit inverse
automorphisms and the base surface relation.

Represent a homomorphism $F\to S_4$ by the tuple of its values on
$(a,b,d,e)$, and compose permutations right to left. Let $\Phi$ be the
operation on such tuples obtained by successive substitution using
$\alpha$, then $\beta$, then $\gamma$. This is the action induced by
the element $w=A_1B_2A_2$ of $B$. The following small table is an exact
certificate. In its last column $h_i$ satisfies
$h_i\Phi(T_i)h_i^{-1}=T_{i+1}$, with subscripts modulo five:
\[
\begin{array}{c|cccc|c}
i&a&b&d&e&h_i\\ \hline
0&(34)&(234)&(134)&(143)&(24)\\
1&(34)&(234)&(123)&(234)&(13)\\
2&(34)&(234)&(123)&(143)&(234)\\
3&(34)&(123)&(234)&(123)&(243)\\
4&(34)&(234)&(134)&(124)&(13).
\end{array}
\]
Each equality is a multiplication of permutations of four letters.
The first tuple generates $S_4$: its first two entries generate the
stabilizer of $1$, and its third entry moves $1$. Its peripheral image
is $(234)$.

The five tuples belong to different simultaneous-conjugacy classes.
For rows $0,1,2,4$ the first pair is the same and has trivial common
centralizer in $S_4$, while the last pair differs. The first pair of
row $3$ generates $S_4$, whereas each of the other four first pairs
generates a subgroup of order six, so row $3$ cannot be conjugate to
one of them. Consequently the action of $w$ on
$\operatorname{Hom}(F,S_4)/S_4$ has a cycle of length exactly five.

Suppose, for a contradiction, that
$\widehat F^{(5')}\to\widehat B^{(5')}$ were injective. An injective
continuous map from a compact group to a Hausdorff group identifies
its source with a closed subgroup. Therefore the topology on
$\widehat F^{(5')}$ is the induced topology. Since $|S_4|=24$, the
homomorphism $\rho:F\to S_4$ specified by $T_0$ extends continuously
to $\widehat F^{(5')}$. There is an open normal subgroup
$N\triangleleft\widehat B^{(5')}$ whose inverse image in
$\widehat F^{(5')}$ is contained in $\ker\rho$.

Let $m$ be the order of the image of $w$ in the finite group
$\widehat B^{(5')}/N$. It is prime to five. For every $x\in F$, the
commutator of $w^m$ with $x$ lies in $N\cap\widehat F^{(5')}$, so
\[
 \rho(w^m xw^{-m})=\rho(x).
\]
The $w$-orbit of the simultaneous-conjugacy class of $\rho$ therefore
has length dividing $m$. Its length is five by the displayed table,
a contradiction. This proves the topological noninjectivity.

To transfer the calculation, take the smooth proper model of
$y^2=x^5-x+1$ over $\mathbb Z_5$, together with its infinity section.
Its special fibre is smooth because the derivative of the right-hand
side is $-1$ in characteristic five. At infinity the coordinates
$u=1/x$, $v=y/x^3$ give $v^2=u-u^5+u^6$, whose derivative in $u$
is $-1$ at $(u,v)=(0,0)$. The product of this
model with itself is smooth and proper, and its diagonal is a smooth
relative Cartier divisor. Thus the configuration space and the
punctured fibre are complements of relative normal-crossing divisors
in smooth proper schemes. Prime-to-five specialization is an
isomorphism for both, compatibly with the fibre-inclusion map, by
\citet[Expos\'e XIII, Corollary 2.9 and \S2.10]{A_SGA1}.
The characteristic-zero comparison with topological fundamental
groups is \citet[Expos\'e XII, Corollary 5.2]{A_SGA1}. Apply these
comparisons to the commutative square formed by the fibre and total
space. The noninjectivity just proved gives the asserted
noninjectivity over $\overline{\mathbb F}_5$.
\end{proof}

The complete small certificate is
\texttt{code/A\_configuration\_obstruction\_certificate.py}, with
its exact output in the adjacent JSON file. It uses only the Python
standard library and verifies 34 free-word and finite-permutation
identities. The proof above supplies the group-topological and
specialization arguments; no inferred permutation-group order is
needed. The larger exploratory $A_5$ calculation is retained only
as an auxiliary calculation and is not a proof input.

This also explains the limitation of a pro-$\ell$ argument based on
the trivial action on punctured first homology. Such an argument
does not control finite quotients involving several primes. The
relative homotopy theorem of
\citet[Theorem C, Corollary 1.14 and Lemma 3.5]{A_HoshiSawadaTsujimura2024}
retains the full tame group of the base and treats a relative
prime-to-characteristic quotient of the fibre. Its pro-$\ell$
comparison does not assert the left exactness refuted here.

\paragraph{Why retaining the full base does not complete the point argument.}
There is a precise obstruction to descending the relative extension.
Write $\Pi=\pi_1(C_{\overline{\mathbb F}_5})$ and
$H=\ker(\Pi\to\Pi^{(5')})$. The outer point-pushing action on the
prime-to-five punctured fibre induces an action on the finite set
$\operatorname{Hom}(F,S_4)/S_4$. To transport this action, apply
SGA1 XIII, Corollary 2.9 to the punctured-fibre family over the
smooth lifted base curve: since $|S_4|=24$ is invertible, the sheaf
$R^1f_*S_4$ is finite locally constant and commutes with
specialization. Thus full-base specialization carries the finite
torsor action compatibly; no identification of the full
characteristic-zero and characteristic-five base groups is being
assumed. Let $\bar w\in\Pi$ be the specialization of the word
used above. In the procyclic closure of $\bar w$, take
its $5$-primary component: equivalently use the exponent
$\varepsilon\in\widehat{\mathbb Z}$ whose $5$-adic component is
$1$ and whose other prime-adic components are zero. Then
$\bar w^\varepsilon\in H$, but it still permutes the displayed
five classes in a five-cycle. Thus $H$ acts nontrivially through
\emph{outer} automorphisms of the prime-to-five fibre.

If the relative extension with this full prime-to-five fibre were
the pullback of an extension over $\Pi^{(5')}$, its outer action
would factor through $\Pi^{(5')}$, and $H$ would act trivially on
the same five classes. Hence such descent is impossible. In the
arithmetic setting, the relative extension retaining the full base
can be pulled back along a full point section. Equality of the
prime-to-$p$ images of two sections does not identify their full
lifts, so it does not identify these relative pullbacks by the
proposed argument.

This is a loss in the section comparison, not a claim that the
abstract curve or relative group is undetermined by its arithmetic
prime-to-$p$ group: the finite-field anabelian theorem already
reconstructs the curve. To finish this route one would need either
a means of identifying the relevant full-base lift of an individual
prime-to-$p$ section without assuming point separation, or an
anabelian reconstruction theorem for the strictly smaller fibre
quotient surviving in the full prime-to-$p$ configuration group.
Neither input is established here. The ordinary punctured-curve
theorem in \citet[Corollary 3.10]{A_SaidiTamagawa2009} does not
supply the latter input.

\subsection{All tame abelian covers of elliptic curves in the stated characteristics}

\begin{claim}{A.29}{Proved; abelian prime-to-$p$ covers of genus one without an exponent bound}
Let $K$ be finitely generated over $\F_p$, and let $C/K$ be
smooth proper geometrically connected of genus at least two.
Suppose that $C_{\overline K}$ has a faithful action of a finite
abelian group $G$ of order prime to $p$, with genus-one quotient
$E$. If $p$ is odd, or if $p=2$ and $E$ is supersingular, then
the geometrically prime-to-$p$ point-section map of $C$ is
injective. The same holds for every hyperbolic open subcurve.
There is no restriction on the exponent of $G$. In particular,
the odd-characteristic conclusion subsumes A.25 and A.27.
\end{claim}
\begin{proof}
We first work over a finite field, allowing finite constant
extensions to define all the finite geometric data used below.
It suffices initially to treat a cyclic cover $C\to E$ of degree
$n>1$, with $(n,p)=1$. It is ramified since $g(C)\ge2$.
Choose an origin on $E$ and a Kummer equation
\[
 z^n=c\in\overline k(E)^*.
\]
Set $N=2n$ if $p$ is odd and $N=n$ if $p=2$. In the latter
case $n$ is odd. Pull back by the \'etale isogeny
$[N]:E'\to E$, where $E'$ denotes another copy of $E$, and
choose a connected component $C'$ of the normalization. Then
$C'\to C$ is a geometrically prime-to-$p$ Galois \'etale cover:
the full pullback is an $E[N]$-torsor, and a component is a
torsor under its stabilizer. It has genus at least two.
The map $C'\to E'$ is cyclic of degree dividing $n$, and the
chosen root $z$ still satisfies
\[
 z^n=c':=[N]^*c.
\]
No connectedness of the entire pullback is assumed.

Let $a$ be the nonzero residue modulo $n$ of the valuation of
$c$ at one of its branch points. The inverse image of this
point is a coset of $E'[N]\simeq(\Z/N)^2$, and the valuation
of $c'$ has the same residue $a$ at every point of that coset.
At each such point the local ramification index of every
component is
\[
 e=\frac n{\gcd(n,a)}.                                      \tag{A.16}
\]
This formula is valid even when the chosen component has
degree smaller than $n$.

\emph{An explicit second-difference identity.}
For integers $i$ put $q_i=i(i-1)/2$. This sequence is periodic
modulo $n$ with period $N$, because
\[
 q_{i+N}-q_i=Ni+\frac{N(N-1)}2\equiv0\pmod n.
\]
Here $N=2n$ always works, while $N=n$ works for odd $n$.
Moreover $q_{i-1}+q_{i+1}-2q_i=1$. Therefore, on a cyclic
row of $N$ points with generator $d$,
\[
 \sum_{i\bmod N}\delta_{x+id}
 =\sum_{i\bmod N}q_i
  \bigl(\delta_{x+(i+1)d}+\delta_{x+(i-1)d}
                  -2\delta_{x+id}\bigr)\pmod n.             \tag{A.17}
\]
Choose a basis $d,d_2$ of $E'[N]$ and sum (A.17) over the
$N$ rows of each branch coset. Multiplying by its coefficient
$a$ expresses the full branch divisor of $c'$ modulo $n$ as
a sum of symmetric triples, each coefficient being an integer
multiple of the original coefficient $a$ of its coset.

For a centre $x\in E'$ let $u_x:E'\to\mathbf P^1$ be the
elliptic inversion quotient with origin $x$. The generator
$d$ has order $N>2$, so the two points $x+d,x-d$ are distinct
and are unramified for this double map. After scaling its
coordinate, the rational function
\[
 h_x=u_x-u_x(x+d)
\]
has exactly the divisor
\[
 \operatorname{div}(h_x)=(x+d)+(x-d)-2(x).                  \tag{A.18}
\]
In characteristic two the inversion map is still separable
of degree two, and nonzero odd-torsion points are unramified,
so the same divisor formula holds.

List the triples with their coefficients as $(h_j,c_j)$,
where $c_j$ is a multiple of the coefficient $a$ on its
branch coset. Formula (A.17) gives
\[
 \operatorname{div}\left(c'\Big/\prod_jh_j^{c_j}\right)
                         \equiv0\pmod n.                 \tag{A.19}
\]
Terms with $c_j\equiv0\pmod n$ may be omitted. For each
remaining term, choose a root
$w_j^n=h_j^{c_j}$ and let $B_j/E'$ be its connected smooth
cover. Put $r_j=\gcd(n,c_j)$ and $m_j=n/r_j>1$.
Since $h_j$ has simple zeros, this cover has degree $m_j$.
After rescaling constants and using B\'ezout's identity it
has the equivalent equation
\[
 s_j^{m_j}=h_j.                                           \tag{A.20}
\]
Indeed, writing $c_j=r_jb_j$ with $\gcd(b_j,m_j)=1$, a
component has $w_j^{m_j}=\zeta h_j^{b_j}$; remove $\zeta$
by a constant rescaling and choose $u b_j+v m_j=1$ to set
$s_j=w_j^u h_j^v$.

\emph{The auxiliary curves separate their point sections.}
Suppose first that $p$ is odd. In coordinates with origin
the centre of the triple, write $E'$ as
$v^2=R(u)$, where $R$ is a separable cubic. If
$h_j=u-b$, then $b$ is not a branch value, because the zeros
in (A.18) are unramified. Thus (A.20) gives
\[
 B_j:\quad v^2=R(s_j^{m_j}+b).
\]
The polynomial on the right is squarefree of degree $3m_j$:
its three values before extracting $m_j$-th roots are
nonzero, and $(m_j,p)=1$. Hence $B_j$ is hyperelliptic with
\[
 g(B_j)=\frac{3m_j-\gcd(m_j,2)}2\ge2.
\]
A.19 supplies point-section injectivity for every $B_j$.
In particular, no genus-one auxiliary exception is needed.

If $p=2$ and $E'$ is supersingular, it has a Weierstrass
equation
\[
 v^2+v=u^3+A u+B
\]
over the algebraic closure. For completeness, in a general
characteristic-two Weierstrass equation, a nonzero coefficient
$a_1$ gives a nonzero two-torsion point at $a_1u+a_3=0$,
using the formula for inversion. Supersingularity therefore
forces $a_1=0$; smoothness forces $a_3\ne0$. Scaling and
translating the $u$-coordinate gives the displayed form.
After (A.20) the auxiliary curve is
\[
 v^2+v=(s_j^{m_j}+b)^3+A(s_j^{m_j}+b)+B.
\]
It is an Artin--Schreier curve with one pole, of odd order
$3m_j$. The differential calculation in A.21 applies
verbatim to any polynomial of odd degree: if its genus is
$g$, then $s_j^i\,ds_j$, $0\le i\le g-1$, form a basis
of regular differentials, and Cartier kills even indices
and sends an odd index $i$ to $(i-1)/2$. It is nilpotent.
As in A.21, a nonzero geometric two-torsion class would
produce a nonzero regular Cartier-fixed logarithmic
differential. Thus $J_{B_j}(\overline k)[2^\infty]=0$.
A.1 and the Abel map give point-section injectivity for
$B_j$ over every finite constant extension.

\emph{An \'etale common cover and an unramified residual character.}
Let $V$ be the connected normalization in the compositum of
the function fields of $C'$ and all the $B_j$. It is an
abelian Galois cover of $E'$ of exponent dividing $n$.
At any point in a branch coset with coefficient $a$, each
auxiliary local exponent is either $c_j$, $-2c_j$, or zero,
with $a\mid c_j$ as integers. Its ramification index
therefore divides $n/\gcd(n,a)$. No auxiliary branches
outside the branch cosets. Tame cyclic extensions of the
completed local field over an algebraically closed residue
field have compositum ramification index the least common
multiple of their indices. Consequently $V\to C'$ is
\'etale, and the inertia order of $V/E'$ at any point is
exactly that of $C'/E'$ there.

The cover $V\to C'$ is Galois of prime-to-$p$ order.
The composite $V\to C$ has prime-to-$p$ geometric Galois
closure: take all conjugates over the Galois cover
$C'\to C$. Their compositum has group a subgroup of a
finite product of the prime-to-$p$ groups of $V/C'$, extended
by the prime-to-$p$ group of $C'/C$. Equivalently it embeds
in the corresponding finite wreath product. This checks
the exact lifting hypothesis of A.3.

Inside $\overline k(V)$ put
\[
 w_H=z\Big/\prod_jw_j.
\]
By (A.19), the connected cover $H/E'$ defined by $w_H$ is
unramified; its degree divides $n$, so $H$ has genus one.
This also includes the degree-one case. Since $C'/E'$ has
a nonempty branch set, choose a branch point $t$. The cover
$H/E'$ is unramified there, while $V/E'$ has the nontrivial
prime-to-$p$ index $e$ of $C'/E'$. Thus all indices of
$V\to H$ above any point of $H$ over $t$ equal $e$.
A.21 shows that every collision on $V$ has equal images
on $H$.

\emph{Inertia retains exactly the point information required.}
Suppose $P,Q$ is a collision on $C$. By A.3 and the checked
Galois-closure condition, after a finite constant extension
it has compatible lifts $P_V,Q_V$ forming a collision on
$V$. Their images on every $B_j$ agree by the injectivity
proved above, and their images on $H$ agree by A.21.
In particular they have the same image $t_0$ on $E'$.
Write $Q_V=gP_V$ with
$g\in\Gamma:=\operatorname{Gal}(V/E')$, and let
$I\subset\Gamma$ be the inertia group at $P_V$.

Use the chosen Kummer roots to define characters
$\chi_0,\chi_j,\chi_H:\Gamma\to\mu_n$ for $z,w_j,w_H$.
They satisfy
\[
 \chi_0=\chi_H\prod_j\chi_j.                              \tag{A.21}
\]
Let $e_0$ be the ramification index of $C'/E'$ at $t_0$,
with $e_0=1$ allowed. Its inertia image is
$\chi_0(I)=\mu_{e_0}$. Equal $B_j$-images mean
$\chi_j(g)\in\chi_j(I)$; by the ramification divisibility,
\[
 \chi_j(I)\subseteq\mu_{e_0}.
\]
The residual cover is unramified, so $\chi_H(I)=1$.
Equal $H$-images therefore give $\chi_H(g)=1$. Equation
(A.21) now yields
$\chi_0(g)\in\mu_{e_0}=\chi_0(I)$. Equivalently
$g\in(\ker\chi_0)I$, so the two points have the same
image on $C'$, and hence $P=Q$ on $C$.
This argument includes every partially ramified fibre.
It does not require the projections to the auxiliary
curves to distinguish all points of the normalization $V$.

For a general finite abelian group $G$, take all its
nontrivial cyclic character quotients $C_\chi$. They map
separably to $E$, so none has genus zero. If
$g(C_\chi)\ge2$, the cyclic theorem gives equal
$C_\chi$-images of a collision on $C$. If its genus is
one, the Galois map $C\to C_\chi$ is ramified by
Riemann--Hurwitz, and A.21 gives the same equality.
These equalities imply the same $E$-image. Write $Q=gP$
and $I=\operatorname{Stab}_G(P)$. For every character,
$g\in(\ker\chi)I$. Characters of $G/I$ separate its
elements, so $g\in I$ and $P=Q$.

Finally define the finite tame action and genus-one quotient
over a finite extension of $K$, spread them over a smooth
finite-type model, and apply A.8. In characteristic two
the supersingular condition persists: the elliptic quotient
has identically zero $j$-invariant. The resulting finite-field
fibres therefore remain in the proved class. Specialization
excludes a collision over $K$, and functoriality to the
proper compactification proves the assertion for open
subcurves.
\end{proof}

The exact second-difference, ramification-divisibility, genus,
and inertia-subgroup calculations are independently reproducible
with \path{code/A_elliptic_divisor_certificate.py}; its
adjacent JSON output includes even and composite degrees and
all partial branch coefficients in the tested degrees.
The displayed identities above prove these facts for all $n$;
the finite computation is a check, not a replacement for them.

\begin{claim}{A.29a}{Proved; an obstruction over an ordinary elliptic quotient in characteristic two}
Let $c\in\F_4$ satisfy $c^2+c+1=0$, let
\[
 E:\quad y^2+y=x+c/x,
\]
and let $B$ be the smooth proper curve with function field
\[
 B:\quad y^2+y=s^3+c+\frac{c}{s^3+c}.                       \tag{A.22}
\]
The elliptic curve $E$ is ordinary, and $B\to E$, given by
$x=s^3+c$, is exactly a degree-three auxiliary of the
construction in A.29, with its symmetric triple centred at
the origin and its two zeros at nonzero three-torsion points.
Nevertheless, over $\overline{\F}_2$ there is no nonconstant
morphism from $B$ to any abelian variety of $2$-rank zero.
Thus replacing each ordinary-elliptic auxiliary by a direct
map to such an abelian variety cannot finish the missing
characteristic-two case. This is not a counterexample to A1.
\end{claim}
\begin{proof}
The Artin--Schreier map $E\to\mathbf P^1_x$ has two simple
poles, at zero and infinity, so $g(E)=1$. Its two ramification
points differ by a nonzero two-torsion class, as
$\operatorname{div}(x)=2(P_0-P_\infty)$. Thus $E$ is ordinary.
Set $Y=xy+c$. Then
\[
 Y^2+xY=x^3+c^2.
\]
Choose $d$ with $x(d)=c$, so $y(d)^2+y(d)=c^2$.
The tangent line at $d$ is
\[
 Y=(y(d)+c^2)x+c^2.
\]
Since $(y(d)+c^2)^2+(y(d)+c^2)=c$, substitution into the
Weierstrass equation gives $(x+c)^3=0$. Hence $d$ is a
nonzero flex, and has exact order three. Equivalently, the
three-division polynomial is $x^4+x^3+c^2$, which vanishes
at $x=c$. Thus $s^3=x+c$ is the auxiliary associated to
$(d)+(-d)-2(P_\infty)$.

In (A.22) there are three simple finite poles, at the distinct
roots $a_i$ of $s^3+c$, and one pole of order three at
infinity. Their different exponents are respectively
$2,2,2,4$. Riemann--Hurwitz gives $g(B)=4$. More explicitly,
$ds$ has order two at each finite ramification point and
order zero at infinity. The four regular differentials
\[
 ds,\quad \frac{ds}{s-a_1},\quad
 \frac{ds}{s-a_2},\quad\frac{ds}{s-a_3}
\]
are independent, hence form a basis. Cartier kills $ds$ and
fixes the three logarithmic forms. Its stable rank is three,
so $B$ has $2$-rank three. One can also read the lower bound
from the three independent classes $[P_{a_i}-P_\infty]$ of
order two: a principal combination with coefficients zero
or one would have an invariant defining function under the
Artin--Schreier involution, hence a function of $s$, whose
valuation at every ramification point is even. The logarithmic
differential injection used in A.21 gives the matching upper
bound from the displayed Cartier-fixed space.

The following exact counts include the unique normalized point
over each rational pole and the point at infinity:
\[
\begin{array}{c|rrrr}
 r&1&2&3&4\\ \hline
 |B(\F_{4^r})|&7&33&82&225.
\end{array}
\]
The certificate
\path{code/A_ordinary_elliptic_obstruction.py} verifies
them in two ways: exhaustive enumeration of $y$ for each
nonpole $s$, and the Artin--Schreier trace criterion. It uses
explicit irreducible binary polynomials of degrees $2,4,6,8$,
whose irreducibility is itself checked by exhaustive trial
division. Thus the counts are reproducible finite arithmetic.
Newton's identities and the genus-four functional equation
give the Frobenius characteristic polynomial
\citep[II, Theorem 1.1]{A_Milne2008}
\[
\begin{split}
 P(X)&=(X^2-X+4)\,
 \bigl(X^6+3X^5+9X^4+20X^3+36X^2+48X+64\bigr)\\
 &=X^8+2X^7+10X^6+23X^5+52X^4+92X^3
                          +160X^2+128X+256.
\end{split}
\tag{A.23}
\]
The four power sums used here are $-2,-16,-17,32$.

There is no root of (A.23) of the form $2\zeta$ with
$\zeta$ a root of unity. If such a root has order $m$, then
$\varphi(m)\le8$. Every prime divisor $\ell$ of $m$ must
satisfy $\ell-1\le8$, so it is $2,3,5,$ or $7$; their
exponents are at most $4,2,1,1$, respectively. Filtering
these possibilities by the totient condition gives exactly
\[
 1,2,3,4,5,6,7,8,9,10,12,14,15,16,18,20,24,30.
\]
For every member of this finite list the certificate performs
the rational Euclidean algorithm and verifies
\[
 \gcd\bigl(P(2X),\Phi_m(X)\bigr)=1.
\]
It constructs each cyclotomic polynomial recursively from
$X^m-1$, so no numerical approximation or root recognition
is involved.

If $J_B$ had a nonzero quotient of $2$-rank zero and dimension
$r$, Poincar\'e reducibility and additivity of the $2$-rank
under isogeny would give $3\le4-r$. Thus $r=1$, and this
quotient would be a supersingular elliptic curve. Over a
finite extension of $\F_4$ it would be isomorphic to the
supersingular curve in A.10, whose $\F_4$-Frobenius is
$[-2]$. The corresponding quotient of the Tate module would
force some root $\alpha$ of $P$ to satisfy
$\alpha^a=(-2)^a$ for an integer $a>0$. This says that
$\alpha/2$ is a root of unity, just excluded. Finally a
nonconstant map from $B$ to an abelian variety of $2$-rank
zero induces a positive-dimensional such quotient of $J_B$,
after translating one point to the origin. No such map exists.
\end{proof}

\begin{claim}{A.29b}{Proved; the obstruction survives the translated cubic refinement}
Use $E,B,d$ from A.29a and put $h=x+c$, with
$\operatorname{div}(h)=(d)+(2d)-2(O)$. Let
$\tau(t)=t+d$, and write $h_i=\tau^{i*}h$ for $i=0,1,2$.
Let $V$ be the connected normalization obtained by adjoining
cube roots of $h_0$ and $h_1$ to $\overline k(E)$.
Then $V\to B$ is \'etale of degree three and $g(V)=10$.
It has a Heisenberg group of order $27$ over $E/\langle d\rangle$,
with central inertia of order three. Nevertheless, $V$ also
admits no nonconstant map to any abelian variety of $2$-rank
zero, even over $\overline{\F}_2$. Thus this particular first
cubic refinement does not remove the obstruction in A.29a.
\end{claim}
\begin{proof}
In the order $(O,d,2d)$ the three divisor vectors are
\[
 (-2,1,1),\qquad(1,1,-2),\qquad(1,-2,1).                  \tag{A.24}
\]
Their sum is zero, so $h_0h_1h_2=\kappa$ is a nonzero
constant. Moreover
\[
 \operatorname{div}(h_1/h_0)=3(O)-3(2d).
\]
The corresponding unramified cubic character is nontrivial:
if $h_1/h_0$ were a cube, its cube root would have divisor
$(O)-(2d)$, impossible since this is a nonzero order-three
class on the elliptic curve. The classes of $h_0,h_1$ in
$\overline k(E)^*/\overline k(E)^{*3}$ are independent.
Indeed their branch residues force the sum of the coefficients
of any relation to be zero modulo three, leaving a power of
this nontrivial unramified character. Hence $V/E$ has group
\[
 A=\{(\epsilon_0,\epsilon_1,\epsilon_2)\in\mu_3^3:
                            \epsilon_0\epsilon_1\epsilon_2=1\}
                    \simeq(\Z/3)^2.
\]
Choose $u_i^3=h_i$ with
$u_0u_1u_2=\kappa^{1/3}$. At each of the three branch
points, all residues in (A.24) equal one modulo three.
Consequently the inertia subgroup is
\[
 Z=\{(\zeta,\zeta,\zeta):\zeta\in\mu_3\}\subset A,
\]
of order three, and there is no other ramification.
The indices on $V/E$ and $B/E$ therefore agree at every
point, proving that $V/B$ is \'etale. Its degree is three,
and $g(V)-1=3(g(B)-1)=9$.

Translation $\tau$ lifts by cycling $u_0,u_1,u_2$. This
lift has order three, since it restores each generator after
three iterations and preserves their prescribed product.
It cyclically permutes the three coordinates of $A$.
Thus $V\to E/\langle d\rangle$ has group
$A\rtimes\langle\tau\rangle$, of order $27$ and exponent
three. Its centre and commutator subgroup are both $Z$;
these assertions follow directly from the permutation action
on triples with coordinate sum zero modulo three. This is
the Heisenberg group of order $27$. Its inertia remains the
central subgroup $Z$.

Let $B_i$ be the quotient retaining $u_i$ and the coordinate
on $E$. Each $B_i$ is isomorphic to $B$ by translation.
Its degree-three kernel $H_i\subset A$ is the subgroup with
$i$-th coordinate one. Choose a generator $z$ of $Z$ and a
generator $b$ of one of the $H_i$. The three kernels are,
in some order,
\[
 H_i=\langle bz^i\rangle,\qquad i=0,1,2.
\]
The quotient $E_2=V/Z$ is an unramified degree-three cover
of $E$, hence has genus one. The action of $b$ on $E_2$ is
a deck translation of this elliptic isogeny, and so acts
trivially on $J_{E_2}$.

Write $N_H=\sum_{a\in H}a$ for the corresponding
endomorphism of $J_V$. The equality of integral group-ring
elements
\[
 \sum_{i=0}^2N_{H_i}=3\operatorname{Id}+(b+b^2)N_Z
\]
follows by expanding the three order-three subgroup sums.
The image of $N_Z$ lies in the pullback of $J_{E_2}$,
because $N_Z$ is pullback after the norm for $V\to E_2$.
The translation observation therefore gives
\[
 \sum_{i=0}^2N_{H_i}=3\operatorname{Id}+2N_Z,
 \qquad N_Z^2=3N_Z
\]
as endomorphisms of $J_V$. In particular,
\[
 \left(\sum_{i=0}^2N_{H_i}\right)
            (9\operatorname{Id}-2N_Z)=27\operatorname{Id}. \tag{A.25}
\]

For $q_i:V\to B_i$ we have
$q_i^*q_{i,*}=N_{H_i}$. Thus the image of the homomorphism
\[
 \Phi:\prod_{i=0}^2J_{B_i}\longrightarrow J_V,
                 \qquad (a_i)\longmapsto\sum_iq_i^*a_i
\]
contains the image of the left-hand side of (A.25).
Multiplication by $27$ is surjective on $J_V$, so $\Phi$
is surjective. This is the needed norm-image assertion;
it is not inferred by counting dimensions.

If a nonconstant map from $V$ to an abelian variety $T$ of
$2$-rank zero existed, translating one value to the origin
would give a nonzero homomorphism $J_V\to T$. Its composite
with each $q_i^*$ is zero by A.29a, since $B_i\simeq B$.
Surjectivity of $\Phi$ would make the homomorphism zero,
a contradiction. The finite group and integral group-ring
calculations are verified in
\path{code/A_cubic_refinement_certificate.py} and its
adjacent JSON. The code separates the group-ring identity
from the geometric translation relation $(b-1)N_Z=0$,
which is justified above. No conclusion is drawn about
other odd \'etale covers of $B$ or $V$, or about A1 itself.
\end{proof}

\begin{claim}{A.29c}{Proved; separation despite the direct-target obstruction}
For the ordinary-elliptic auxiliary
\[
 B:\quad y^2+y=s^3+c+\frac{c}{s^3+c},\qquad c^2+c+1=0,
\]
of A.29a, A1 holds over every finite extension of $\F_4$.
Thus the absence of a nonconstant map from either $B$ or its
first translated refinement $V$ to an abelian variety of
$2$-rank zero does not prevent point-section injectivity.
\end{claim}
\begin{proof}
Put $\delta(s,y)=(cs,y)$ and $\iota(s,y)=(s,y+1)$.
These automorphisms commute. Write $O_B$ for infinity and
$W_0,W_1,W_2$ for the finite Weierstrass points, lying over
$s^3+c=0$, and number them cyclically under $\delta$.
The $\delta$-fixed points are $O_B$ and the two points with
$s=0$. Under $B\to E$, all three $W_i$ map to the nonzero
two-torsion point $T$ of $E$ above $x=0$.

Let
\[
 A=\bigl(\ker(\operatorname{Nm}_{B/E}:J_B\to E)\bigr)^0,
 \qquad \gamma(P)=[P-\delta P].
\]
The image of $\gamma$ lies in $A$: it is connected, has
zero norm, and contains zero. The degree-three norm and
pullback give an isogeny decomposition, so
$\dim A=3$ and the $2$-rank of $A$ is two. By A.29a its
$\F_4$-Frobenius polynomial is
\[
 R(X)=X^6+3X^5+9X^4+20X^3+36X^2+48X+64.                \tag{A.26}
\]

\emph{Two divisor facts.}
First, every nonzero value of $\gamma$ has a unique
geometric preimage. Indeed, $\gamma(P)=\gamma(Q)$ implies
\[
 P+\delta Q\sim Q+\delta P.
\]
If these effective divisors are equal, either $P=Q$ or
both points are $\delta$-fixed. Otherwise they are two
fibres of the hyperelliptic pencil. Consequently
$\delta Q=\iota P$ and $\delta P=\iota Q$, which imply
$\delta^2P=P$, hence $\delta P=P$. Then $Q=\iota P$
is fixed as well. The zero fibre consists exactly of the
$\delta$-fixed points, by injectivity of the Abel map.

Second, for non-Weierstrass points $P,Q$,
\[
 \gamma(P)-\gamma(Q)\notin A[2](\overline{\F}_2)
                                      \setminus\{0\}.    \tag{A.27}
\]
The nonzero points of this two-dimensional $\F_2$-space
are precisely $[W_i-W_j]$ for $i\ne j$. They lie in $A$
because adjacent differences are $\gamma(W_i)$; the
independent classes $[W_i-O_B]$ from A.29a show that
their span has rank two. If (A.27) failed, we would have
\[
 P+\delta Q+W_j\sim Q+\delta P+W_i,\qquad i\ne j.        \tag{A.28}
\]
Each side has exactly one Weierstrass summand. If neither
side contains a hyperelliptic pair, uniqueness of reduced
divisors of degree three makes the divisors equal, which
their different Weierstrass summands exclude. If exactly
one contains such a pair, reduction gives degrees one and
three, again impossible. If both contain pairs, they are
$(P,\delta Q)$ and $(Q,\delta P)$; removing them would
give $[W_j-O_B]=[W_i-O_B]$, also impossible.

Only the elementary degree-two and degree-three cases of
hyperelliptic reduction are needed here. They can also be
proved without an algorithm: a function of degree $d\le3$
outside the hyperelliptic rational subfield gives a
birational curve of bidegree $(2,d)$ in
$\mathbf P^1\times\mathbf P^1$. Its arithmetic genus is
$d-1\le2$, contradicting $g(B)=4$. Thus every moving
system of those degrees is the hyperelliptic pencil plus
its fixed divisor. The quadratic hyperelliptic extension
is separable in characteristic two, so this argument
applies without changing characteristic.

\emph{The Frobenius bound.}
The norm relation gives $1+\delta+\delta^2=0$ on $A$.
The geometric \'etale Tate module $T_2A$ has rank two
over $\Z_2$ and is a torsion-free rank-one module over
\[
 \mathcal O=\Z_2[\zeta_3],
\]
the unramified quadratic discrete valuation ring. It is
therefore free of rank one. Since $\delta$ is defined over
$\F_4$, Frobenius $\Phi$ commutes with this action and
acts on that module by a unit $\lambda\in\mathcal O^*$.
This module argument is essential; point counts alone
would not establish the asserted module structure.

Exact determinants of powers of the integral companion
matrix of (A.26) give
\[
 \begin{aligned}
 |A(\F_{4^3})|&=291772=2^2\cdot72943,\\
 |A(\F_{4^6})|&=70459436736=2^6\cdot1100928699.
 \end{aligned}                                         \tag{A.29}
\]
The geometric $2$-primary point group is
$\operatorname{Frac}(\mathcal O)/\mathcal O$. The kernel
of $\Phi^r-1$ on it has order
$2^{2v_{\mathcal O}(\lambda^r-1)}$. Consequently
\[
 v_{\mathcal O}(\lambda^3-1)=1,
 \qquad v_{\mathcal O}(\lambda^6-1)=3.                  \tag{A.30}
\]
Suppose $\gamma(P)$ has exact order $2^a$, $a\ge4$.
For $\sigma=\Phi^{6\cdot2^{a-4}}$ the identity
$(1+v)^2-1=2v+v^2$ shows
\[
 v_{\mathcal O}\bigl(\lambda^{6\cdot2^{a-4}}-1\bigr)
                                                          =a-1.
\]
Here one starts with $\lambda^6=1+8u$, $u$ a unit;
every subsequent squaring increases the valuation by
exactly one. Hence $\sigma\gamma(P)-\gamma(P)$ is a
nonzero point of $A[2]$. Since Weierstrass points have
$\gamma$-value of order at most two, both $P$ and
$\sigma P$ are non-Weierstrass. This contradicts (A.27).
Every $2$-primary value of $\gamma$ therefore has order
at most eight. By (A.30), $\Phi^6$ fixes $A[8]$; the
unique-preimage assertion proves
\[
 0\ne\gamma(P)\in A[2^\infty]
                  \quad\Longrightarrow\quad P\in B(\F_{4096}).
                                                               \tag{A.31}
\]
No bound uniform in curves or covers is being assumed.

\emph{The complete finite certificate.}
Set $H=s^3+c$ and $Y=Hy$. Then
\[
 Y^2+H(s)Y=F(s),\qquad
 F=H^3+cH=s^9+c s^6+s^3+c.                              \tag{A.32}
\]
Both independent programs use
$\F_2[t]/(t^{12}+t^3+1)$ and verify that $t+1$ has a
cycle of all $4095$ nonzero residues. This verifies that
the quotient is a field; in binary coefficient encoding
$c=72$. Exhaustive enumeration gives $4206$ points,
including infinity. Because $\gamma(\delta P)=\delta\gamma(P)$
and $\gamma(\iota P)=-\gamma(P)$, one representative
from each orbit under $\delta$ and $\iota$ suffices.
There are $702$ affine orbits, of which two are exceptional:
the two fixed points with $s=0$, and the three finite
Weierstrass points. Their $\gamma$-orders are respectively
one and two. Infinity is fixed as well.

Here is a direct linear-algebra certificate excluding
$8\gamma(P)=0$ on every other orbit. For $Q=\delta P$,
the principal divisor of $(s-s(Q))^8$ is
$8Q+8\iota Q-16O_B$. Thus $8(P-Q)$ is principal
if and only if there is a nonzero element of
$L(16O_B)$ vanishing to order at least eight at both
$P$ and $\iota Q$. The pole orders of $s$ and $Y$ are
two and nine. Riemann--Roch gives $\dim L(16O_B)=13$,
and its following $13$ linearly independent elements
are a basis:
\[
 1,s,s^2,\ldots,s^8,\quad Y,sY,s^2Y,s^3Y.              \tag{A.33}
\]
Independence also follows by separating the rational
subfield from its $Y$-multiple. At either non-Weierstrass
point, $s-s(P)$ is a uniformizer and $H(P)\ne0$.
Expanding (A.32) uniquely determines the first eight
coefficients of $Y$ recursively, since the coefficient
of each new unknown is $H(P)$. The sixteen vanishing
conditions on the basis (A.33) form a $16$ by $13$
matrix over $\F_{4096}$.

The self-contained program
\path{code/root_A29c_linear_certificate.py} verifies
the series equation coefficient by coefficient and
performs exact Gaussian elimination. Every one of the
$700$ nonexceptional orbit matrices has rank $13$.
Its adjacent JSON records the ranks, point count, and
a hash of all matrix entries. A nonzero numerator with
the specified zeros would have a pole of order exactly
$16$ at infinity: its required zero divisor already
has degree $16$. Thus the matrix criterion is an
equivalence, including the reverse implication, and
proves the desired exclusion.

An independent certificate,
\path{code/A_genus4_prym_experiment.py}, computes the
same packet by exact Mumford divisor arithmetic. Its
composition and reduction formulas, with their ideal
and principal-divisor proofs, are given in
\path{audit/A_genus4_prym_certificate.md}. It verifies
every division and curve congruence, compares $2103$
Cantor doublings with independent Hensel doublings,
and checks $702$ inverse additions and $702$ deck-cycle
sums. It also computes (A.29) by exact integer
determinants. The two methods agree. Together with
(A.31), their conclusion is
\[
 0\ne\gamma(P)\in A[2^\infty]
                  \quad\Longrightarrow\quad
                  P\in\{W_0,W_1,W_2\}.                \tag{A.34}
\]

\emph{The translated cover eliminates the packet.}
Suppose distinct $P,Q$ have equal prime-to-two section
classes over a finite extension. The ramified cyclic
map $B\to E$ and A.21 force their $E$-images to agree.
Thus $Q=\delta^jP$, $j=1$ or $2$, and A.1 makes
$[P-Q]$ $2$-primary. If $j=2$, use
$[P-\delta^2P]=-\delta^2\gamma(P)$; in either case
(A.34) puts $P,Q$ in the finite Weierstrass triple.
Their common image on $E$ is $T$.

Lift this collision through the \'etale cubic cover
$V\to B$ of A.29b, using A.3 and a finite constant
extension. Its central inertia subgroup is $Z$ and
$E_2=V/Z$ is elliptic. The cyclic cubic map
$V\to E_2$ is ramified, with nine branch points, so
A.21 forces the two lifts to have equal $E_2$-images.
They differ by a nonidentity element $z^j\in Z$,
since their images on $B$ are distinct.

For $i=0,1,2$, the quotient $B_i$ retaining the root
of $h_i=\tau^{i*}h$ is identified with $B$ by the
translation $t\mapsto t+id$ on $E$. This identification
sends $z$ to the same deck transformation $\delta$.
The two lifts therefore give a nontrivial $\delta$-pair
on each $B_i$ with $2$-primary divisor difference.
The hyperelliptic involution on $B_i$ covers
$t\mapsto-t-2id$; its Weierstrass points lie precisely
over $-id$ and $T-id$. For $i=1,2$, neither base point
equals $T$, because $id$ is nonzero of order three
and $T$ has order two. Hence the images over the actual
base point $T$ are non-Weierstrass. This contradicts
(A.34) under the indicated identification $B_i\simeq B$.
The collision is impossible.
\end{proof}

\subsection{What remains of A1 and A2}

\begin{claim}{A.30}{Unresolved; exact remaining assertion}
For every finite field $k$ of characteristic $p$, every smooth proper
geometrically connected $C/k$ of genus at least two, and distinct
$P,Q\in C(k)$, the two classes of sections in $E_C^{(p')}$ are
distinct. Equivalently, the conditions of A.3 force $P=Q$ in this
range. The remaining cases include arbitrary curves outside the
abelian prime-to-$p$ genus-zero quotient class of A.24,
the genus-one quotient class of A.29, the
$S_3$ class of A.26, and
the other cases proved above; in characteristic two, this
includes hyperelliptic curves outside
the cases proved above. No general proof or counterexample is
supplied here. Thus A1 and the final implication in A2 are not
finished in their full scope.
\end{claim}

The inspected result of \citet[Proposition 1.8(vi) and Corollary
1.11(iii)]{A_SaidiTamagawa2009} is finiteness of an exceptional set,
not its emptiness. Its proof uses a finiteness theorem for
prime-supported torsion on the difference surface in the Jacobian;
it supplies no sublinear bound as the curve varies in an \'etale
tower. The refinement \citep{A_SaidiTamagawa2017} still retains
exceptional sets. Our searches are recorded separately; we make
no priority claim for A.2--A.4.

The tame-Belyi theorem of \citet[Theorems 1.2 and
1.5]{A_KedlayaLittWitaszek2023} does not by itself reduce
the remaining assertion to A.24. A tame map need not have
prime-to-$p$ Galois-closure order, and the map from its
Galois closure to the original curve need not be \'etale.
Even for a nonabelian prime-to-$p$ Galois cover, cyclic
character quotients only determine a point modulo the
commutator subgroup and its stabilizer. The argument of
A.24 does not remove that remaining fibre ambiguity.

On scan 36 the phrase ``de $X$'' in the covering sentence has to be
interpreted mathematically as a cover of $Y$, the curve containing
$P,Q$: otherwise its fibres above these points are not defined.
The visible wording was inspected in the supplied image; this is
an editorial correction supported by the types of the maps, not
an assertion that the ink reads $Y$. Likewise ``$p$-torsion'' has
to mean $p$-primary torsion, with no uniform exponent asserted.
